% =============================================================================
% tesis.tex · Documento maestro de la tesis
% «Detección evolutiva de regímenes del ciclo financiero peruano mediante
%  Clasificación Factorial Jerárquica y aprendizaje automático»
% Compilar desde la raíz del repositorio:  latexmk -lualatex -cd tesis_latex/tesis.tex
% =============================================================================
\nonstopmode % no detenerse a pedir entrada (latexmk -cd no lee tesis_latex/latexmkrc)
\documentclass[12pt, a4paper, oneside, openany]{book}

% =============================================================================
% preambulo.tex · Paquetes, formato y entornos de la tesis maestra
% Formato provisional (decisión B, 20-10-2026): clase book con los márgenes y el
% espaciado de tesis/plantilla_uni.docx. Debe revisarse cuando se disponga de las
% «Pautas para redactar el Trabajo de Tesis de Maestría» de la UNI.
% =============================================================================

% ---- Página: A4, 3 cm a la izquierda y 2,5 cm en los demás lados ----------------
% La plantilla Word no fija el tamaño de papel; se adopta A4, el estándar en el Perú.
\usepackage[a4paper, left=3cm, right=2.5cm, top=2.5cm, bottom=2.5cm,
            footskip=1.25cm, headheight=15pt]{geometry}

% ---- Codificación y tipografía ----------------------------------------------
% El preámbulo no usa fontspec: en este entorno TeX Live no incluye luaotfload (paquete
% texlive-luatex), sin el cual fontspec no funciona. Se usan fuentes Type 1 con
% codificación T1, que compilan igual en cualquier instalación de TeX Live.
% Texto: Nimbus Sans (paquete helvet), clon de Helvetica; Arial reproduce los anchos de
% Helvetica, de modo que es el equivalente métrico de Arial 12. Matemáticas: Computer
% Modern (amsfonts), porque no hay una fuente matemática sin serifas instalada.
\usepackage[T1]{fontenc}
\usepackage{textcomp}
\usepackage{helvet}
\renewcommand{\familydefault}{\sfdefault}
\renewcommand{\rmdefault}{phv}       % evita sustituciones: no hay Latin Modern T1 instalada
\renewcommand{\ttdefault}{pcr}       % monoespaciada: Courier (Type 1; las EC solo existen en mapa de bits)
% Sin luaotfload, LuaLaTeX pasa el código Unicode a la fuente T1, que coincide con Latin-1
% en las letras acentuadas pero no en los signos siguientes; se definen explícitamente.
\usepackage{newunicodechar}
\newunicodechar{«}{\guillemotleft}
\newunicodechar{»}{\guillemotright}
\newunicodechar{¿}{\textquestiondown}
\newunicodechar{¡}{\textexclamdown}
\newunicodechar{—}{\textemdash}
\newunicodechar{–}{\textendash}
\newunicodechar{“}{\textquotedblleft}
\newunicodechar{”}{\textquotedblright}
\newunicodechar{‘}{\textquoteleft}
\newunicodechar{’}{\textquoteright}
\newunicodechar{·}{\textperiodcentered}
\newunicodechar{ń}{\'n}                % Caliński, en el .bib
\usepackage{amsmath, amssymb, amsthm, mathtools, bm}
\usepackage{icomma}                  % 0,7 en modo matemático sin espacio tras la coma

% ---- Idioma -----------------------------------------------------------------
% es-tabla: «Tabla» en lugar de «Cuadro»; es-noshorthands evita conflictos de «"» y «~»;
% es-nodecimaldot conserva el punto en modo matemático (la coma decimal la da icomma).
\usepackage[spanish, es-tabla, es-noshorthands, es-nodecimaldot]{babel}
\usepackage{csquotes}
\DeclareQuoteStyle{spanish}{\guillemotleft}{\guillemotright}{\textquotedblleft}{\textquotedblright}

% ---- Interlineado 1,5 (equivalente al w:line="360" de la plantilla Word) -------
\usepackage{setspace}
\onehalfspacing
\setlength{\parindent}{0pt}
\setlength{\parskip}{6pt plus 2pt minus 1pt}
\setlength{\emergencystretch}{3em}   % evita líneas desbordadas en párrafos con términos largos

% ---- Encabezados y pie: número de página centrado, como la plantilla ----------
\usepackage{fancyhdr}
\pagestyle{fancy}
\fancyhf{}
\fancyhead[R]{\small\nouppercase{\leftmark}}
\fancyfoot[C]{\thepage}
\renewcommand{\headrulewidth}{0pt}
\fancypagestyle{plain}{\fancyhf{}\fancyfoot[C]{\thepage}}

% ---- Títulos ----------------------------------------------------------------
\usepackage{titlesec}
\titleformat{\chapter}[display]{\normalfont\bfseries\Large}
  {\MakeUppercase{\chaptertitlename}\ \thechapter}{0.6em}{\LARGE}
\titlespacing*{\chapter}{0pt}{0pt}{1.5em}
\titleformat{\section}{\normalfont\bfseries\large}{\thesection}{0.8em}{}
\titleformat{\subsection}{\normalfont\bfseries\normalsize}{\thesubsection}{0.8em}{}
\titleformat{\subsubsection}{\normalfont\bfseries\itshape\normalsize}{\thesubsubsection}{0.8em}{}
\setcounter{secnumdepth}{3}
\setcounter{tocdepth}{2}

% ---- Tablas, figuras y listas ------------------------------------------------
\usepackage{booktabs, longtable, array}
\usepackage{graphicx}
\graphicspath{{../reports/figures/}{../reports/figures/niveles/}}
\usepackage[font=small, labelfont=bf]{caption}
\usepackage{enumitem}
\setlist{itemsep=2pt, topsep=4pt}
\usepackage{xcolor}

% ---- Bibliografía: APA 7 con biblatex-apa y el .bib compartido con el informe --
\usepackage[style=apa, backend=biber]{biblatex}
\addbibresource{../tesis/referencias.bib}

% ---- Enlaces y referencias cruzadas (hyperref antes de cleveref) --------------
\usepackage[hidelinks, unicode, pdfusetitle]{hyperref}
\usepackage[spanish, noabbrev]{cleveref}

% ---- Entornos de amsthm, numerados por capítulo -----------------------------
\theoremstyle{plain}
\newtheorem{teorema}{Teorema}[chapter]
\newtheorem{proposicion}[teorema]{Proposición}
\newtheorem{lema}[teorema]{Lema}
\newtheorem{corolario}[teorema]{Corolario}
\theoremstyle{definition}
\newtheorem{definicion}[teorema]{Definición}
\newtheorem{supuesto}[teorema]{Supuesto}
\newtheorem{algoritmo}[teorema]{Algoritmo}
\theoremstyle{remark}
\newtheorem{observacion}[teorema]{Observación}
\renewcommand{\qedsymbol}{$\blacksquare$}
\addto\captionsspanish{\renewcommand{\proofname}{Demostración}}

\crefname{teorema}{teorema}{teoremas}
\Crefname{teorema}{Teorema}{Teoremas}
\crefname{proposicion}{proposición}{proposiciones}
\Crefname{proposicion}{Proposición}{Proposiciones}
\crefname{lema}{lema}{lemas}
\Crefname{lema}{Lema}{Lemas}
\crefname{corolario}{corolario}{corolarios}
\Crefname{corolario}{Corolario}{Corolarios}
\crefname{definicion}{definición}{definiciones}
\Crefname{definicion}{Definición}{Definiciones}
\crefname{supuesto}{supuesto}{supuestos}
\Crefname{supuesto}{Supuesto}{Supuestos}
\crefname{algoritmo}{algoritmo}{algoritmos}
\Crefname{algoritmo}{Algoritmo}{Algoritmos}
\crefname{observacion}{observación}{observaciones}
\Crefname{observacion}{Observación}{Observaciones}
\crefname{equation}{ecuación}{ecuaciones}
\Crefname{equation}{Ecuación}{Ecuaciones}
\creflabelformat{equation}{#2(#1)#3}
\crefname{chapter}{capítulo}{capítulos}
\Crefname{chapter}{Capítulo}{Capítulos}
\crefname{section}{sección}{secciones}
\Crefname{section}{Sección}{Secciones}
\crefname{subsection}{subsección}{subsecciones}
\Crefname{subsection}{Subsección}{Subsecciones}
\crefname{enumi}{objetivo}{objetivos}
\Crefname{enumi}{Objetivo}{Objetivos}
\numberwithin{equation}{chapter}

% ---- Marcas editoriales -----------------------------------------------------
% \pendiente{qué falta}: obra sin ficha en papers/ o cifra aún no generada. Se resalta en
%   el PDF y cada uso se lista en la nota del redactor y en tesis/registro_verificacion.md.
% \verificar{motivo}: obra con ficha y entrada en el .bib, pero con un dato por verificar
%   (equivale a la marca [VERIFICAR: motivo] del informe). Va como nota al pie.
\newcommand{\pendiente}[1]{%
  \textcolor{red!75!black}{\textbf{[PENDIENTE:}~\textit{#1}\textbf{]}}}
\newcommand{\verificar}[1]{%
  \footnote{\textcolor{orange!70!black}{\textbf{VERIFICAR:}}~#1}}

% =============================================================================
% notacion.tex · Notación única de la tesis (macros y lista de símbolos)
% Regla: antes de introducir un símbolo se comprueba que no exista aquí. Si se cambia
% uno, se actualiza en toda la tesis y en el informe (tesis/), que deriva de esta lista.
% =============================================================================

% ---- Conjuntos y operadores --------------------------------------------------
\newcommand{\RR}{\mathbb{R}}                       % números reales
\newcommand{\NN}{\mathbb{N}}                       % números naturales
\newcommand{\T}{\top}                              % transpuesta
\DeclareMathOperator{\tr}{tr}                      % traza
\DeclareMathOperator{\diag}{diag}                  % matriz diagonal / diagonal de una matriz
\DeclareMathOperator{\rango}{rango}                % rango de una matriz
\DeclareMathOperator*{\argmin}{arg\,min}
\DeclareMathOperator*{\argmax}{arg\,max}
\DeclareMathOperator{\RV}{RV}                      % coeficiente RV
\DeclareMathOperator{\CH}{CH}                      % índice de Caliński y Harabasz
\DeclareMathOperator{\KMO}{KMO}                    % medida de Kaiser-Meyer-Olkin
\DeclareMathOperator{\congr}{cong}                 % congruencia de Tucker
\DeclareMathOperator{\ctr}{ctr}                    % contribución a un eje
\DeclareMathOperator{\qlt}{cal}                    % calidad de representación en un plano
\newcommand{\norm}[1]{\lVert #1\rVert}             % norma euclídea
\newcommand{\normF}[1]{\lVert #1\rVert_{F}}        % norma de Frobenius
\newcommand{\prodF}[2]{\langle #1,#2\rangle_{F}}   % producto interno de Frobenius
\newcommand{\abs}[1]{\lvert #1\rvert}
\newcommand{\prom}[1]{\langle #1\rangle}           % media temporal de una sucesión en un tramo
\DeclareMathOperator{\vart}{var}                   % varianza temporal de una sucesión en un tramo

% ---- Datos -------------------------------------------------------------------
\newcommand{\mX}{\mathbf{X}}       % matriz de datos observados, n x p
\newcommand{\mZ}{\mathbf{Z}}       % matriz de datos centrados y estandarizados, n x p
\newcommand{\vx}{\mathbf{x}}       % fila de X (vector del mes t)
\newcommand{\vz}{\mathbf{z}}       % fila o columna de Z
\newcommand{\mI}{\mathbf{I}}       % identidad
\newcommand{\vuno}{\mathbf{1}}     % vector de unos

% ---- PCA ----------------------------------------------------------------------
\newcommand{\mR}{\mathbf{R}}           % matriz de correlaciones
\newcommand{\mV}{\mathbf{V}}           % vectores propios de R
\newcommand{\vv}{\mathbf{v}}           % vector propio
\newcommand{\mLambda}{\boldsymbol{\Lambda}} % valores propios en diagonal
\newcommand{\mE}{\mathbf{E}}           % matriz p x d con columnas ortonormales (candidata)
\newcommand{\mP}{\mathbf{P}}           % proyector ortogonal
\newcommand{\vc}{\mathbf{c}}           % componente principal (scores)
\newcommand{\mM}{\mathbf{M}}           % aplicación lineal de reconstrucción (rango <= d)
\newcommand{\mO}{\mathbf{O}}           % matriz ortogonal
\newcommand{\mN}{\mathbf{N}}           % matriz genérica (teorema de Eckart-Young)

% ---- EPCA ---------------------------------------------------------------------
\newcommand{\nw}{n_{w}}                % ancho de la ventana
\newcommand{\Wt}{\mathcal{W}}          % conjunto de meses de una ventana
\newcommand{\vm}{\mathbf{m}}           % media ponderada de la ventana
\newcommand{\mC}{\mathbf{C}}           % covarianza ponderada de la ventana
\newcommand{\mD}{\mathbf{D}}           % diagonal de C

% ---- HFC ----------------------------------------------------------------------
\newcommand{\gA}{\mathcal{A}}          % grupo de variables
\newcommand{\gB}{\mathcal{B}}          % grupo de variables
\newcommand{\fac}{\boldsymbol{\varphi}}% factor (representante) de un grupo
\newcommand{\dif}{\boldsymbol{\psi}}   % componente de diferencias de un nodo

% ---- Partición y comparación -------------------------------------------------
\newcommand{\vy}{\mathbf{y}}           % punto de una trayectoria a particionar
\newcommand{\mF}{\mathbf{F}}           % configuración (RV, Procrustes)
\newcommand{\estr}{\mathcal{S}}        % estructura estimada de un conjunto de meses
\newcommand{\interv}{\mathcal{I}}      % intervalo de un régimen

% ---- Episodios de referencia y validación externa (iteración 3b) --------------
\newcommand{\epi}{\mathsf{e}}          % episodio de referencia
\newcommand{\Epis}{\mathcal{E}}        % conjunto de episodios de referencia
\newcommand{\tini}[1]{\underline{t}_{#1}} % mes de inicio de un episodio
\newcommand{\tfin}[1]{\overline{t}_{#1}}  % mes de fin (efectivo) de un episodio

% ---- Ventanas, nulos y criterios de la parte dinámica (iteración 3b) ---------
\newcommand{\nomega}{n_{\omega}}       % n efectivo de Kish de los pesos de una ventana
\newcommand{\nhorn}{n_{\mathrm{H}}}    % filas del análisis paralelo de Horn por ventana
\newcommand{\Kest}{K^{\ast}}           % número de segmentos (clases) elegido por un criterio
\newcommand{\Kef}{K_{\mathrm{ef}}}     % mayor número de segmentos factible con n_min
\newcommand{\Lb}{L_{\mathrm{b}}}       % longitud de bloque del remuestreo (no confundir con L del LSTM)
\newcommand{\mb}{m_{\mathrm{b}}}       % número de bloques de una réplica
\newcommand{\RVad}[1]{\RV^{\mathrm{ad}}_{#1}}            % RV entre ventanas adyacentes no solapadas
\newcommand{\umbralRV}[1]{\underline{\RV}^{\,\mathrm{ad}}_{#1}} % p5 por posición del nulo
\newcommand{\RVadr}[2]{\RV^{\mathrm{ad},(#2)}_{#1}}       % RV adyacente de la réplica #2
\newcommand{\umbralRVr}[2]{\underline{\RV}^{\,\mathrm{ad},(#2)}_{#1}} % p5 sin la réplica
\newcommand{\ind}{\mathbb{I}}          % función indicadora
\newcommand{\pib}{\pi^{\mathrm{b}}}    % sucesión de índices por bloques (no confundir con la constante pi)
\DeclareMathOperator{\fr}{fr}          % frontera de un par de ventanas adyacentes
\DeclareMathOperator{\red}{red}        % reducción relativa del costo de Fisher
\newcommand{\pval}{\mathrm{pv}}        % p-valor de Monte Carlo
\newcommand{\paso}{\mathsf{s}}         % paso entre ventanas de la regla del 06

% ---- Aprendizaje profundo ----------------------------------------------------
\newcommand{\mA}{\mathbf{A}}           % codificador lineal
\newcommand{\mB}{\mathbf{B}}           % decodificador lineal
\newcommand{\veta}{\boldsymbol{\eta}}  % código latente
\newcommand{\vh}{\mathbf{h}}           % estado oculto del LSTM
\newcommand{\vcel}{\mathbf{s}}         % estado de la celda de memoria del LSTM
\newcommand{\vig}{\mathbf{g}^{\mathrm{e}}} % compuerta de entrada
\newcommand{\vfg}{\mathbf{g}^{\mathrm{o}}} % compuerta de olvido
\newcommand{\vog}{\mathbf{g}^{\mathrm{s}}} % compuerta de salida
\newcommand{\mW}{\mathbf{W}}           % pesos de entrada del LSTM
\newcommand{\mU}{\mathbf{U}}           % pesos recurrentes del LSTM
\newcommand{\vb}{\mathbf{b}}           % sesgo

% =============================================================================
% Lista de símbolos (se imprime en las páginas preliminares con \listadesimbolos)
% =============================================================================
\newcommand{\gruposimbolos}[1]{\multicolumn{2}{@{}l}{\textbf{#1}}\\[2pt]}
\newcommand{\listadesimbolos}{%
\chapter*{Lista de símbolos}
\addcontentsline{toc}{chapter}{Lista de símbolos}
\markboth{Lista de símbolos}{Lista de símbolos}
{\small
\begin{longtable}{@{}p{0.24\textwidth}p{0.70\textwidth}@{}}
\toprule
Símbolo & Significado \\
\midrule
\endhead
\bottomrule
\endfoot
\gruposimbolos{Generales}
$\RR$, $\NN$, $\mathbb{Z}$ & Números reales, naturales y enteros. \\
$\mI_p$, $\vuno$ & Matriz identidad de orden $p$; vector de unos. \\
$\mathbf{a}^{\T}$, $\tr$, $\diag$, $\rango$ & Transpuesta, traza, matriz diagonal (o diagonal de una matriz) y rango. \\
$\norm{\cdot}$, $\normF{\cdot}$, $\prodF{\cdot}{\cdot}$ & Norma euclídea; norma y producto interno de Frobenius, $\prodF{\mathbf{G}}{\mathbf{H}}=\tr(\mathbf{G}^{\T}\mathbf{H})$. \\
$\lfloor\cdot\rfloor$, $\lceil\cdot\rceil$ & Parte entera inferior y superior. \\
$\abs{\cdot}$, $\operatorname{sgn}$ & Valor absoluto o módulo de un número complejo; función signo. \\
$\mathrm{i}$, $e$ & Unidad imaginaria; base de los logaritmos naturales ($\ln$). \\
$O(\cdot)$ & Orden del número de operaciones de un algoritmo. \\
$\ind\{\cdot\}$, $\Pr$, $\mathbb{E}$, $\#$ & Función indicadora; probabilidad; esperanza; número de elementos de un conjunto. \\
$\mathbf{a}$, $\mathbf{b}$, $\xi$ & Vectores y escalares genéricos en enunciados y demostraciones. \\
$\mathbf{G}\succeq 0$ & Matriz simétrica semidefinida positiva. \\
$\odot$, $\sigma(\cdot)$, $\tanh$ & Producto elemento a elemento; función logística; tangente hiperbólica. \\[4pt]
\gruposimbolos{Datos}
$n$, $p$ & Número de meses (observaciones) y de series (variables); en el análisis oficial, $n=\nMesesNiveles$ y $p=\nSeries$ (\cref{sec:met-datos}). \\
$t$, $s$ & Índices de mes, $t,s\in\{1,\dots,n\}$. \\
$j$, $j'$ & Índices de serie, $j,j'\in\{1,\dots,p\}$. \\
$\mX$, $\vx_t$, $x_{tj}$ & Matriz de datos $n\times p$ (series sin transformar o transformadas); su fila $t$; su elemento. \\
$\mZ$, $\vz_t$, $\vz_j$, $z_{tj}$ & Datos centrados y estandarizados con divisor $n$; fila $t$; columna $j$; elemento. \\
$\bar x_j$, $\hat\sigma_j$ & Media y desviación típica (divisor $n$) de la serie $j$. \\
$\Delta$ & Primera diferencia, $\Delta x_t=x_t-x_{t-1}$; para vectores, $\Delta\vx_t=\vx_t-\vx_{t-1}$. \\
$\hat x_t$, $\hat\epsilon_t$ & Recta de mínimos cuadrados de una serie sobre el tiempo y sus residuos (regla D5, nulo mixto). \\
$P$, $\vartheta$, $\gamma(P)$ & Período de una oscilación en meses; frecuencia angular $2\pi/P$; ganancia de $\Delta$. \\
$\varrho$, $\psi_0$ & Amplitud y fase de una oscilación. \\
$t_0$, $T$ & Mes anterior al primero de un tramo de meses consecutivos y longitud del tramo (\cref{prop:ganancia}). \\
$\prom{y}_{t_0,T}$, $\vart_{t_0,T}(y)$ & Media y varianza temporales de la sucesión $(y_t)$ en los meses $t_0+1,\dots,t_0+T$. \\
$y$, $\varrho_y$, $\psi_y$ & Oscilación genérica, su amplitud y su fase, en la demostración de la \cref{prop:ganancia}. \\[4pt]
\gruposimbolos{Análisis de componentes principales}
$\mR$, $r_{jj'}$ & Matriz de correlaciones $\mZ^{\T}\mZ/n$ y su elemento. \\
$\lambda_\ell$, $\vv_\ell$, $v_{j\ell}$ & Valor propio $\ell$ de $\mR$ (en orden decreciente), vector propio unitario asociado y su coordenada $j$. \\
$\ell$, $d$ & Índice de componente; número de componentes retenidas. \\
$\mV$, $\mLambda$, $\mV_d$ & Matriz de vectores propios; matriz diagonal de valores propios; primeras $d$ columnas de $\mV$. \\
$\mE$, $\mP$ & Matriz $p\times d$ con columnas ortonormales; proyector ortogonal $\mE\mE^{\T}$. \\
$J(\mE)$ & Error cuadrático medio de reconstrucción del PCA. \\
$w_\ell$ & En la demostración de la \cref{prop:pca}, $\vv_\ell^{\T}\mP\vv_\ell$. \\
$\vc_\ell$, $c_{t\ell}$ & Componente principal $\ell$ ($\mZ\vv_\ell$) y su valor en el mes $t$. \\
$\rho_{j\ell}$ & Correlación entre la serie $j$ y la componente $\ell$. \\
$\cos^2_{j\ell}$, $\ctr_{j\ell}$, $\ctr_{t\ell}$ & Calidad de representación de la serie $j$ en el eje $\ell$; contribución de la serie $j$ y del mes $t$ al eje $\ell$. \\
$\qlt_j(\ell,\ell')$ & Calidad de representación de la serie $j$ en el plano $(\ell,\ell')$. \\
$\mN$, $\varsigma_\ell$ & Matriz genérica que se aproxima (teorema de Eckart y Young) y sus valores singulares. \\
$\mM$, $\mO$ & Aplicación lineal de reconstrucción de rango a lo sumo $d$; matriz ortogonal. \\
$\lambda^{\mathrm{H}}_\ell$, $\lambda^{\mathrm{H}}_\ell(n',p)$, $\beta_\ell$ & Umbral del análisis paralelo de Horn, en general o con una tabla de $n'$ filas y $p$ columnas; valor esperado de la vara quebrada. \\
$r^{jj'}$, $q_{jj'}$ & Elemento de $\mR^{-1}$; correlación anti-imagen (parcial) entre las series $j$ y $j'$. \\
$\KMO$, $\KMO_j$ & Medida global e individual de Kaiser-Meyer-Olkin. \\
$\chi^2_{\mathrm{B}}$ & Estadístico de esfericidad de Bartlett. \\[4pt]
\gruposimbolos{Clasificación Factorial Jerárquica}
$\gA$, $\gB$ & Grupos (clases) de variables. \\
$\fac_{\gA}$ & Factor (representante) del grupo $\gA$, vector de $\RR^n$. \\
$\dif_{\gA\cup\gB}$ & Componente de diferencias del nodo que fusiona $\gA$ y $\gB$. \\
$\mu_1^{\gA\gB}\ge\mu_2^{\gA\gB}$ & Valores propios de la matriz de covarianzas de $[\fac_{\gA}\ \fac_{\gB}]$; $\mu_2^{\gA\gB}$ es el índice de fusión. \\
$(\alpha,\beta)$ & Vector propio unitario asociado a $\mu_1^{\gA\gB}$ (coeficientes de la fusión). \\
$\kappa_{\gA}$, $\kappa_{\gB}$, $\kappa_{\gA\gB}$, $\mathbf{K}_{\gA\gB}$ & Varianzas de $\fac_{\gA}$ y $\fac_{\gB}$, su covarianza y la matriz $2\times2$ que forman. \\
$\mu_1^{\text{raíz}}$ & Primer valor propio de la última fusión. \\
$\mR_{\gA}$, $\lambda_1(\mR_{\gA})$, $\mathbf{w}^{\gA}$ & Submatriz de $\mR$ de las series de $\gA$; su mayor valor propio; coeficientes de $\fac_{\gA}$ sobre las $\vz_j$, $j\in\gA$. \\[4pt]
\gruposimbolos{Análisis evolutivo}
$\nw$ & Ancho de la ventana móvil, en meses. \\
$N$ & Número de ventanas, $N=n-\nw+1$; en general, longitud de una trayectoria. \\
$i$, $t_i$ & Índice de la ventana que empieza en el mes $i$; mes al que se asigna, $t_i=i+\lfloor\nw/2\rfloor$. \\
$\Wt_t$ & Conjunto de meses de la ventana asignada al mes $t$. \\
$u$, $\omega_u$ & Posición dentro de la ventana, $u=1,\dots,\nw$, y su peso. \\
$\nu$ & Parámetro de forma de los pesos gaussianos (argumento \texttt{nsd} de \texttt{aede}). \\
$\vm_t$, $\mC_t$, $\mD_t$, $\mR_t$ & Media, covarianza, diagonal de la covarianza y correlación ponderadas de la ventana. \\
$\lambda_\ell(t)$, $\vv_\ell(t)$ & Valor y vector propio $\ell$ de $\mR_t$. \\
$\tilde\omega_u$, $\tilde\vz_s$ & Peso sin normalizar; mes $s$ estandarizado con la media y la varianza ponderadas. \\
$\mathbf{G}_b$ & Matriz diagonal de factores de escala $b_j>0$ (con desplazamientos $a_j$). \\
$\nomega$ & Tamaño efectivo de Kish de los pesos de la ventana, $(\sum_u\omega_u^2)^{-1}$ (\cref{def:kish}). \\
$c(\nu)$, $a_\nu$ & Límite de $\nomega/\nw$ con pesos gaussianos y $a_\nu=\nu/\sqrt2$ (\cref{prop:kish}). \\
$\xi_u$, $S_k$, $I_k$ & Puntos equiespaciados, sumas de Riemann e integrales de la demostración de la \cref{prop:kish}. \\
$\nhorn$ & Filas del análisis paralelo por ventana, $\lceil\nomega\rceil$ (\cref{def:umbral-ventana}). \\
$\paso$, $M$, $\ell_0$ & Paso entre ventanas, número de ventanas y eje que sigue la regla del 06 (\cref{def:regla06}). \\
$\mathbf{a}_k$, $\tilde{\mathbf{a}}_k$ & Vector propio unitario del eje $\ell_0$ en la ventana $k$ de la regla del 06, sin alinear y alineado. \\
$\rho(\cdot,\cdot)$, $\operatorname{sgn}^{+}$ & Correlación de Pearson entre las coordenadas de dos vectores; signo con $\operatorname{sgn}^{+}(0)=1$. \\
$\rho^{(06)}_k$, $\bar\rho$, $s_\rho$ & Estabilidad de la regla del 06 en la ventana $k$; su media y su desviación típica. \\
$c^{\mathrm{ent}}_{\epi}$, $c^{\mathrm{sal}}_{\epi}$ & Centros de entrada y de salida del episodio $\epi$ en la ventana (\cref{def:borde}). \\[4pt]
\gruposimbolos{Partición de Fisher y número de clases}
$\vy_t\in\RR^{q}$, $q$ & Punto $t$ de la trayectoria que se particiona; su dimensión. \\
$K$, $k$ & Número de segmentos (clases); índice de segmento. \\
$\tau_k$ & Último punto del segmento $k$, con $0=\tau_0<\tau_1<\dots<\tau_K=N$. \\
$n_{\min}$ & Longitud mínima de un segmento. \\
$Q(a,b)$, $\bar\vy_{a:b}$ & Suma de cuadrados intrasegmento de los puntos $a,\dots,b$ y su media. \\
$Q^\ast_K(b)$ & Mínimo de la suma de cuadrados intra al dividir los puntos $1,\dots,b$ en $K$ segmentos. \\
$T$ & Suma de cuadrados total de la trayectoria. \\
$K_{\max}$, $Q_K$ & Número máximo de segmentos; $Q^\ast_K(N)$. \\
$\Kef$ & Mayor número de segmentos factible, $\min(K_{\max},\lfloor N/n_{\min}\rfloor)$ (\cref{cor:fisher-factible}). \\
$\Kest$, $\Kest_{\mathrm{codo}}$, $\Kest_{\mathrm K}$ & Número de segmentos que elige un criterio (el $k^\ast$ del código); el del codo por caída máxima y el de la regla K (\cref{def:criterios-k}). \\
$\red_K$, $\upsilon$ & Reducción relativa del costo al pasar de $K$ a $K+1$ segmentos; umbral de la regla K. \\
$w_{\mathrm m}$, $\bar c_{t\ell}$ & Longitud de la media móvil causal y \emph{score} suavizado de la datación PCA-MM12 (\cref{def:datacion}). \\
$\CH^{\mathrm{h}}_K$, $\CH^{\mathrm{a}}_K$ & Índice de Caliński y Harabasz en su forma habitual y en la variante de \texttt{aede}. \\
$K^{\mathrm h}$, $K^{\mathrm a}$ & Número de segmentos elegido por cada variante. \\[4pt]
\gruposimbolos{Régimen, transición y comparación de estructuras}
$\estr(\cdot)$ & Estructura estimada de un conjunto de meses. \\
$\tau$ & Mes de una transición (\cref{def:transicion}); no confundir con $\tau_k$. \\
$\Wt_\tau^{-}$, $\Wt_\tau^{+}$ & Conjuntos de meses contiguos adyacentes, antes y después de $\tau$. \\
$\delta(\cdot,\cdot)$, $\varepsilon$ & Disimilitud entre estructuras; tolerancia. \\
$\interv_k$, $\Wt$, $\Wt'$ & Intervalo de un régimen; conjuntos de meses contiguos (ventanas). \\
$\mF_1$, $\mF_2$, $m$ & Configuraciones no nulas de las mismas $m$ filas (centradas o no, \cref{def:rv}). \\
$\mF_t$, $\mR^{[d]}_t$ & Matriz $p\times d$ de cargas de la ventana de centro $t$ y su producto $\mF_t\mF_t^{\T}$ (\cref{def:rv-cargas}). \\
$\RV^{(d)}(t,t')$, $\rho_{\mR}(t,t')$ & RV de cargas sin centrar con $d$ ejes; coeficiente entre las matrices de correlación completas. \\
$\fr(i)$, $\RVad{i}$ & Frontera y RV del par $i$ de ventanas adyacentes no solapadas (\cref{def:pares}). \\
$\RV(\mF_1,\mF_2)$ & Coeficiente RV. \\
$c$ & Constante de proporcionalidad positiva en la \cref{prop:rv}(a). \\
$a_1$, $a_2$ & Factores de escala no nulos en la \cref{prop:rv}(b). \\
$\mathbf{G}$, $\mathbf{H}$ & En la demostración de la \cref{prop:rv}, $\mF_1\mF_1^{\T}$ y $\mF_2\mF_2^{\T}$. \\
$\mO_1$, $\boldsymbol{\Sigma}$, $\mO_2$, $\boldsymbol{\Theta}$ & Descomposición en valores singulares de $\mF_1^{\T}\mF_2$; matriz ortogonal auxiliar de la \cref{prop:procrustes}. \\
$\mO^\ast$ & Rotación de Procrustes. \\
$\congr(\mathbf{a},\mathbf{b})$ & Coeficiente de congruencia de Tucker entre dos vectores. \\[4pt]
\gruposimbolos{Episodios de referencia y nulos}
$\epi$, $\Epis$ & Episodio de referencia; conjunto de episodios. \\
$\tini{\epi}$, $\tfin{\epi}$ & Mes de inicio y mes de fin (efectivo) del episodio $\epi$ (\cref{def:episodio}). \\
$h$ & Tolerancia en meses (\cref{def:tolerancia}). \\
$\tau^\circ$ & Mes de una transición que se quiere confirmar (\cref{def:confirmada}). \\
$\Lb$, $\mb$, $\iota_k$ & Longitud de bloque, número de bloques e inicios de bloque (\cref{def:bloques}); $\Lb$ no es la $L$ del LSTM. \\
$\pib(t)$, $n'$ & Sucesión de índices por bloques; su longitud. \\
$\dot x_t$, $\overline{\Delta x}_j$ & Incrementos centrados y media de los incrementos de la serie $j$ (\cref{def:hac}). \\
$\hat\Gamma_k$, $m_{\mathrm{HAC}}$, $\bar m$ & Autocovarianza de rezago $k$ de los incrementos; número de rezagos; $m_{\mathrm{HAC}}+1$. \\
$\hat\sigma^2_{\mathrm{HAC},j}$, $t^{\mathrm{HAC}}_j$ & Varianza de largo plazo y estadístico $t$ de la deriva de la serie $j$. \\
$x^\ast_{tj}$ & Serie simulada del nulo de tendencias mixto (\cref{def:nulo-mixto}). \\
$B$, $\zeta$, $\mX^{(\zeta)}$ & Número de réplicas; índice de réplica; réplica del nulo de estructura (\cref{def:nulo-estructura}). \\
$\RVadr{i}{\zeta}$, $\umbralRV{i}$, $\umbralRVr{i}{-\zeta}$ & RV adyacente de la réplica $\zeta$; cuantil de 5\,\% por posición; el mismo sin la réplica $\zeta$ (\cref{def:oe4}). \\
$\varpi_0$, $\varpi_\zeta$, $\pval$ & Proporción de pares bajo el umbral, observada y de la réplica $\zeta$; p-valor del criterio de OE4. \\
$\mathcal{P}_j$, $\tilde\varpi_j$, $\widetilde{\pval}$, $\mathcal{J}$, $\alpha_0$ & Perfiles, proporciones y p-valor de la versión simétrica; conjunto y nivel auxiliares (\cref{prop:oe4}). \\
$\mathcal{T}_a$, $E_t$, $E^{(95)}_a$ & Meses de entrenamiento del año $a$; error de reconstrucción del mes $t$; su percentil 95 en $\mathcal{T}_a$ (\cref{def:walkforward}). \\[4pt]
\gruposimbolos{Aprendizaje profundo}
$f_\theta$, $g_\phi$ & Codificador y decodificador, con parámetros $\theta$ y $\phi$. \\
$\veta_t\in\RR^{d}$ & Código latente del mes $t$. \\
$\Omega(\theta,\phi)$ & Término de regularización. \\
$\mA$, $\mB$ & Codificador lineal $d\times p$ y decodificador lineal $p\times d$. \\
$L$ & Longitud de la ventana que lee el LSTM. \\
$\vh_t\in\RR^{n_h}$, $n_h$ & Estado oculto del LSTM y su dimensión. \\
$\vcel_t$, $\tilde\vcel_t$ & Estado de la celda de memoria y candidato. \\
$\vig_t$, $\vfg_t$, $\vog_t$ & Compuertas de entrada, de olvido y de salida. \\
$\mW_\bullet$, $\mU_\bullet$, $\vb_\bullet$ & Pesos de entrada, pesos recurrentes y sesgos de cada compuerta. \\
\end{longtable}
}%
}


% Cifras: cifras.tex lo genera herramientas/generar_cifras.py desde data/results/.
% Mientras no exista, o si le falta alguna macro, cifras_stub.tex la define como
% \pendiente{cifra} (\providecommand no sobrescribe las macros ya generadas).
\IfFileExists{cifras.tex}{% Archivo generado por herramientas/generar_cifras.py a partir de herramientas/cifras_mapa.csv.
% No editar a mano: cada cifra se lee de data/results/ (regla 3 de CLAUDE.md).
% Para regenerarlo: python3 herramientas/generar_cifras.py

% Meses de la vía en niveles | data/results/05_dimensiones.csv | via=niveles | n
\newcommand{\nMesesNiveles}{228}
% Meses de la vía transformada | data/results/05_dimensiones.csv | via=transformadas | n
\newcommand{\nMesesTransformadas}{227}
% Series del análisis (D1) | data/results/05_dimensiones.csv | via=niveles | p
\newcommand{\nSeries}{17}
% Primer mes de la vía en niveles | data/results/05_dimensiones.csv | via=niveles | fecha_inicio
\newcommand{\fechaInicioNiveles}{enero de 2007}
% Primer mes de la vía transformada | data/results/05_dimensiones.csv | via=transformadas | fecha_inicio
\newcommand{\fechaInicioTransformadas}{febrero de 2007}
% Último mes de los datos (ambas vías) | data/results/05_dimensiones.csv | via=niveles | fecha_fin
\newcommand{\fechaFinDatos}{diciembre de 2025}
% Año inicial del período de estudio | data/results/05_dimensiones.csv | via=niveles | fecha_inicio
\newcommand{\anioInicioDatos}{2007}
% Año final del período de estudio | data/results/05_dimensiones.csv | via=niveles | fecha_fin
\newcommand{\anioFinDatos}{2025}
% Inercia (%) de PC1 en niveles | data/results/niveles/05_pca_valores_propios.csv | componente=1 | inercia
\newcommand{\inerciaPCUnoNiveles}{67{,}2}
% Inercia (%) de PC2 en niveles | data/results/niveles/05_pca_valores_propios.csv | componente=2 | inercia
\newcommand{\inerciaPCDosNiveles}{18{,}3}
% Inercia (%) de PC3 en niveles | data/results/niveles/05_pca_valores_propios.csv | componente=3 | inercia
\newcommand{\inerciaPCTresNiveles}{7{,}1}
% Inercia (%) de PC4 en niveles | data/results/niveles/05_pca_valores_propios.csv | componente=4 | inercia
\newcommand{\inerciaPCCuatroNiveles}{2{,}9}
% Inercia (%) de PC1 en transformadas | data/results/05_pca_valores_propios.csv | componente=1 | inercia
\newcommand{\inerciaPCUnoTransformadas}{23{,}3}
% Inercia (%) de PC2 en transformadas | data/results/05_pca_valores_propios.csv | componente=2 | inercia
\newcommand{\inerciaPCDosTransformadas}{16{,}0}
% Inercia (%) de PC3 en transformadas | data/results/05_pca_valores_propios.csv | componente=3 | inercia
\newcommand{\inerciaPCTresTransformadas}{11{,}2}
% Inercia (%) de PC4 en transformadas | data/results/05_pca_valores_propios.csv | componente=4 | inercia
\newcommand{\inerciaPCCuatroTransformadas}{8{,}6}
% Componentes retenidos por Horn en niveles | data/results/niveles/05_pca_retencion.csv | criterio=horn | componentes
\newcommand{\hornNiveles}{2}
% Componentes retenidos por Horn en transformadas | data/results/05_pca_retencion.csv | criterio=horn | componentes
\newcommand{\hornTransformadas}{4}
% Componentes con valor propio > 1 en niveles | data/results/niveles/05_pca_retencion.csv | criterio=kaiser | componentes
\newcommand{\kaiserNiveles}{3}
% Componentes con valor propio > 1 en transformadas | data/results/05_pca_retencion.csv | criterio=kaiser | componentes
\newcommand{\kaiserTransformadas}{6}
% KMO global en niveles | data/results/niveles/05_adecuacion.csv | ambito=global;medida=kmo_global | valor
\newcommand{\kmoGlobalNiveles}{0{,}86}
% KMO global en transformadas | data/results/05_adecuacion.csv | ambito=global;medida=kmo_global | valor
\newcommand{\kmoGlobalTransformadas}{0{,}70}
% Chi cuadrado de Bartlett con n en niveles | data/results/niveles/05_adecuacion.csv | ambito=global;medida=bartlett_chi2_n | valor
\newcommand{\chiBartlettNiveles}{12\,210}
% Chi cuadrado de Bartlett con n en transformadas | data/results/05_adecuacion.csv | ambito=global;medida=bartlett_chi2_n | valor
\newcommand{\chiBartlettTransformadas}{1\,831}
% Grados de libertad de Bartlett (p(p-1)/2) | data/results/niveles/05_adecuacion.csv | ambito=global;medida=bartlett_gl | valor
\newcommand{\glBartlett}{136}
% Chi cuadrado de Bartlett con la mediana de n_ef en niveles | data/results/niveles/05_adecuacion.csv | ambito=global;medida=bartlett_chi2_nef | valor
\newcommand{\chiBartlettNefNiveles}{162}
% p-valor de Bartlett con la mediana de n_ef en niveles | data/results/niveles/05_adecuacion.csv | ambito=global;medida=bartlett_p_nef | valor
\newcommand{\pBartlettNefNiveles}{0{,}06}
% Mediana de n_ef entre pares en niveles | data/results/niveles/05_n_efectivo_resumen.csv | version=original | n_ef_mediana
\newcommand{\nEfMedianaNiveles}{10{,}4}
% Mediana de n_ef entre pares en transformadas | data/results/05_n_efectivo_resumen.csv | version=original | n_ef_mediana
\newcommand{\nEfMedianaTransformadas}{182{,}0}
% Número de condición de R en niveles | data/results/niveles/05_adecuacion.csv | ambito=global;medida=num_condicion | valor
\newcommand{\numCondicionNiveles}{67\,647}
% Número de condición de R en transformadas | data/results/05_adecuacion.csv | ambito=global;medida=num_condicion | valor
\newcommand{\numCondicionTransformadas}{68}
% Correlación de PC1 con el tiempo en niveles | data/results/niveles/05_pca_tendencia_componentes.csv | nombre=PC1 | r_tiempo
\newcommand{\rTiempoPCUnoNiveles}{0{,}993}
% Clase D5 de PC1 en niveles | data/results/niveles/05_pca_tendencia_componentes.csv | nombre=PC1 | clase_d5
\newcommand{\claseDCincoPCUnoNiveles}{DS}
% Clase D5 de PC2 en niveles | data/results/niveles/05_pca_tendencia_componentes.csv | nombre=PC2 | clase_d5
\newcommand{\claseDCincoPCDosNiveles}{DS}
% Clase D5 de PC3 en niveles | data/results/niveles/05_pca_tendencia_componentes.csv | nombre=PC3 | clase_d5
\newcommand{\claseDCincoPCTresNiveles}{DS}
% Clase D5 de PC4 en niveles | data/results/niveles/05_pca_tendencia_componentes.csv | nombre=PC4 | clase_d5
\newcommand{\claseDCincoPCCuatroNiveles}{DS}
% Clase D5 de PC1 en transformadas | data/results/05_pca_tendencia_componentes.csv | nombre=PC1 | clase_d5
\newcommand{\claseDCincoPCUnoTransformadas}{TS}
% Clase D5 de PC2 en transformadas | data/results/05_pca_tendencia_componentes.csv | nombre=PC2 | clase_d5
\newcommand{\claseDCincoPCDosTransformadas}{TS}
% Clase D5 de PC3 en transformadas | data/results/05_pca_tendencia_componentes.csv | nombre=PC3 | clase_d5
\newcommand{\claseDCincoPCTresTransformadas}{TS}
% Clase D5 de PC4 en transformadas | data/results/05_pca_tendencia_componentes.csv | nombre=PC4 | clase_d5
\newcommand{\claseDCincoPCCuatroTransformadas}{indeterminada}
% Series en niveles clasificadas TS (D5) | data/results/niveles/05_clasificacion_d5_resumen.csv | conjunto=series_niveles | TS
\newcommand{\nSeriesTS}{1}
% Series en niveles clasificadas DS (D5) | data/results/niveles/05_clasificacion_d5_resumen.csv | conjunto=series_niveles | DS
\newcommand{\nSeriesDS}{13}
% Series en niveles con clase indeterminada (D5) | data/results/niveles/05_clasificacion_d5_resumen.csv | conjunto=series_niveles | indeterminada
\newcommand{\nSeriesIndeterminadas}{3}
% Proporción del nulo con deriva con PC1 >= observado | data/results/niveles/05_nulo_tendencias.csv | version=con_deriva;medida=inercia_pc1 | proporcion
\newcommand{\nuloConDerivaPCUno}{0{,}902}
% Proporción del nulo con deriva con PC2 >= observado | data/results/niveles/05_nulo_tendencias.csv | version=con_deriva;medida=inercia_pc2 | proporcion
\newcommand{\nuloConDerivaPCDos}{0{,}015}
% Proporción del nulo sin deriva con PC1 >= observado | data/results/niveles/05_nulo_tendencias.csv | version=sin_deriva;medida=inercia_pc1 | proporcion
\newcommand{\nuloSinDerivaPCUno}{0{,}004}
% Proporción del nulo sin deriva con PC2 >= observado | data/results/niveles/05_nulo_tendencias.csv | version=sin_deriva;medida=inercia_pc2 | proporcion
\newcommand{\nuloSinDerivaPCDos}{0{,}592}
% Proporción del nulo sin tendencia con PC1 >= observado | data/results/niveles/05_nulo_tendencias.csv | version=sin_tendencia;medida=inercia_pc1 | proporcion
\newcommand{\nuloSinTendenciaPCUno}{0{,}124}
% Proporción del nulo sin tendencia con PC2 >= observado | data/results/niveles/05_nulo_tendencias.csv | version=sin_tendencia;medida=inercia_pc2 | proporcion
\newcommand{\nuloSinTendenciaPCDos}{0{,}595}
% Mediana de la inercia de PC1 bajo el nulo con deriva | data/results/niveles/05_nulo_tendencias.csv | version=con_deriva;medida=inercia_pc1 | p50
\newcommand{\nuloConDerivaMedianaPCUno}{74{,}6}
% Inercia (%) de PC1 de las series en niveles sin tendencia lineal | data/results/niveles/05_pca_contraste_tendencia.csv | version=niveles_sin_tendencia;medida=inercia;componente=1 | valor
\newcommand{\inerciaPCUnoSinTendencia}{40{,}0}
% KMO de las series en niveles sin tendencia lineal | data/results/niveles/05_pca_contraste_tendencia.csv | version=niveles_sin_tendencia;medida=kmo | valor
\newcommand{\kmoSinTendencia}{0{,}70}
% Percentil 5 de la congruencia bootstrap de PC2 (niveles; L = 24) | data/results/niveles/05_bootstrap_congruencia.csv | L=24;componente=2 | p5
\newcommand{\congruenciaBootPCDosNiveles}{0{,}793}
% Percentil 5 de la congruencia bootstrap de PC4 (niveles; L = 24) | data/results/niveles/05_bootstrap_congruencia.csv | L=24;componente=4 | p5
\newcommand{\congruenciaBootPCCuatroNiveles}{0{,}608}
}{}
% =============================================================================
% cifras_stub.tex · Sustituto provisional de cifras.tex
% NO contiene números: cada macro del contrato con el desarrollador se define como
% \pendiente{cifra}. tesis.tex carga cifras.tex (generado por
% herramientas/generar_cifras.py a partir de data/results/) si existe, y este archivo
% si no. No editar cifras.tex a mano.
% Contrato acordado (iteración 3a): \nMesesNiveles, \nMesesTransformadas, \nSeries,
% \fechaInicioNiveles, \fechaInicioTransformadas, \fechaFinDatos.
% =============================================================================
\providecommand{\nMesesNiveles}{\pendiente{cifra}}
\providecommand{\nMesesTransformadas}{\pendiente{cifra}}
\providecommand{\nSeries}{\pendiente{cifra}}
\providecommand{\fechaInicioNiveles}{\pendiente{cifra}}
\providecommand{\fechaInicioTransformadas}{\pendiente{cifra}}
\providecommand{\fechaFinDatos}{\pendiente{cifra}}

% Macros solicitadas al desarrollador fuera del contrato inicial (nota del redactor,
% 2026-10-20): años inicial y final del período de estudio, que el capítulo 1 usa en el
% problema, el objetivo general y la HT1. Deben salir de 05_dimensiones.csv.
\providecommand{\anioInicioDatos}{\pendiente{cifra}}
\providecommand{\anioFinDatos}{\pendiente{cifra}}

% Macros de la iteración 3b (redactor LaTeX, tareas L1 a L3): fechas de los episodios de
% referencia (data/referencia/episodios_referencia.csv, D7) y parámetros de método leídos de
% data/results/final/ (n de Kish por ventana, filas de Horn, nulo de estructura).
\providecommand{\episodioCrisisInicio}{\pendiente{cifra}}
\providecommand{\episodioCrisisFin}{\pendiente{cifra}}
\providecommand{\episodioTaperInicio}{\pendiente{cifra}}
\providecommand{\episodioTaperFin}{\pendiente{cifra}}
\providecommand{\episodioMateriasInicio}{\pendiente{cifra}}
\providecommand{\episodioMateriasFin}{\pendiente{cifra}}
\providecommand{\episodioKuczynskiInicio}{\pendiente{cifra}}
\providecommand{\episodioGuerraInicio}{\pendiente{cifra}}
\providecommand{\episodioGuerraFin}{\pendiente{cifra}}
\providecommand{\episodioCovidInicio}{\pendiente{cifra}}
\providecommand{\episodioCovidFin}{\pendiente{cifra}}
\providecommand{\nOmegaGaussVeinticuatro}{\pendiente{cifra}}
\providecommand{\nOmegaGaussTreintayseis}{\pendiente{cifra}}
\providecommand{\nOmegaGaussCuarentayocho}{\pendiente{cifra}}
\providecommand{\nOmegaGaussSesenta}{\pendiente{cifra}}
\providecommand{\nHornGaussVeinticuatro}{\pendiente{cifra}}
\providecommand{\nHornGaussTreintayseis}{\pendiente{cifra}}
\providecommand{\nHornGaussCuarentayocho}{\pendiente{cifra}}
\providecommand{\nHornGaussSesenta}{\pendiente{cifra}}
\providecommand{\largoBloqueNulo}{\pendiente{cifra}}
\providecommand{\replicasNuloEstructura}{\pendiente{cifra}}
\providecommand{\largoBloqueNuloCorto}{\pendiente{cifra}}
\providecommand{\largoBloqueNuloLargo}{\pendiente{cifra}}
\providecommand{\replicasNuloSensibilidad}{\pendiente{cifra}}
% Parámetros de diseño provisionales (vuelta 1 de la auditoría 3b, N-03): salen de
% data/results/final/parametros_diseno.csv; no están en decisiones.qmd.
\providecommand{\nuGauss}{\pendiente{cifra}}
\providecommand{\pasoReglaSeis}{\pendiente{cifra}}
\providecommand{\ejeReglaSeis}{\pendiente{cifra}}
\providecommand{\ventanaMediaMovil}{\pendiente{cifra}}
\providecommand{\clasesHfcMenor}{\pendiente{cifra}}
\providecommand{\clasesHfcMayor}{\pendiente{cifra}}
\providecommand{\percentilWalkforward}{\pendiente{cifra}}


\title{Detección evolutiva de regímenes del ciclo financiero peruano mediante
  Clasificación Factorial Jerárquica y aprendizaje automático}
\author{Roberto Carlos Samaniego Salcedo}
\date{}

\begin{document}

% ---- Páginas preliminares ----------------------------------------------------
\frontmatter
% Portada provisional: el formato definitivo depende de las pautas de la UNI.
\begin{titlepage}
\centering
\onehalfspacing
{\large\bfseries UNIVERSIDAD NACIONAL DE INGENIERÍA\par}
{\large FACULTAD DE CIENCIAS\par}
\vspace{1.5cm}
{\normalsize Maestría en Ciencias en Modelización Matemática y Computacional\par}
\vspace{2.5cm}
{\Large\bfseries Detección evolutiva de regímenes del ciclo financiero peruano mediante
Clasificación Factorial Jerárquica y aprendizaje automático\par}
\vspace{2cm}
{\normalsize TESIS DE MAESTRÍA\par}
\vspace{0.3cm}
{\small\pendiente{denominación exacta del grado según las pautas de la UNI}\par}
\vspace{2cm}
{\normalsize\textbf{Tesista:} Roberto Carlos Samaniego Salcedo\par}
\vspace{0.4cm}
{\normalsize\textbf{Asesor:} Dr. Sergio Camiz\par}
\vfill
{\normalsize Lima, Perú\par}
\vspace{0.4cm}
{\small\color{red!75!black} Versión de trabajo: formato provisional, pendiente de las pautas de la UNI\par}
\end{titlepage}


\chapter*{Resumen}
\addcontentsline{toc}{chapter}{Resumen}
\markboth{Resumen}{Resumen}
\pendiente{el resumen (un solo párrafo de 250 palabras como máximo, con 2 a 4 palabras
clave) y el abstract se redactan cuando los capítulos 1 a 5 estén aprobados}

\cleardoublepage
\tableofcontents
\listadesimbolos

% ---- Cuerpo ------------------------------------------------------------------
\mainmatter
% =============================================================================
% Capítulo 1 · Introducción
% Fuente: tesis/_capitulos/01_introduccion.qmd (aprobado en la iteración 2). Se traslada
% sin afirmaciones nuevas; las cifras de los datos salen de cifras.tex.
% =============================================================================
\chapter{Introducción}
\label{cap:introduccion}

\section{Planteamiento y descripción del problema}
\label{sec:planteamiento}

El ciclo financiero es la fluctuación de mediano plazo entre fases de expansión y contracción
del crédito, de los precios de activos y de la disposición a asumir riesgos.
\textcite{drehmann_borio_tsatsaronis_2012} lo distinguen del ciclo económico por su duración, de
8 a 30 años, y documentan que sus picos coinciden con las crisis bancarias sistémicas. No se
observa de forma directa: se manifiesta en el movimiento conjunto de muchas variables, cuya
estructura de relaciones puede cambiar en el tiempo. Medirlo es, por tanto, un problema de
reducción dimensional sobre series de tiempo multivariadas.

Para el Perú, el antecedente más directo es
\textcite{camiz_cerda_atoche_2026}\verificar{manuscrito fechado el 25 de abril de 2026, sin
datos de revista, volumen ni DOI; confirmar con el Dr.~Camiz su estado de publicación.}, que
continúa la tesis de maestría de Atoche Murrieta \pendiente{Atoche Murrieta (2022), tesis de
maestría, PUCP; sin ficha}. Ese trabajo aplica el Análisis de Componentes Principales (PCA), la
Clasificación Factorial Jerárquica (HFC) y el Análisis Evolutivo de Componentes Principales
(EPCA) a 17 indicadores mensuales del Banco Central de Reserva del Perú (BCRP). Deja abiertas,
sin embargo, cuatro cuestiones: no evalúa el aprendizaje profundo, no clasifica fases del ciclo
contrastadas con episodios fechados de forma independiente, su muestra termina en 2020 y fija
el ancho de la ventana del EPCA sin un criterio objetivo.

% Fuente de las cifras: cifras.tex (herramientas/generar_cifras.py), a partir de las
% dimensiones que escribe notebooks/05_pca.ipynb en data/results/.
Esta tesis analiza \nSeries{} series mensuales del BCRP de crédito, liquidez, tasas, precios,
tipo de cambio, bolsa, riesgo soberano, actividad y reservas. El análisis oficial, sobre las
series sin transformar (decisión D1), cubre \nMesesNiveles{} meses, de \fechaInicioNiveles{} a
\fechaFinDatos{}. La vía transformada tiene \nMesesTransformadas{} meses, de
\fechaInicioTransformadas{} a \fechaFinDatos{}, porque la diferenciación pierde el primer mes.
Ambas vías se ejecutan con el mismo código (decisión D2), y la forma de combinarlas está
pendiente de discusión con el asesor.

\textbf{Problema de investigación:} ¿cómo identificar, fechar y caracterizar los regímenes del
ciclo financiero peruano entre \anioInicioDatos{} y \anioFinDatos{} a partir de la estructura
multivariada de los indicadores del BCRP, integrando los métodos exploratorios evolutivos de la
escuela de Camiz con autoencoders y redes LSTM, sin renunciar a la interpretabilidad ni al
principio de que la estructura emerge de los datos?

\section{Justificación de la investigación}
\label{sec:justificacion}

\textbf{Metodológica.} La tesis amplía la cadena PCA--HFC--EPCA de tres maneras. Fecha los
regímenes con la partición óptima de
\textcite{fisher_1958}\verificar{páginas (789--798), solo en fuentes secundarias; número y
DOI ya verificados. Ficha preliminar, texto no leído: solo se cita su aporte general.}, como hacen
\textcite{camiz_roig_2011}\verificar{el nombre del editor de las actas (O.~Papú) se infirió por
analogía; el año figura como 2010/2011 en la ficha.}. Propone un criterio objetivo para el ancho
de la ventana del EPCA. E incorpora autoencoders y redes LSTM anidados en la cadena clásica,
porque un autoencoder lineal con pérdida cuadrática genera el subespacio de los primeros
componentes principales
\parencite{bourlard_kamp_1988,baldi_hornik_1989}\verificar{Bourlard y Kamp (1988), número del
fascículo; ficha preliminar, texto no consultado. Baldi y Hornik (1989), solo el DOI; el texto
completo ya se leyó.}; véase la \cref{prop:ae-lineal}. Como los antecedentes de aprendizaje profundo
usan conjuntos de datos mucho mayores
\parencite{hinton_salakhutdinov_2006,tuhin_et_al_2025}\verificar{Hinton y Salakhutdinov (2006),
se procesó el material suplementario; Tuhin et al. (2025), lista completa de autores.},
documentar cómo se adaptan a series macrofinancieras cortas es en sí un resultado.

\textbf{Económica y social.} \textcite{drehmann_borio_tsatsaronis_2012} acuñan la noción de
«recesión inconclusa»: si la política ignora la expansión del crédito y de los precios de los
inmuebles, el ajuste se posterga y la crisis posterior resulta más severa. Un indicador de
régimen transparente y reproducible podría servir como insumo exploratorio para los análisis de
estabilidad financiera del BCRP y de la SBS.

\section{Objetivos}
\label{sec:objetivos}

\subsection{Objetivo general}

Formalizar y evaluar una metodología exploratoria evolutiva, que integre el Análisis de
Componentes Principales, la Clasificación Factorial Jerárquica y el Análisis Evolutivo de
Componentes Principales con autoencoders y redes LSTM, para detectar, fechar y caracterizar los
regímenes del ciclo financiero peruano entre \anioInicioDatos{} y \anioFinDatos{} a partir de
indicadores macrofinancieros del BCRP.

\subsection{Objetivos específicos (compromisos firmes)}
\label{sec:oe}

\begin{enumerate}[label=\textbf{OE\arabic*.}, ref=OE\arabic*, leftmargin=*]
  \item\label{oe:estado-arte} Sistematizar el estado del arte sobre el ciclo financiero y sobre
    la reducción dimensional de series de tiempo multivariadas.
  \item\label{oe:datos} Construir y versionar el conjunto de series mensuales del BCRP, con
    trazabilidad de códigos, transformaciones y exclusiones.
  \item\label{oe:clasicos} Aplicar PCA, HFC y EPCA en las dos vías del diseño y contrastar los
    resultados con los de \textcite{camiz_cerda_atoche_2026}.
  \item\label{oe:ventana} Formalizar un criterio objetivo para el ancho de ventana del EPCA,
    basado en la similitud entre ventanas consecutivas (coeficiente RV o distancia de
    Procrustes).
  \item\label{oe:robustez-hfc} Evaluar la robustez de los grupos de la HFC mediante remuestreo
    por bloques y correlación cofenética.
  \item\label{oe:dimension} Triangular el número de componentes con la vara quebrada, el
    criterio de Rencher y el análisis paralelo de Horn
    \parencite{horn_1965}\verificar{ficha preliminar: se leyó el resumen, no el texto; solo se
    cita el aporte general.} \pendiente{fuentes de la vara quebrada y del criterio de Rencher}.
  \item\label{oe:dl} Diseñar autoencoders y redes LSTM adaptados a la escala de los datos, con
    pesos atados, preentrenamiento por capas, regularización y validación por bloques temporales
    y \emph{walk-forward}.
  \item\label{oe:incertidumbre} Cuantificar la incertidumbre del error de reconstrucción
    mediante Monte Carlo Dropout
    \parencite{gal_ghahramani_2016}\verificar{editores del volumen; PDF no consultado.}.
  \item\label{oe:datacion} Fechar regímenes con la partición óptima de Fisher sobre
    trayectorias de representaciones (primer valor propio del EPCA, \emph{scores} del PCA o
    códigos latentes) y clasificarlos en expansión, contracción o transición.
  \item\label{oe:episodios} Contrastar regímenes y fases con episodios fechados por fuentes
    externas, nunca por las propias series ni por un indicador derivado de ellas.
  \item\label{oe:simulado} Comparar los métodos sobre un contexto simulado, con regímenes
    conocidos por construcción, y calibrar con ello el número de factores y de clases, el ancho
    de ventana y los pesos de los análisis evolutivos.
  \item\label{oe:comparacion} Comparar los métodos clásicos y los de aprendizaje profundo según
    interpretabilidad, estabilidad temporal, capacidad descriptiva y robustez.
\end{enumerate}

\subsection{Objetivos exploratorios condicionales}

Se abordarán si los resultados de los objetivos firmes lo justifican:

\begin{enumerate}[label=\textbf{OC\arabic*.}, ref=OC\arabic*, leftmargin=*]
  \item\label{oc:ehfc} Evaluar la viabilidad de una HFC evolutiva (EHFC) que aplique la
    jerarquía completa en cada ventana móvil, en lugar de la aplicación secuencial de HFC y EPCA
    de la literatura.
  \item\label{oc:statis} Explorar el STATIS dual como detector independiente de quiebres
    estructurales, siguiendo a
    \textcite{camiz_maulucci_roig_2010}\verificar{datos de edición de las actas (ISBN o DOI).}.
  \item\label{oc:anticipar} Explorar la capacidad del LSTM para anticipar transiciones, con
    validación estrictamente fuera de muestra y declarando su escaso poder estadístico.
  \item\label{oc:extension} Extender la muestra a años anteriores a \anioInicioDatos{}, si la
    homogeneidad de las series lo permite, para evaluar la crisis financiera global fuera de
    muestra.
\end{enumerate}

\section{Hipótesis}
\label{sec:hipotesis}

El estudio es exploratorio y no formula hipótesis contrastables en sentido estricto, por tres
razones. Los métodos centrales describen la muestra sin postular un modelo que la genere, de
modo que no hay una hipótesis nula. Los pocos episodios de referencia no dan poder estadístico.
Y los resultados no se generalizan a otras economías. Se formulan, en cambio, \textbf{hipótesis
de trabajo}, que se evalúan por la convergencia de la evidencia en las dos vías. Cada una indica
qué resultado contaría en su contra.

\begin{description}[style=nextline, leftmargin=1.2em]
  \item[HT1 (no estacionariedad estructural).] La estructura de correlación no es
    constante entre \anioInicioDatos{} y \anioFinDatos{}: hay intervalos estables separados por
    transiciones. \emph{En contra:} que las ventanas sucesivas resulten tan similares como
    cabría esperar bajo una estructura constante.
  \item[HT2 (organización en grupos).] Las series se organizan en pocos grupos
    homogéneos que se corresponden con los bloques identificados por
    \textcite{camiz_cerda_atoche_2026}: salud económica y financiera, tasas de interés y riesgo
    del mercado financiero. Como el conjunto de series no es idéntico, la correspondencia se
    evalúa por bloque. \emph{En contra:} grupos inestables bajo remuestreo o sin
    correspondencia con esos bloques.
  \item[HT3 (correspondencia con episodios).] Los regímenes se corresponden,
    dentro de una tolerancia declarada de antemano, con episodios fechados externamente.
    \emph{En contra:} cortes sin relación con ningún episodio, o que dependan de la regla que
    fija su número.
  \item[HT4 (complementariedad).] Los autoencoders y las redes LSTM aportan
    información sobre la dinámica y la atipicidad del sistema que el PCA no capta por
    construcción. \emph{En contra:} que no mejoren a su anclaje lineal en ningún criterio fuera
    de muestra, o que sus resultados cambien con la semilla aleatoria.
\end{description}

% =============================================================================
% Capítulo 2 · Antecedentes y marco teórico
% Fuente: tesis/_capitulos/02_marco_teorico.qmd (aprobado), llevado al rigor de la tesis:
% definiciones, supuestos y proposiciones. Para la EPCA ponderada, la fuente de verdad es
% R/camiz_port.R (aede_r, .bstatw_ac, .ustatw_ac). Las demostraciones marcadas como
% propias no dependen de ninguna cita. Ninguna cifra de los datos figura en este capítulo.
% =============================================================================
\chapter{Antecedentes y marco teórico}
\label{cap:marco}

Este capítulo reúne los antecedentes y formaliza los métodos que usa la tesis. La
\cref{sec:antecedentes} resume el antecedente peruano. Las
\cref{sec:ciclo,sec:epistemologia} fijan el objeto de estudio y la posición epistemológica. Las
\cref{sec:datos,sec:pca,sec:hfc,sec:epca,sec:fisher,sec:comparacion,sec:nulos,sec:dl} formalizan
los métodos: cada uno con su definición, sus supuestos y, cuando corresponde, una proposición. Las
\cref{sec:criterios,sec:vacios} cierran con los criterios de comparación y los vacíos que
justifican el aporte. Todas las demostraciones de este capítulo son propias, salvo que se indique
la fuente; los resultados que se citan sin demostrar se señalan como tales.

% -----------------------------------------------------------------------------
\section{Antecedentes peruanos}
\label{sec:antecedentes}

El antecedente más directo es el de
\textcite{camiz_cerda_atoche_2026}\verificar{manuscrito fechado el 25 de abril de 2026 sin datos
de revista, volumen ni DOI. Debe confirmarse con el Dr.~Camiz su estado de publicación.}, que,
según esa misma fuente, continúa la tesis de maestría de
\textcite{atoche_murrieta_2022}\verificar{la portada y el registro del repositorio difieren en el
período del título; se usa el de la portada.}. Los autores aplican PCA, HFC y EPCA a 17
series mensuales de indicadores macrofinancieros del BCRP, entre enero de 1999 y finales de 2020.
El PCA global muestra un primer componente dominante. Dos series, el tipo de cambio interbancario
y el apalancamiento global bancario, se correlacionan poco con ese componente y quedan fuera del
análisis principal. Aplicada a las 17 series, la HFC deja aisladas esas dos series y agrupa las
otras 15 en tres grupos: salud económica y financiera, con 7 series (activos internos,
circulante, ahorro total, crédito total, fondos de pensiones, liquidez y reservas internacionales
netas); tasas de interés, con 4 series (tasas activas y pasivas, en moneda nacional y
extranjera); y riesgo del mercado financiero, con 4 series (reservas, índice bursátil de Lima,
EMBIG y coeficiente de dolarización)\verificar{la ficha distingue «reservas» (grupo de riesgo) de
«reservas internacionales netas» (grupo de salud). Confirmar en la Tabla~1 del manuscrito a qué
series del BCRP corresponden.}. A continuación, el EPCA se aplica al conjunto completo y a cada
grupo por separado, con una ventana móvil gaussiana de 48 meses. La serie del primer valor propio
se particiona en subperíodos de estabilidad mediante el algoritmo de
\textcite{fisher_1958}\verificar{páginas (789--798), que constan solo en fuentes secundarias;
volumen, número y DOI ya verificados. La ficha es preliminar: el texto no se ha leído y solo se
cita su aporte general.}, y el número de cortes se
elige con el criterio de \textcite{calinski_harabasz_1974}\verificar{ficha preliminar: texto no
consultado; páginas 1--27 solo en fuentes secundarias. Solo se cita su aporte general.}.

Los autores interpretan sus resultados por grupo. El grupo de salud financiera resulta el más
estable en el tiempo; el de tasas pierde estabilidad antes de los episodios de crisis, por lo que
podría funcionar como señal temprana; y el de riesgo de mercado es el menos unidimensional y el
más sensible a perturbaciones puntuales. El estudio recomienda además un procedimiento en dos
pasos: aplicar la HFC para identificar grupos suficientemente unidimensionales y luego el EPCA a
cada grupo para verificar su estabilidad temporal. Tres rasgos de este antecedente son decisivos
para la presente tesis. Primero, según los propios autores, la ventana de 48 meses y su
ponderación resultaron de un proceso de experimentación, no de un procedimiento formal de
selección. Segundo, la muestra no cubre el período posterior a la pandemia. Tercero, los
resultados se interpretan como indicadores por grupo, sin una clasificación explícita de fases
del ciclo.

% -----------------------------------------------------------------------------
\section{El ciclo financiero y los regímenes como objetos matemáticos}
\label{sec:ciclo}

\subsection{El ciclo financiero}

Siguiendo a \textcite{drehmann_borio_tsatsaronis_2012}, el ciclo financiero se define como el
componente de mediano plazo que comparten el crédito, la razón entre crédito y producto y los
precios de los inmuebles, con oscilaciones de 8 a 30 años. Los autores lo caracterizan con dos
métodos complementarios, los filtros de frecuencia y la datación de puntos de giro, y encuentran
que los precios bursátiles siguen una dinámica distinta, más volátil y menos sincronizada con el
crédito, por lo que los excluyen de la definición. La distinción respecto del ciclo económico, de
1 a 8 años, es también de duración, y las recesiones que coinciden con la fase de contracción del
ciclo financiero resultan, en promedio, más profundas.

Esta definición plantea tres dificultades para el caso peruano. La primera es la longitud:
identificar ciclos de 8 a 30 años exige series de varias décadas, y la muestra mensual disponible
abarca menos de veinte años. Por ello, la tesis no pretende estimar ciclos completos, sino
\textbf{regímenes}, entendidos como configuraciones estructurales dentro de la muestra. La
segunda son los datos: el índice de precios de inmuebles del BCRP es trimestral y no forma parte
del conjunto mensual, de modo que el crédito, la liquidez, las tasas y el riesgo soberano son las
aproximaciones disponibles. La tercera es el enfoque: en lugar de fijar \emph{a priori} qué
variables componen el ciclo, la escuela de Camiz deja que la estructura de correlación determine
los grupos. Ese es el sentido de la exclusión de series en \textcite{camiz_cerda_atoche_2026},
que responde a su comportamiento estadístico y no a un criterio temático.

\subsection{Régimen, transición y detección}
\label{sec:regimen}

El sistema macrofinanciero se observa como una trayectoria $\{\vx_t\}_{t=1}^{n}$, con
$\vx_t\in\RR^{p}$. Sea $\estr(\cdot)$ una función que asigna a un conjunto de meses una
\textbf{estructura} estimada: la matriz de correlaciones, el subespacio principal, la jerarquía de
la HFC o la configuración de códigos latentes. Sobre esta base se adoptan las siguientes
definiciones de trabajo, desarrolladas en el marco conceptual del proyecto.

\begin{definicion}[Régimen]\label{def:regimen}
Sea $\delta$ una disimilitud entre estructuras (por ejemplo, $1-\RV$, \cref{def:rv}) y
$\varepsilon>0$ una tolerancia. Un \emph{régimen} es un conjunto de meses consecutivos
$\interv_k$ en el que la estructura del sistema es aproximadamente constante:
\begin{equation}\label{eq:regimen}
  \delta\bigl(\estr(\Wt),\estr(\Wt')\bigr)\le\varepsilon
  \quad\text{para todo par de conjuntos de meses contiguos } \Wt,\Wt'\subset\interv_k .
\end{equation}
\end{definicion}

Un régimen es, por tanto, una propiedad de la configuración conjunta de las variables y no del
nivel de una variable aislada.

\begin{definicion}[Transición]\label{def:transicion}
Hay una \emph{transición} en el mes $\tau$ cuando la estructura cambia de manera significativa y
persistente: $\delta\bigl(\estr(\Wt_\tau^{-}),\estr(\Wt_\tau^{+})\bigr)>\varepsilon$, donde
$\Wt_\tau^{-}$ y $\Wt_\tau^{+}$ son conjuntos de meses contiguos adyacentes, antes y después de
$\tau$, y la nueva estructura se mantiene en conjuntos posteriores. Según su forma, una
transición puede ser abrupta o suave (una rotación progresiva de factores); según su alcance,
total o parcial (cuando afecta solo a un bloque de variables).
\end{definicion}

\begin{definicion}[Detección]\label{def:deteccion}
\emph{Detectar regímenes} es estimar la partición $0=\tau_0<\tau_1<\dots<\tau_K=N$ de una
trayectoria de representaciones $\vy_1,\dots,\vy_N$ ---la serie de valores propios del EPCA, los
\emph{scores} del PCA o los códigos latentes--- mediante el problema de la \cref{def:fisher}.
\end{definicion}

Así planteada, la detección es una clasificación con restricción de contigüidad temporal, que
puede contrastarse con la HFC, una clasificación de variables sin esa restricción.

\begin{definicion}[Fase]\label{def:fase}
Una \emph{fase} es una etiqueta asignada a cada régimen detectado: expansión, contracción o
transición. La segmentación usa toda la estructura multivariada, y la etiqueta se asigna después,
a partir del comportamiento del crédito dentro de cada intervalo.
\end{definicion}

Esta separación evita que la clasificación se reduzca a una relectura del crecimiento del
crédito. La regla concreta de etiquetado y sus umbrales son una propuesta preliminar del
proyecto, que se presenta en la metodología y está pendiente de discusión con el asesor.

\subsection{Validación externa}
\label{sec:validacion}

Los episodios con los que se contrastan los regímenes se fechan con fuentes externas al análisis,
nunca a partir de las propias series ni de un indicador derivado de ellas. De lo contrario, la
validación sería circular. La bibliografía revisada sirve aquí \textbf{solo para identificar
eventos candidatos, no para fecharlos}. \textcite{tuhin_et_al_2025}\verificar{lista completa de
autores contra la fuente original.} tratan la crisis financiera
global y la pandemia de COVID-19 como crisis del mercado estadounidense.
\textcite{camiz_cerda_atoche_2026} asocian a eventos como la renuncia de Kuczynski y la guerra
comercial entre Estados Unidos y China transiciones que obtienen del EPCA y de la partición de
Fisher sobre series del BCRP muy parecidas a las de esta tesis. Por esa razón, sus cortes no
pueden usarse como fechas externas. Cada evento se fecha, en cambio, con un documento externo al
análisis y con la convención de la decisión D7.

\begin{definicion}[Episodio de referencia y convención de fechado]\label{def:episodio}
Un \emph{episodio de referencia} $\epi\in\Epis$ es un par de meses $\tini{\epi}\le\tfin{\epi}$.
El inicio $\tini{\epi}$ es el mes del \emph{evento o anuncio externo} que lo abre, no el de su
impacto en el Perú, tomado del texto de una fuente con ficha. El fin $\tfin{\epi}$ es el mes del
evento externo que lo cierra, si existe una fuente para él. Si no la hay, el episodio es
\emph{puntual} y se toma $\tfin{\epi}=\tini{\epi}$ como fin efectivo: el fin no se infiere ni de las
series ni de un indicador. Todas las fechas se leen de un único archivo,
\texttt{data/referencia/episodios\_referencia.csv}.
\end{definicion}

\begin{definicion}[Compatibilidad con tolerancia]\label{def:tolerancia}
Para $h\ge0$, un mes $\tau$ es \emph{compatible} con el episodio $\epi$ con tolerancia $h$ si su
distancia al intervalo $[\tini{\epi},\tfin{\epi}]$ no supera $h$ meses, es decir, si
$\tini{\epi}-h\le\tau\le\tfin{\epi}+h$. Por la decisión D7, la tolerancia principal es $h=3$ y la de
sensibilidad, $h=6$.
\end{definicion}

\begin{observacion}[Alcance de la tolerancia]\label{obs:tolerancia}
Para un episodio puntual, la \cref{def:tolerancia} se reduce a $\abs{\tau-\tini{\epi}}\le h$. Para un
período, un mes interior está a distancia cero. D7 fija el valor de $h$, pero no dice si la
distancia se mide al intervalo o solo a sus extremos. La definición adopta la primera lectura, que
es la que implementa el código cuando amplía los episodios en el anexo de robustez
(\cref{sec:met-validacion}); la segunda queda como pregunta abierta y no se usa.
\end{observacion}

La \cref{tab:episodios} reúne los seis episodios con sus fuentes. Ninguna fecha se derivó de las
series, de los cortes de \textcite{camiz_cerda_atoche_2026} ni de un indicador.

\begin{table}[htbp]
\centering
\small
\caption{Episodios de referencia (decisión D7)}
\label{tab:episodios}
\begin{tabular}{@{}p{0.27\textwidth}p{0.33\textwidth}p{0.33\textwidth}@{}}
\toprule
Episodio & Inicio (evento externo) & Fin (evento externo) \\
\midrule
Crisis financiera global & \episodioCrisisInicio{} \parencite{lehman_brothers_2008} &
  \episodioCrisisFin{} \parencite{nber_2010} \\
\emph{Taper tantrum} & \episodioTaperInicio{} \parencite{bernanke_2013} &
  \episodioTaperFin{} \parencite{reserva_federal_2013} \\
Colapso de los precios de las materias primas & \episodioMateriasInicio{}
  \parencite{stocker_et_al_2015} & \episodioMateriasFin{} \parencite{fmi_2016} \\
Renuncia de Kuczynski (caso Odebrecht) & \episodioKuczynskiInicio{} \parencite{congreso_2018} &
  --- (episodio puntual) \\
Guerra comercial entre Estados Unidos y China & \episodioGuerraInicio{} \parencite{ustr_2018} &
  \episodioGuerraFin{} \parencite{ustr_2020} \\
Pandemia de COVID-19 & \episodioCovidInicio{} \parencite{oms_2020b} &
  \episodioCovidFin{} \parencite{nber_2021} \\
\bottomrule
\end{tabular}

\smallskip
\begin{minipage}{\textwidth}\footnotesize
Fuente: \texttt{data/referencia/episodios\_referencia.csv}, a partir de las fuentes citadas.
\end{minipage}
\end{table}

Cada fecha corresponde a un hecho que consta en su fuente. La crisis financiera global se abre con
el anuncio de la solicitud de quiebra de Lehman Brothers \parencite{lehman_brothers_2008} y se
cierra con el valle de la recesión estadounidense que fijó el comité de fechado del NBER
\parencite{nber_2010}. El \emph{taper tantrum} se abre con el testimonio del presidente de la
Reserva Federal ante el Comité Económico Conjunto del Congreso de los Estados Unidos
\parencite{bernanke_2013}\verificar{la ficha registra la fecha y la descripción del testimonio,
pero no se leyó el pasaje sobre la reducción del ritmo de compras de activos.} y se cierra con el
comunicado del FOMC que decidió reducir ese ritmo \parencite{reserva_federal_2013}. El colapso de
las materias primas se fecha por la caída del precio del petróleo, que según
\textcite{stocker_et_al_2015} empezó a mediados de 2014, y termina el mes anterior al primer
aumento de los precios de los metales y del índice general que informa el FMI
\parencite{fmi_2016}. Para los metales que exporta el Perú, la caída empieza antes, desde 2011
\parencite{fmi_2015b}, pero esa fuente no da el mes. La renuncia de Kuczynski es un episodio puntual: se fecha
con la resolución del Congreso que la acepta \parencite{congreso_2018}\verificar{la ficha se
elaboró con la nota oficial del Congreso; no se leyó el texto de la resolución ni su publicación
en \emph{El Peruano}.}, y no hay una fuente que fije su fin. La guerra comercial se abre con el
anuncio de los aranceles de la Oficina del Representante Comercial de los Estados Unidos
\parencite{ustr_2018} y se cierra con la firma del acuerdo de la «fase uno»
\parencite{ustr_2020}. La pandemia se abre con la declaración de la OMS que la caracteriza como
tal \parencite{oms_2020b} y se cierra con el valle de la actividad mensual estadounidense que fijó
el NBER \parencite{nber_2021}, con el mismo criterio que el fin de la crisis global.

Dos advertencias acompañan la tabla. Primero, el fin de la guerra comercial y el inicio de la
pandemia distan menos que la tolerancia principal, de modo que una transición en ese tramo no
puede atribuirse a uno solo de los dos episodios. Segundo, las fechas son las de eventos
externos, casi todos estadounidenses o globales: miden cuándo ocurrió el choque, no cuándo se
transmitió a las series del BCRP. Con pocos episodios, algunos puntuales, la correspondencia de la
hipótesis HT3 es descriptiva y no un contraste con poder estadístico.

% -----------------------------------------------------------------------------
\section{Posición epistemológica: «datos primero»}
\label{sec:epistemologia}

La tesis se inscribe en la tradición del análisis de datos que sigue el asesor
\parencite{benzecri_1973}\verificar{ficha preliminar: obra no leída. No se ha ubicado el pasaje
que sostiene el principio de dejar hablar a los datos (\emph{laisser parler les données}), de
modo que la fórmula no se atribuye textualmente a la obra; faltan además autoría del tomo,
editorial, tomo y página.}, cuyo principio rector suele resumirse como dejar que los datos
hablen. Según este principio, la estructura debe
emerger de los datos con el mínimo de supuestos impuestos por el analista. El PCA solo exige
elegir la métrica y la HFC solo exige elegir el criterio de fusión. Ninguno de los dos postula un
modelo generador. Los antecedentes de la escuela se declaran explícitamente exploratorios, y sus
resultados describen la muestra sin generalizarla
\parencite{camiz_maulucci_roig_2010,camiz_et_al_2020}.

Esta posición tiene dos consecuencias para el diseño. La primera se refiere a la transformación
de las series. Diferenciar es un supuesto: la primera diferencia es un filtro lineal que atenúa
las oscilaciones lentas, como precisa la proposición siguiente.

\begin{proposicion}[Ganancia de la primera diferencia]\label{prop:ganancia}
Sea $x_t=\varrho\cos(\vartheta t+\psi_0)$, con $t\in\mathbb{Z}$, amplitud $\varrho>0$, fase
$\psi_0\in\RR$ y frecuencia angular $\vartheta=2\pi/P$, donde $P\ge 2$ es el período en meses.
Entonces
\begin{equation}\label{eq:ganancia}
  \Delta x_t \;=\; \gamma(P)\,\varrho\cos\!\Bigl(\vartheta t+\psi_0+\frac{\pi-\vartheta}{2}\Bigr),
  \qquad
  \gamma(P)\;=\;\abs{1-e^{-2\pi\mathrm{i}/P}}\;=\;2\sin(\pi/P).
\end{equation}
Para una sucesión $(y_t)_{t\in\mathbb{Z}}$, un mes $t_0\in\mathbb{Z}$ y una longitud
$T\in\NN$, sean $\prom{y}_{t_0,T}=\frac1T\sum_{t=t_0+1}^{t_0+T}y_t$ y
$\vart_{t_0,T}(y)=\prom{y^2}_{t_0,T}-\prom{y}_{t_0,T}^2$ la media y la varianza temporales de
$y$ en los $T$ meses consecutivos $t_0+1,\dots,t_0+T$. En consecuencia:
\textup{(a)} si $P$ es entero, $\vart_{t_0,P}(\Delta x)=\gamma(P)^2\,\vart_{t_0,P}(x)$ para todo
$t_0\in\mathbb{Z}$, es decir, sobre cualquier período completo de $P$ meses consecutivos la
varianza temporal se multiplica por $\gamma(P)^2$; y, para todo $P\ge2$ real, los límites
$\lim_{T\to\infty}\vart_{t_0,T}(x)$ y $\lim_{T\to\infty}\vart_{t_0,T}(\Delta x)$ existen, no
dependen de $t_0$ y el segundo es $\gamma(P)^2$ veces el primero, que vale $\varrho^2/2$ si $P>2$
y $\varrho^2\cos^2\psi_0$ si $P=2$. La razón de varianzas es, por tanto, $\gamma(P)^2$ en los dos
sentidos, salvo en el caso degenerado $P=2$, $\cos\psi_0=0$, en que $x\equiv0$; \textup{(b)} $\gamma$ es decreciente en $P$ y
$\gamma(P)\to 0$ cuando $P\to\infty$, con $\gamma(P)\sim 2\pi/P$; \textup{(c)} $\gamma(P)<1$
si y solo si $P>6$.
\end{proposicion}

\begin{proof}
Por la identidad $\cos\xi_1-\cos\xi_2=-2\sin\frac{\xi_1+\xi_2}{2}\sin\frac{\xi_1-\xi_2}{2}$,
con $\xi_1=\vartheta t+\psi_0$ y $\xi_2=\vartheta(t-1)+\psi_0$,
\[
  \Delta x_t=-2\varrho\sin\frac{\vartheta}{2}\,\sin\Bigl(\vartheta t+\psi_0-\frac{\vartheta}{2}\Bigr)
  =2\varrho\sin\frac{\vartheta}{2}\,\cos\Bigl(\vartheta t+\psi_0-\frac{\vartheta}{2}+\frac{\pi}{2}\Bigr),
\]
porque $-\sin\xi=\cos(\xi+\pi/2)$ para todo $\xi$. Como $P\ge2$, $\vartheta/2=\pi/P\in(0,\pi/2]$ y
$\sin(\pi/P)>0$, de modo que la amplitud es $2\varrho\sin(\pi/P)=\gamma(P)\varrho$. Además,
$\abs{1-e^{-\mathrm{i}\vartheta}}^2=(1-\cos\vartheta)^2+\sin^2\vartheta=2-2\cos\vartheta
=4\sin^2(\vartheta/2)$, lo que prueba la segunda igualdad de \cref{eq:ganancia}.
\textup{(a)} Por \cref{eq:ganancia}, $x$ y $\Delta x$ son de la forma
$y_t=\varrho_y\cos(\vartheta t+\psi_y)$, con $(\varrho_y,\psi_y)=(\varrho,\psi_0)$ y
$(\varrho_y,\psi_y)=\bigl(\gamma(P)\varrho,\,\psi_0+(\pi-\vartheta)/2\bigr)$, respectivamente.
Para $k\in\{1, 2\}$ y $\xi\in\RR$, si $e^{\mathrm{i}k\vartheta}\ne1$, la suma geométrica da
\[
  \Bigl\lvert\sum_{t=t_0+1}^{t_0+T}e^{\mathrm{i}(k\vartheta t+\xi)}\Bigr\rvert
  =\frac{\abs{e^{\mathrm{i}k\vartheta T}-1}}{\abs{e^{\mathrm{i}k\vartheta}-1}}
  \le\frac{2}{\abs{e^{\mathrm{i}k\vartheta}-1}},
\]
y la suma es nula si además $k\vartheta T\in2\pi\mathbb{Z}$. Tomando la parte real y usando
$y_t^2=\frac{\varrho_y^2}{2}+\frac{\varrho_y^2}{2}\cos(2\vartheta t+2\psi_y)$ se distinguen dos
casos. \emph{Caso $P>2$.} Entonces $\vartheta\in(0,\pi)$ y $2\vartheta\in(0,2\pi)$, de modo que
$e^{\mathrm{i}\vartheta}\ne1$ y $e^{2\mathrm{i}\vartheta}\ne1$. Si $P$ es entero, $\vartheta P=2\pi$ y
$2\vartheta P=4\pi$, las dos sumas sobre $T=P$ meses son nulas y
$\prom{y}_{t_0,P}=0$, $\prom{y^2}_{t_0,P}=\varrho_y^2/2$, luego
$\vart_{t_0,P}(y)=\varrho_y^2/2$ exactamente. Para $P$ real cualquiera, las sumas están acotadas
uniformemente en $t_0$ y $T$, de modo que $\prom{y}_{t_0,T}\to0$ y
$\prom{y^2}_{t_0,T}\to\varrho_y^2/2$, luego $\vart_{t_0,T}(y)\to\varrho_y^2/2$ sin depender de
$t_0$. En ambos sentidos la varianza pasa de $\varrho^2/2$ a $\gamma(P)^2\varrho^2/2$.
\emph{Caso $P=2$.} Entonces $\vartheta=\pi$, $e^{2\mathrm{i}\vartheta}=1$ y el promedio de
$\cos^2$ ya no es $1/2$: $x_t=\varrho(-1)^t\cos\psi_0$, de modo que $x_t^2=\varrho^2\cos^2\psi_0$
para todo $t$ y $\prom{x}_{t_0,T}$ vale $0$ si $T$ es par y tiende a $0$ en general; luego
$\vart_{t_0,2}(x)=\varrho^2\cos^2\psi_0$ y $\vart_{t_0,T}(x)\to\varrho^2\cos^2\psi_0$. Además
$x_{t-1}=-x_t$, de donde $\Delta x_t=2x_t=\gamma(2)\,x_t$ y
$\vart_{t_0,T}(\Delta x)=\gamma(2)^2\vart_{t_0,T}(x)$ para todo $t_0$ y todo $T$.
\textup{(b)} $\pi/P$ decrece en $P$ dentro de
$(0,\pi/2]$, donde el seno es creciente; y $\sin x\sim x$ cuando $x\to0$. \textup{(c)}
$2\sin(\pi/P)<1$ equivale a $\sin(\pi/P)<1/2$, es decir, a $\pi/P<\pi/6$ dentro de
$(0,\pi/2]$, o sea, $P>6$.
\end{proof}

La diferencia, por tanto, atenúa justamente las oscilaciones lentas que definen el ciclo
financiero: todas las de período mayor de seis meses salen atenuadas, y tanto más cuanto más
lentas. Trabajar sin transformar respeta los datos tal como los publica la fuente, pero mezcla la
tendencia de profundización financiera con la estructura de régimen. Por eso el proyecto mantiene
las series sin transformar como análisis oficial (decisión D1). Hasta la iteración 3b, las analizó
junto con las transformadas en \textbf{dos vías}, con el mismo código y los mismos parámetros
(decisión D2). Desde la decisión D15, el análisis principal se hace solo con las series sin
transformar, y las transformadas pasan a un anexo de robustez que conserva el sustento de no
diferenciar y la detección de episodios fuera de muestra. Ese diseño se presenta al asesor como
propuesta, no como hecho consumado (\cref{sec:met-diseno}).

La segunda consecuencia, y la tensión central del marco, se refiere a los autoencoders y las
redes LSTM. Estos modelos introducen \textbf{supuestos arquitectónicos} que no provienen de los
datos: el número de capas y su ancho, la función de activación, la dimensión latente, la longitud
de la ventana, la regularización y la semilla aleatoria. Cada uno es una decisión del analista.
La tesis no oculta esa tensión, sino que la acota con cuatro principios:

\begin{enumerate}[leftmargin=*]
  \item \textbf{Anidamiento.} Por la \cref{prop:ae-lineal}, el PCA es un caso particular del
    autoencoder. El aprendizaje profundo no sustituye el criterio de reconstrucción del análisis
    de datos, sino que amplía la clase de funciones admitida.
  \item \textbf{Anclaje lineal.} Cada modelo no lineal se reporta junto con su versión lineal,
    con la misma dimensión y en la misma escala de error. La no linealidad se acepta solo si
    mejora fuera de muestra.
  \item \textbf{Estabilidad como condición de validez.} Una estructura que cambia con la semilla
    aleatoria no puede considerarse emergida de los datos. Por eso se exige estabilidad entre
    réplicas, medida con el coeficiente RV (\cref{def:rv}), y se reporta la réplica central, no
    la de mejor ajuste.
  \item \textbf{Subordinación a la evidencia.} El aprendizaje profundo se justifica por lo que el
    PCA no puede hacer por construcción: representar la dinámica de las trayectorias (LSTM) y
    señalar observaciones que no comparten la estructura típica (error de reconstrucción). No se
    justifica por una supuesta superioridad general.
\end{enumerate}

Con estos principios, los autoencoders y las redes LSTM se presentan como un
\textbf{complemento} de la cadena PCA--HFC--EPCA, no como su reemplazo.

% -----------------------------------------------------------------------------
\section{Datos y convenciones generales}
\label{sec:datos}

\begin{definicion}[Datos y estandarización]\label{def:datos}
Los datos de una vía forman la matriz $\mX=(x_{tj})\in\RR^{n\times p}$, cuya fila
$\vx_t^{\T}$ reúne las $p$ series en el mes $t$. Para cada serie se definen la media
$\bar x_j=\frac1n\sum_{t}x_{tj}$ y la desviación típica con divisor $n$,
$\hat\sigma_j=\bigl(\frac1n\sum_t(x_{tj}-\bar x_j)^2\bigr)^{1/2}$. La matriz estandarizada es
$\mZ=(z_{tj})$, con $z_{tj}=(x_{tj}-\bar x_j)/\hat\sigma_j$; su fila $t$ es $\vz_t^{\T}$ y su
columna $j$ es $\vz_j$. La matriz de correlaciones es $\mR=\mZ^{\T}\mZ/n$, con elementos
$r_{jj'}$ y $r_{jj}=1$.
\end{definicion}

\begin{supuesto}[Datos completos y no degenerados]\label{sup:datos}
La matriz $\mX$ no tiene valores faltantes, sus meses forman un calendario continuo y ninguna
serie es constante ($\hat\sigma_j>0$ para todo $j$).
\end{supuesto}

Los métodos de la escuela de Camiz son descriptivos: no suponen un modelo probabilístico de los
datos. Los supuestos distribucionales aparecen solo en las pruebas de la
\cref{sec:adecuacion} y se declaran allí. El divisor $n$ es el del código del asesor y el de
\texttt{StandardScaler}; las correlaciones no dependen del divisor.

% -----------------------------------------------------------------------------
\section{Análisis de componentes principales}
\label{sec:pca}

\subsection{Definición y propiedad de reconstrucción}

\begin{definicion}[PCA de correlaciones]\label{def:pca}
Sea $\mR=\mV\mLambda\mV^{\T}$ la descomposición espectral de $\mR$, con
$\mLambda=\diag(\lambda_1,\dots,\lambda_p)$, $\lambda_1\ge\dots\ge\lambda_p\ge0$, y
$\mV=[\vv_1\cdots\vv_p]$ ortogonal. La \emph{componente principal} $\ell$ es
$\vc_\ell=\mZ\vv_\ell\in\RR^{n}$, de elementos $c_{t\ell}$. Para $d<p$, $\mV_d$ denota las
primeras $d$ columnas de $\mV$. La \emph{inercia} retenida por las $d$ primeras componentes es
$\sum_{\ell\le d}\lambda_\ell/p$, porque $\tr\mR=p$.
\end{definicion}

Para $\mE\in\RR^{p\times d}$ con $\mE^{\T}\mE=\mI_d$, sea $\mP=\mE\mE^{\T}$ el proyector ortogonal
sobre las columnas de $\mE$, y sea
\begin{equation}\label{eq:pca}
  J(\mE)=\frac1n\sum_{t=1}^{n}\norm{\vz_t-\mE\mE^{\T}\vz_t}^2
\end{equation}
el error cuadrático medio de reconstrucción.

\begin{proposicion}[Reconstrucción óptima]\label{prop:pca}
Para todo $\mE$ con $\mE^{\T}\mE=\mI_d$,
\[
  J(\mE)=p-\tr(\mE^{\T}\mR\mE)\;\ge\;\sum_{\ell=d+1}^{p}\lambda_\ell ,
\]
con igualdad si $\mE=\mV_d$. Si además $\lambda_d>\lambda_{d+1}$, la igualdad se cumple si y solo
si $\mE\mE^{\T}=\mV_d\mV_d^{\T}$, es decir, si las columnas de $\mE$ generan el subespacio de los
$d$ primeros vectores propios. En particular, el error mínimo por celda es
$J(\mV_d)/p=1-\sum_{\ell\le d}\lambda_\ell/p$.
\end{proposicion}

\begin{proof}
Como $\mP$ es simétrico e idempotente, $\norm{\vz_t-\mP\vz_t}^2=\vz_t^{\T}\vz_t-\vz_t^{\T}\mP\vz_t$.
Promediando en $t$ y usando $\frac1n\sum_t\vz_t\vz_t^{\T}=\mR$,
$J(\mE)=\tr\mR-\tr(\mP\mR)=p-\tr(\mE^{\T}\mR\mE)$. Sean $w_\ell=\vv_\ell^{\T}\mP\vv_\ell
=\norm{\mE^{\T}\vv_\ell}^2$. Como $\mP$ es un proyector, $0\le w_\ell\le1$, y
$\sum_\ell w_\ell=\tr(\mP\mV\mV^{\T})=\tr\mP=d$. Además,
$\tr(\mP\mR)=\sum_\ell\lambda_\ell\,\vv_\ell^{\T}\mP\vv_\ell=\sum_\ell\lambda_\ell w_\ell$. Entonces
\begin{align*}
  \sum_{\ell}\lambda_\ell w_\ell-\sum_{\ell\le d}\lambda_\ell
  &=\sum_{\ell\le d}\lambda_\ell(w_\ell-1)+\sum_{\ell>d}\lambda_\ell w_\ell
  \le\lambda_d\sum_{\ell\le d}(w_\ell-1)+\lambda_{d+1}\sum_{\ell>d}w_\ell\\
  &=(\lambda_{d+1}-\lambda_d)\sum_{\ell>d}w_\ell\le0,
\end{align*}
donde se usó $w_\ell-1\le0$, $\lambda_\ell\ge\lambda_d$ para $\ell\le d$, $\lambda_\ell\le
\lambda_{d+1}$ para $\ell>d$ y $\sum_{\ell\le d}(w_\ell-1)=-\sum_{\ell>d}w_\ell$. Luego
$\tr(\mE^{\T}\mR\mE)\le\sum_{\ell\le d}\lambda_\ell$, que es la cota, y con $\mE=\mV_d$ se tiene
$w_\ell=1$ para $\ell\le d$ y $0$ en otro caso, con igualdad. Si $\lambda_d>\lambda_{d+1}$, la
igualdad exige $\sum_{\ell>d}w_\ell=0$, es decir, $\mP\vv_\ell=\mathbf{0}$ para $\ell>d$. Entonces la
imagen de $\mP$ está contenida en el subespacio generado por $\vv_1,\dots,\vv_d$, y como
$\rango\mP=d$, coincide con él; luego $\mP=\mV_d\mV_d^{\T}$. El recíproco es inmediato.
\end{proof}

La \cref{prop:pca} es un caso particular del teorema de aproximación de rango bajo, que se
enuncia sin demostración porque la tesis no lo usa en ninguna prueba.

\begin{teorema}[Eckart y Young; enunciado citado, no demostrado]\label{teo:eckart-young}
Sea $\mN\in\RR^{n\times p}$ con valores singulares $\varsigma_1\ge\varsigma_2\ge\dots\ge0$. Para
$d<\rango\mN$, el mínimo de $\normF{\mN-\mM}^2$ entre las matrices $\mM$ de rango a lo sumo $d$ es
$\sum_{\ell>d}\varsigma_\ell^2$, y se alcanza con la descomposición en valores singulares
truncada en $d$ términos; si $\varsigma_d>\varsigma_{d+1}$, el minimizador es único
\parencite{eckart_young_1936}\verificar{texto leído parcialmente (copia con OCR defectuoso, falta
la última página): confirmar en la fuente la formulación de la condición de unicidad.}.
\end{teorema}

Una exposición general del PCA, que la tesis toma como obra de referencia del método, es la de
\textcite{jolliffe_2002}\verificar{metadatos verificados; texto no consultado. Se cita solo como
obra general, sin atribuirle ningún resultado concreto.}.

\begin{observacion}[Identificación y signo]\label{obs:signo-pca}
$J$ depende de $\mE$ solo a través del proyector $\mE\mE^{\T}$: el PCA identifica un subespacio, y
la base que lo genera, incluido el signo de cada eje, se fija por convención. Con valores propios
simples, cada $\vv_\ell$ está determinado salvo el signo. La convención de esta tesis, que es la
de \texttt{Rcamiz} y la de la función \texttt{pca\_basico} del documento en R, es la
\emph{cuadrática}: se invierte $\vv_\ell$ si
$\sum_{j:\,v_{j\ell}<0}v_{j\ell}^2>\sum_{j:\,v_{j\ell}>0}v_{j\ell}^2$. El EPCA del asesor usa otra
regla (\cref{alg:aede}).
\end{observacion}

\subsection{Calidad de representación y contribuciones}
\label{sec:cos2}

La escuela francesa acompaña cada eje con dos tipos de ayudas a la interpretación: la calidad con
que el eje representa cada variable y la contribución de cada variable o de cada observación a la
construcción del eje \pendiente{fuente de la terminología de calidad de representación
($\cos^2$) y contribuciones de la escuela francesa; sin ficha}.

\begin{definicion}[Calidad y contribución]\label{def:cos2}
Sea $\lambda_\ell>0$ y sea $\rho_{j\ell}$ la correlación entre la serie $\vz_j$ y la componente
$\vc_\ell$. Se definen:
\begin{enumerate}[label=(\alph*)]
  \item la \emph{calidad de representación} de la serie $j$ en el eje $\ell$,
    $\cos^2_{j\ell}=\rho_{j\ell}^2$, y en el plano $(\ell,\ell')$,
    $\qlt_j(\ell,\ell')=\cos^2_{j\ell}+\cos^2_{j\ell'}$;
  \item la \emph{contribución} de la serie $j$ al eje $\ell$, $\ctr_{j\ell}=v_{j\ell}^2$;
  \item la \emph{contribución} del mes $t$ al eje $\ell$, $\ctr_{t\ell}=c_{t\ell}^2/(n\lambda_\ell)$.
\end{enumerate}
\end{definicion}

\begin{proposicion}[Propiedades de $\cos^2$ y de las contribuciones]\label{prop:cos2}
Si $\lambda_1,\dots,\lambda_p>0$:
\begin{enumerate}[label=(\alph*)]
  \item $\rho_{j\ell}=\sqrt{\lambda_\ell}\,v_{j\ell}$, de modo que $\cos^2_{j\ell}=\lambda_\ell
    v_{j\ell}^2$ y $\ctr_{j\ell}=\cos^2_{j\ell}/\lambda_\ell$;
  \item $\sum_{\ell=1}^{p}\cos^2_{j\ell}=1$ para cada serie $j$;
  \item $\sum_{j=1}^{p}\ctr_{j\ell}=1$ y $\sum_{t=1}^{n}\ctr_{t\ell}=1$ para cada eje $\ell$;
  \item $\cos^2_{j\ell}$ es el cuadrado del coseno del ángulo entre $\vz_j$ y $\vc_\ell$ en
    $\RR^n$, y $\qlt_j(\ell,\ell')$ es el cuadrado de la norma de la proyección ortogonal de
    $\vz_j/\norm{\vz_j}$ sobre el plano generado por $\vc_\ell$ y $\vc_{\ell'}$.
\end{enumerate}
\end{proposicion}

\begin{proof}
Las columnas de $\mZ$ están centradas, luego también $\vc_\ell$; además
$\frac1n\vc_\ell^{\T}\vc_\ell=\vv_\ell^{\T}\mR\vv_\ell=\lambda_\ell$ y
$\frac1n\vz_j^{\T}\vz_j=1$. Por tanto
$\rho_{j\ell}=\frac1n\vz_j^{\T}\mZ\vv_\ell/\sqrt{\lambda_\ell}=(\mR\vv_\ell)_j/\sqrt{\lambda_\ell}
=\sqrt{\lambda_\ell}\,v_{j\ell}$, lo que prueba (a). (b): $\sum_\ell\lambda_\ell v_{j\ell}^2=
(\mV\mLambda\mV^{\T})_{jj}=r_{jj}=1$. (c): $\norm{\vv_\ell}=1$ y
$\sum_t c_{t\ell}^2=n\lambda_\ell$. (d): como ambos vectores están centrados, la correlación es
el coseno del ángulo entre ellos. Para $\ell\ne\ell'$,
$\vc_\ell^{\T}\vc_{\ell'}=n\,\vv_\ell^{\T}\mR\vv_{\ell'}=0$, de modo que
$\{\vc_\ell/\norm{\vc_\ell},\vc_{\ell'}/\norm{\vc_{\ell'}}\}$ es una base ortonormal del plano y la
norma al cuadrado de la proyección de un vector unitario es la suma de los cuadrados de sus
cosenos con esa base.
\end{proof}

Una serie con $\qlt_j(1,2)$ baja está mal representada en el primer plano factorial, y su
posición en ese plano no debe interpretarse.

\subsection{Número de componentes}
\label{sec:retencion}

\Cref{oe:dimension} compromete triangular el número $d$ con tres criterios. Se formalizan los dos
que tienen una definición operativa en el código; las fuentes están pendientes.

\begin{definicion}[Vara quebrada y análisis paralelo]\label{def:retencion}
\begin{enumerate}[label=(\alph*)]
  \item \emph{Vara quebrada}, cuya aplicación al PCA \textcite{camiz_pillar_2018} atribuyen a
    Frontier (1976) \pendiente{Frontier (1976), fuente primaria de la vara quebrada; sin ficha}: la
    proporción de inercia esperada de la componente $\ell$ cuando la inercia total se reparte al
    azar entre $p$ componentes es $\beta_\ell=\frac1p\sum_{i=\ell}^{p}\frac1i$. Se retienen las
    componentes iniciales consecutivas con $\lambda_\ell/p>\beta_\ell$.
  \item \emph{Análisis paralelo de Horn} \parencite{horn_1965}\verificar{ficha preliminar: se
    leyó el resumen, no el texto. El uso de un percentil, en lugar de la media de los valores
    propios simulados, es una elección de la tesis y no se atribuye a la fuente.}: $\lambda^{\mathrm
    H}_\ell$ es un percentil superior, fijado de antemano, de la distribución del valor propio
    $\ell$ de la matriz de correlaciones de una tabla $n\times p$ de variables gaussianas
    independientes, estimado por simulación. Se retienen las componentes iniciales consecutivas
    con $\lambda_\ell>\lambda^{\mathrm H}_\ell$.
\end{enumerate}
\end{definicion}

El criterio de Kaiser retiene las componentes con $\lambda_\ell>1$
\parencite{kaiser_1960}\verificar{ficha preliminar: texto no consultado; la atribución del criterio
a esta obra está por confirmar, y el volumen, el número y las páginas solo constan en fuentes
secundarias.}. Según \textcite{camiz_pillar_2018}, el criterio de Rencher generaliza la prueba de
esfericidad de Bartlett: contrasta, con un estadístico ji cuadrado, si cada valor propio sucesivo
difiere de los restantes, y cuenta la dimensión como señal si la prueba es significativa al nivel
elegido. Su fórmula no se formaliza hasta leer la fuente directa
\parencite{rencher_christensen_2012}\verificar{ficha preliminar: no se comprobó que la sección 12.6
contenga la prueba que Camiz y Pillar (2018) atribuyen a Rencher, ni qué edición citan.}. El
análisis paralelo compara con un nulo
de variables \emph{independientes en el tiempo}; con series autocorrelacionadas, ese nulo no es el
pertinente, y la tesis lo declara al usarlo.

\subsection{Adecuación muestral: medida KMO y prueba de esfericidad de Bartlett}
\label{sec:adecuacion}

Las dos medidas responden a preguntas distintas. La prueba de Bartlett pregunta si hay alguna
correlación que resumir; la medida KMO, si las correlaciones son comunes al conjunto o se explican
por pares de series. La prueba de esfericidad se remite a
\textcite{bartlett_1950,bartlett_1951} y la medida de adecuación muestral, a
\textcite{kaiser_1970} y \textcite{kaiser_rice_1974}\verificar{las cuatro fichas son preliminares y
los textos no se han consultado: solo se cita el aporte general. Las fórmulas de las
\cref{def:kmo,def:bartlett} son las del código y no se atribuyen a una de estas obras en
particular; según una fuente secundaria, la forma de la \cref{eq:kmo} es la revisada de 1974.}.

\begin{supuesto}[Matriz de correlaciones regular]\label{sup:regular}
$\mR$ es invertible ($\lambda_p>0$). Se escribe $\mR^{-1}=(r^{jj'})$.
\end{supuesto}

\begin{definicion}[Correlación anti-imagen y medida KMO]\label{def:kmo}
Bajo el \cref{sup:regular}, la correlación anti-imagen entre las series $j\ne j'$ es
\[
  q_{jj'}=-\frac{r^{jj'}}{\sqrt{r^{jj}\,r^{j'j'}}} .
\]
La medida global y la individual de Kaiser-Meyer-Olkin son
\begin{equation}\label{eq:kmo}
  \KMO=\frac{\sum_{j\ne j'}r_{jj'}^2}{\sum_{j\ne j'}r_{jj'}^2+\sum_{j\ne j'}q_{jj'}^2},
  \qquad
  \KMO_j=\frac{\sum_{j'\ne j}r_{jj'}^2}{\sum_{j'\ne j}r_{jj'}^2+\sum_{j'\ne j}q_{jj'}^2}.
\end{equation}
\end{definicion}

Cuando sus denominadores no se anulan, ambas están en $[0,1]$. La correlación anti-imagen coincide con la correlación parcial entre las
series $j$ y $j'$ dada el resto \pendiente{fuente de la relación entre la inversa de la matriz de
correlaciones y las correlaciones parciales; sin ficha}. Un $\KMO$ cercano a uno indica que las
correlaciones parciales son pequeñas frente a las totales, es decir, que la correlación entre dos
series se explica por lo que comparten con las demás. Es la fórmula de la función \texttt{kmo} del
documento en R.

\begin{definicion}[Prueba de esfericidad de Bartlett]\label{def:bartlett}
El estadístico de esfericidad es
\begin{equation}\label{eq:bartlett}
  \chi^2_{\mathrm B}=-\Bigl(n-1-\frac{2p+5}{6}\Bigr)\ln\det\mR ,
\end{equation}
que se compara con una distribución ji cuadrado con $p(p-1)/2$ grados de libertad. Es la fórmula
de la función \texttt{bartlett\_esfericidad} del documento en R.
\end{definicion}

\begin{supuesto}[Supuestos de la prueba de Bartlett]\label{sup:bartlett}
La distribución de referencia de $\chi^2_{\mathrm B}$ supone que las filas $\vx_t$ son
observaciones independientes e idénticamente distribuidas de una normal multivariada, y que la
hipótesis nula es que la matriz de correlaciones poblacional es $\mI_p$; la aproximación es
asintótica en $n$.
\end{supuesto}

\begin{lema}[Signo del estadístico]\label{lem:det}
Para toda matriz de correlaciones $\mR$, $0\le\det\mR\le1$, con $\det\mR=1$ si y solo si
$\mR=\mI_p$. En consecuencia, si $n-1>(2p+5)/6$, $\chi^2_{\mathrm B}\ge0$, con igualdad si y
solo si $\mR=\mI_p$.
\end{lema}

\begin{proof}
Los valores propios son no negativos y suman $\tr\mR=p$. Por la desigualdad entre las medias
aritmética y geométrica, $\det\mR=\prod_\ell\lambda_\ell\le\bigl(\frac1p\sum_\ell\lambda_\ell
\bigr)^p=1$, con igualdad si y solo si todos los $\lambda_\ell$ valen uno, es decir, si y solo si
$\mR=\mV\mI_p\mV^{\T}=\mI_p$. La afirmación sobre $\chi^2_{\mathrm B}$ se sigue de que
$-\ln\det\mR\ge0$.
\end{proof}

\begin{observacion}[Lectura con series de tiempo]\label{obs:bartlett-series}
Como $\ln\det\mR=\sum_\ell\ln\lambda_\ell$, el estadístico está dominado por los valores propios
más pequeños: un par de series casi colineales lleva $\lambda_p$ cerca de cero y
$\chi^2_{\mathrm B}$ a valores grandes con independencia del resto de la estructura. Por eso
conviene informarlo junto con $\ln\det\mR$ y con el número de condición
$\lambda_1/\lambda_p$. Con series autocorrelacionadas, el \cref{sup:bartlett} no se cumple y la
distribución ji cuadrado no es la de referencia; reemplazar $n$ por un tamaño muestral efectivo
es una sensibilidad que se apoya en el ajuste de los grados de libertad de una correlación entre
series autocorrelacionadas de \textcite{pyper_peterman_1998}\verificar{ficha preliminar: texto no
consultado; no se ha comprobado que la fórmula de $n$ efectivo del código sea la de este artículo,
que además tiene una fe de erratas. Solo se cita el aporte general.}. En esas condiciones,
$\chi^2_{\mathrm B}$ y $\KMO$ se leen como descripciones de $\mR$, no como pruebas.
\end{observacion}

% -----------------------------------------------------------------------------
\section{Clasificación Factorial Jerárquica}
\label{sec:hfc}

La HFC, cuya formulación original se debe a
\textcite{denimal_2007}\footnote{De esta obra se leyó el resumen; los metadatos se verificaron en
Numdam y corrigen los que daba la ficha de \textcite{camiz_et_al_2020}. \textcite{camiz_roig_2011}
citan además una presentación de Denimal (2001) en un congreso, sin actas localizadas.
\pendiente{Denimal (2001); sin ficha}}, es una clasificación jerárquica ascendente de las
\textbf{variables}, no de las observaciones. Cada grupo de variables tiene un representante, su
\emph{factor}, y las fusiones se deciden con el PCA de dos factores.

\begin{definicion}[Clasificación Factorial Jerárquica]\label{def:hfc}
Para dos vectores centrados $\fac_{\gA},\fac_{\gB}\in\RR^n$, sean
$\kappa_{\gA}=\frac1n\fac_{\gA}^{\T}\fac_{\gA}$, $\kappa_{\gB}=\frac1n\fac_{\gB}^{\T}\fac_{\gB}$ y
$\kappa_{\gA\gB}=\frac1n\fac_{\gA}^{\T}\fac_{\gB}$, y sean $\mu_1^{\gA\gB}\ge\mu_2^{\gA\gB}$ los
valores propios de la matriz
$\mathbf{K}_{\gA\gB}=\bigl(\begin{smallmatrix}\kappa_{\gA}&\kappa_{\gA\gB}\\ \kappa_{\gA\gB}&\kappa_{\gB}
\end{smallmatrix}\bigr)$, con $(\alpha,\beta)$ un vector propio unitario asociado a
$\mu_1^{\gA\gB}$. La HFC procede así:
\begin{enumerate}[label=\arabic*.]
  \item \emph{Inicio.} Cada serie es un grupo, $\gA=\{j\}$, con factor $\fac_{\{j\}}=\vz_j$.
  \item \emph{Fusión.} Entre los pares de grupos vigentes se fusiona
    \begin{equation}\label{eq:hfc}
      (\gA^\ast,\gB^\ast)=\argmin_{\gA\ne\gB}\;\mu_2^{\gA\gB}.
    \end{equation}
    El nuevo grupo $\gA^\ast\cup\gB^\ast$ tiene factor
    $\fac_{\gA^\ast\cup\gB^\ast}=\alpha\fac_{\gA^\ast}+\beta\fac_{\gB^\ast}$ y componente de
    diferencias $\dif_{\gA^\ast\cup\gB^\ast}=-\beta\fac_{\gA^\ast}+\alpha\fac_{\gB^\ast}$, ambas
    definidas salvo el signo. El \emph{índice} del nodo es $\mu_2^{\gA^\ast\gB^\ast}$.
  \item \emph{Fin.} Se repite el paso 2 hasta que queda un solo grupo; la jerarquía tiene $p-1$
    nodos.
\end{enumerate}
\end{definicion}

El índice de fusión mide la heterogeneidad del grupo resultante: cuanto menor es, mayor es la
parte de la varianza conjunta que recoge un único eje. El factor funciona como \textbf{cronología
representativa} de lo que el grupo comparte, y la componente de diferencias mide las divergencias
internas. \textcite{camiz_et_al_2020} cortan la jerarquía cuando el segundo valor propio supera
0,7, umbral que remiten a Jolliffe \parencite[2002, como se cita en][]{camiz_et_al_2020}. Muestran
también, en series dendrocronológicas, que la HFC aproxima de cerca la solución global del PCA y
que, a diferencia de este, entrega una partición explícita de las series en grupos con
representantes propios. Ese caso de estudio, sin embargo, tiene pocas series, que además son
aproximadamente estacionarias.

\begin{lema}[PCA de dos factores]\label{lem:pca2}
Con la notación de la \cref{def:hfc},
\[
  \mu_{1,2}^{\gA\gB}=\frac{\kappa_{\gA}+\kappa_{\gB}\pm
  \sqrt{(\kappa_{\gA}-\kappa_{\gB})^2+4\kappa_{\gA\gB}^2}}{2}.
\]
Si $\kappa_{\gA\gB}\ne0$, $(\alpha,\beta)$ es proporcional a
$\bigl(\kappa_{\gA\gB},\,\mu_1^{\gA\gB}-\kappa_{\gA}\bigr)$. Si $\kappa_{\gA\gB}=0$,
$\mu_1^{\gA\gB}=\max(\kappa_{\gA},\kappa_{\gB})$ y el factor fusionado es el de mayor varianza.
Además, $\frac1n\norm{\fac_{\gA\cup\gB}}^2=\mu_1^{\gA\gB}$,
$\frac1n\norm{\dif_{\gA\cup\gB}}^2=\mu_2^{\gA\gB}$ y $\mu_1^{\gA\gB}+\mu_2^{\gA\gB}=\kappa_{\gA}+
\kappa_{\gB}$.
\end{lema}

\begin{proof}
Los valores propios son las raíces de
$\mu^2-(\kappa_{\gA}+\kappa_{\gB})\mu+\kappa_{\gA}\kappa_{\gB}-\kappa_{\gA\gB}^2=0$. La primera
fila de $(\mathbf{K}_{\gA\gB}-\mu_1\mI_2)(\alpha,\beta)^{\T}=\mathbf{0}$ da $(\kappa_{\gA}-\mu_1)\alpha+\kappa_{\gA\gB}\beta=0$, de donde se obtiene
la dirección si $\kappa_{\gA\gB}\ne0$; si $\kappa_{\gA\gB}=0$, la matriz es diagonal. La varianza de
$\alpha\fac_{\gA}+\beta\fac_{\gB}$ es la forma cuadrática de $\mathbf{K}_{\gA\gB}$ en el vector propio
unitario, es decir, $\mu_1$; la de la componente de diferencias, que usa el vector propio
ortogonal, es $\mu_2$; la suma de los valores propios es la traza.
\end{proof}

Las fórmulas del \cref{lem:pca2} son las de la función \texttt{distancia\_hfc} de
\texttt{R/camiz\_port.R}, y la actualización del factor es la de \texttt{agrega\_r}.

\begin{proposicion}[Fusión de dos variables estandarizadas]\label{prop:hfc-hojas}
Sean $j\ne j'$ dos series con $r=r_{jj'}$. Su PCA de dos factores tiene
\[
  \mu_1^{\{j\}\{j'\}}=1+\abs{r},\qquad\mu_2^{\{j\}\{j'\}}=1-\abs{r},
\]
y, si $r\ne0$, factor $\fac_{\{j,j'\}}=(\vz_j+\operatorname{sgn}(r)\,\vz_{j'})/\sqrt2$ (salvo el
signo). En particular, la primera fusión de la HFC une el par de series con mayor
$\abs{r_{jj'}}$, y su índice es $1-\max_{j\ne j'}\abs{r_{jj'}}$.
\end{proposicion}

\begin{proof}
Con $\kappa_{\gA}=\kappa_{\gB}=1$ y $\kappa_{\gA\gB}=r$, el \cref{lem:pca2} da
$\mu_{1,2}=(2\pm\sqrt{4r^2})/2=1\pm\abs{r}$ y la dirección $(r,\abs{r})\propto(1,\operatorname{sgn}
r)$, de norma $\sqrt2$ tras normalizar. En el primer paso todos los grupos son hojas, y
$1-\abs{r}$ es mínimo donde $\abs{r}$ es máximo.
\end{proof}

\begin{proposicion}[Descomposición de la inercia y cota del factor]\label{prop:hfc-inercia}
Sea $\mu_1^{\text{raíz}}$ el primer valor propio de la última fusión. Entonces:
\begin{enumerate}[label=(\alph*)]
  \item $p=\mu_1^{\text{raíz}}+\sum_{\text{nodos}}\mu_2^{\text{nodo}}$: la inercia total se
    reparte entre la varianza del factor raíz y los índices de los $p-1$ nodos;
  \item el factor de todo grupo $\gA$ es $\fac_{\gA}=\sum_{j\in\gA}w^{\gA}_j\vz_j$ con
    $\sum_{j\in\gA}(w^{\gA}_j)^2=1$; por tanto, si $\mR_{\gA}$ es la submatriz de $\mR$ de
    las series de $\gA$, $\frac1n\norm{\fac_{\gA}}^2\le\lambda_1(\mR_{\gA})$, con igualdad cuando
    $\gA$ tiene dos series.
\end{enumerate}
\end{proposicion}

\begin{proof}
(a) Sea $V_s$ la suma de las varianzas de los factores vigentes tras $s$ fusiones. Al inicio,
$V_0=\sum_j\frac1n\norm{\vz_j}^2=p$. Por el \cref{lem:pca2}, cada fusión reemplaza
$\kappa_{\gA}+\kappa_{\gB}$ por $\mu_1^{\gA\gB}$, de modo que $V_{s}=V_{s-1}-\mu_2^{\text{nodo }s}$.
Tras $p-1$ fusiones queda solo el factor raíz, con $V_{p-1}=\mu_1^{\text{raíz}}$.
(b) Por inducción: es cierto para las hojas, con coeficiente uno. Si $\fac_{\gA}$ y $\fac_{\gB}$
tienen coeficientes de norma unitaria con soportes disjuntos, los de
$\alpha\fac_{\gA}+\beta\fac_{\gB}$ tienen norma al cuadrado $\alpha^2+\beta^2=1$. Entonces
$\frac1n\norm{\fac_{\gA}}^2=(\mathbf{w}^{\gA})^{\T}\mR_{\gA}\mathbf{w}^{\gA}$ con
$\norm{\mathbf{w}^{\gA}}=1$, que es a lo sumo $\lambda_1(\mR_{\gA})$ por el cociente de Rayleigh.
Con dos series, la \cref{prop:hfc-hojas} y la matriz $\bigl(\begin{smallmatrix}1&r\\r&1
\end{smallmatrix}\bigr)$ dan la igualdad.
\end{proof}

\begin{observacion}[Correspondencia con el código]\label{obs:hfc-codigo}
(a) La descomposición (a) es la tabla de inercia de \texttt{clahfac} (columnas
\texttt{Hierarchy}, \texttt{Repr\_Var} y \texttt{Total}). (b) La implementación no recorre todos
los pares en cada paso, como sugiere \cref{eq:hfc}, sino que usa un grafo de vecinos reducibles
con un umbral, técnica que se remonta a \textcite{bruynooghe_1978}\verificar{metadatos
verificados en Numdam; texto no consultado. No se le atribuye la propiedad de reducibilidad que
exigiría el índice de la HFC.}. Su equivalencia con \cref{eq:hfc} depende de una propiedad de
reducibilidad del índice que no se demuestra en esta tesis. (c) Por la \cref{prop:hfc-inercia}(b), la varianza del factor de un grupo
de más de dos series no es, en general, el primer valor propio de la submatriz de correlaciones del
grupo; la salida \texttt{eigentab\_convergido} de \texttt{clahfacsub\_r} es este último.
(d) El signo de cada factor se orienta de modo que domine la suma de los cuadrados de las
correlaciones positivas con las series del grupo.
\end{observacion}

% -----------------------------------------------------------------------------
\section{Análisis evolutivo de componentes principales}
\label{sec:epca}

El PCA y la HFC suponen implícitamente que la estructura de correlación es la misma en toda la
muestra. \textcite{camiz_maulucci_roig_2010}\verificar{datos de edición de las actas (ISBN o
DOI).} relajan ese supuesto con el EPCA: un PCA
independiente sobre cada subtabla de una ventana móvil. La versión que se formaliza aquí es la
función \texttt{aede} del paquete \texttt{AedeClahfac} del asesor, traducida a R en
\texttt{R/camiz\_port.R} (funciones \texttt{aede\_r}, \texttt{.bstatw\_ac} y \texttt{.ustatw\_ac}),
que es la fuente de verdad de esta sección.

\subsection{Ventanas, fechado y pesos}

\begin{definicion}[Ventanas y convención de fechado]\label{def:ventana}
Sea $\nw\le n$ el ancho de la ventana y $N=n-\nw+1$ el número de ventanas. La ventana
$i\in\{1,\dots,N\}$ contiene los meses $i,i+1,\dots,i+\nw-1$ y se asigna al mes
\begin{equation}\label{eq:fechado}
  t_i=i+\lfloor\nw/2\rfloor .
\end{equation}
Equivalentemente, la ventana asignada al mes $t$ es
$\Wt_t=\{t-\lfloor\nw/2\rfloor,\dots,t+\lceil\nw/2\rceil-1\}$, para
$t=1+\lfloor\nw/2\rfloor,\dots,n-\lceil\nw/2\rceil+1$. La posición del mes $s\in\Wt_t$ dentro de
la ventana es $u=s-t+\lfloor\nw/2\rfloor+1\in\{1,\dots,\nw\}$.
\end{definicion}

Si $\nw$ es impar, $t_i$ es el mes central de la ventana. Si es par, la ventana no tiene un mes
central y $t_i$ es el posterior de los dos centrales, la posición $\nw/2+1$. La
\cref{eq:fechado} reproduce la asignación de nombres de fila de \texttt{aede\_r}
(\texttt{rownames(serie)[(1 + nw/2):(wn + nw/2)]}), en la que R trunca los índices no enteros.
Los meses iniciales y finales de la muestra no reciben ninguna ventana.

\begin{definicion}[Pesos de la ventana]\label{def:pesos}
Para $\nu\ge0$, los pesos gaussianos son
\begin{equation}\label{eq:pesos}
  \omega_u=\frac{\tilde\omega_u}{\sum_{u'=1}^{\nw}\tilde\omega_{u'}},\qquad
  \tilde\omega_u=\exp\!\Bigl[-\frac12\Bigl(\nu\,\frac{u-(\nw+1)/2}{(\nw-1)/\sqrt2}\Bigr)^{2}\Bigr],
  \qquad u=1,\dots,\nw .
\end{equation}
El parámetro $\nu$ es el argumento \texttt{nsd} de \texttt{aede}. Con $\nu=0$ se obtienen los
pesos uniformes, $\omega_u=1/\nw$.
\end{definicion}

Los pesos son positivos, suman uno y son simétricos respecto del centro $(\nw+1)/2$. En los
extremos, $u-(\nw+1)/2=\pm(\nw-1)/2$, de modo que $\tilde\omega_1=\tilde\omega_{\nw}=e^{-\nu^2/4}$,
mientras que el máximo de $\tilde\omega$ es uno, en el centro. El parámetro $\nu$ fija, pues, el
cociente entre el peso de los extremos y el del centro, con independencia de $\nw$. Los pesos
uniformes corresponden a la ventana rectangular del trabajo de 2010. La alternativa que siguen
\textcite{camiz_cerda_atoche_2026} privilegia el centro de la ventana \pendiente{Camiz y Diblasi
(2013), fuente del esquema de ponderación gaussiana; sin ficha}. \textcite{atoche_murrieta_2022}
usa pesos gaussianos con un parámetro igual a 3, elegido, según el texto, tras experimentar con
varios valores; la ficha no permite confirmar que ese parámetro sea el $\nu$ de \cref{eq:pesos}\verificar{la portada y el registro del repositorio difieren en el
período del título; se leyó el texto hasta el capítulo 5.2, donde no se cita una fuente del
esquema gaussiano.}.

\subsection{Correlación ponderada}

\begin{definicion}[Estadísticos ponderados de la ventana]\label{def:epca}
Para la ventana $\Wt_t$, con el peso $\omega_s$ del mes $s$ dado por su posición, se definen
\begin{equation}\label{eq:epca}
\begin{aligned}
  \vm_t&=\sum_{s\in\Wt_t}\omega_s\vx_s, &
  \mC_t&=\sum_{s\in\Wt_t}\omega_s(\vx_s-\vm_t)(\vx_s-\vm_t)^{\T},\\
  \mD_t&=\diag(\mC_t), &
  \mR_t&=\mD_t^{-1/2}\mC_t\mD_t^{-1/2}.
\end{aligned}
\end{equation}
Los valores y vectores propios de $\mR_t$ se denotan $\lambda_\ell(t)$ y $\vv_\ell(t)$.
\end{definicion}

\begin{supuesto}[Ninguna serie constante en la ventana]\label{sup:ventana}
Para todo $j$, $(\mC_t)_{jj}>0$. Con pesos positivos, como los de \cref{eq:pesos}, esto equivale a
que ninguna serie sea constante dentro de la ventana.
\end{supuesto}

\begin{proposicion}[La correlación ponderada es una matriz de correlaciones]\label{prop:epca}
Sean $\omega_s\ge0$ con $\sum_{s\in\Wt_t}\omega_s=1$ y supóngase el \cref{sup:ventana}. Sea
$\tilde\vz_s=\mD_t^{-1/2}(\vx_s-\vm_t)$ la estandarización \emph{ponderada} de cada mes de la
ventana. Entonces:
\begin{enumerate}[label=(\alph*)]
  \item $\sum_s\omega_s\tilde\vz_s=\mathbf{0}$, $\sum_s\omega_s\tilde z_{sj}^2=1$ para todo $j$, y
    $\mR_t=\sum_s\omega_s\tilde\vz_s\tilde\vz_s^{\T}$;
  \item $\mR_t$ es simétrica, semidefinida positiva, con diagonal unitaria y elementos en
    $[-1,1]$: es una matriz de correlaciones;
  \item $\lambda_\ell(t)\ge0$ y $\sum_{\ell=1}^{p}\lambda_\ell(t)=p$ en toda ventana;
  \item $\mR_t$ no cambia si cada serie se transforma como $x_{sj}\mapsto a_j+b_jx_{sj}$ con
    $b_j>0$; en particular, es la misma con $\mX$ que con $\mZ$.
\end{enumerate}
\end{proposicion}

\begin{proof}
(a) $\sum_s\omega_s(\vx_s-\vm_t)=\vm_t-\vm_t=\mathbf{0}$ porque los pesos suman uno. La
igualdad $\mR_t=\sum_s\omega_s\tilde\vz_s\tilde\vz_s^{\T}$ es \cref{eq:epca} con el factor
$\mD_t^{-1/2}$ a cada lado, y su diagonal es $(\mC_t)_{jj}/(\mC_t)_{jj}=1$, que es la varianza
ponderada de $\tilde z_{sj}$. (b) La simetría es clara. Para $\mathbf{a}\in\RR^p$,
$\mathbf{a}^{\T}\mR_t\mathbf{a}=\sum_s\omega_s(\mathbf{a}^{\T}\tilde\vz_s)^2\ge0$. La submatriz
principal de orden dos de las series $j$ y $j'$ es semidefinida positiva con diagonal unitaria,
luego su determinante $1-(\mR_t)_{jj'}^2$ es no negativo. (c) Los valores propios de una matriz
semidefinida positiva son no negativos y suman la traza, que es $p$ por (b). (d) La media se
transforma en $a_j+b_j(\vm_t)_j$, de modo que $x_{sj}-(\vm_t)_j$ se multiplica por $b_j$. Con
$\mathbf{G}_b=\diag(b_1,\dots,b_p)$, $\mC_t$ pasa a $\mathbf{G}_b\mC_t\mathbf{G}_b$ y $\mD_t$ a
$\mathbf{G}_b^2\mD_t$, y los factores se cancelan en $\mR_t$.
\end{proof}

\begin{observacion}[Por qué la estandarización debe ser ponderada]\label{obs:m1}
Si la ventana se estandariza con la media y la varianza \emph{sin ponderar} y después se toma el
producto cruzado ponderado, la diagonal resultante es la varianza ponderada de la serie dividida
entre su varianza sin ponderar, que en general no vale uno. Basta un ejemplo de tres meses con
valores $-1,0,1$ y pesos $1/4,1/2,1/4$: la media sin ponderar es cero, la varianza sin ponderar es
$2/3$, el segundo momento ponderado es $1/2$ y el elemento diagonal vale $3/4$. Por eso la
\cref{eq:epca} define la media, la covarianza y la correlación ponderadas, que son las de
\texttt{.bstatw\_ac}. Con pesos uniformes ambas versiones coinciden; el notebook
\texttt{06\_epca.ipynb}, que estandariza cada ventana sin ponderar, corresponde a ese caso.
\end{observacion}

\subsection{El algoritmo \texttt{aede}}

\begin{algoritmo}[EPCA con datación, \texttt{aede}]\label{alg:aede}
Entradas (con la opción \texttt{standard = TRUE}): la matriz $\mX$, el ancho $\nw$, el parámetro de forma $\nu$, el número $d$ de ejes y
el número máximo de clases $K_{\max}$. Para cada ventana $i=1,\dots,N$:
\begin{enumerate}[label=\arabic*.]
  \item calcular $\mR_{t_i}$ por \cref{eq:epca} con los pesos \cref{eq:pesos};
  \item obtener $\lambda_\ell(t_i)$ y $\vv_\ell(t_i)$, $\ell=1,\dots,d$;
  \item orientar cada $\vv_\ell(t_i)$ con la regla \emph{lineal}: invertirlo si
    $\abs{\sum_{j:\,v_{j\ell}<0}v_{j\ell}}>\sum_{j:\,v_{j\ell}>0}v_{j\ell}$, y guardar
    $\sqrt{\lambda_\ell(t_i)}\,\vv_\ell(t_i)$ (salida \texttt{correlationsgn});
  \item si $i>1$, invertir además $\vv_\ell(t_i)$ cuando
    $\norm{\vv_\ell(t_i)+\vv_\ell(t_{i-1})}<\norm{\vv_\ell(t_i)-\vv_\ell(t_{i-1})}$, donde
    $\vv_\ell(t_{i-1})$ es el vector ya alineado de la ventana anterior, y guardar
    $\sqrt{\lambda_\ell(t_i)}\,\vv_\ell(t_i)$ (salida \texttt{correlationlsq}).
\end{enumerate}
Por último, aplicar la partición de Fisher (\cref{def:fisher}) a la trayectoria univariada
$y_i=\lambda_1(t_i)$, $i=1,\dots,N$, con $n_{\min}=1$, para $K=1,\dots,K_{\max}$, y calcular el
índice $\CH^{\mathrm a}_K$ de la \cref{def:ch}. Un corte tras la ventana $\tau_k$ se fecha en
$t_{\tau_k+1}$, el mes asignado a la primera ventana del segmento siguiente.
\end{algoritmo}

\begin{observacion}[Alcance de las convenciones]\label{obs:aede}
(a) Los valores propios no dependen del signo de los vectores, de modo que las reglas de los pasos
3 y 4 no afectan a la datación, que solo usa $\lambda_1$. (b) La alineación del paso 4 compara el
vector $\ell$ con el vector $\ell$ de la ventana anterior. Si $\lambda_\ell(t)$ y
$\lambda_{\ell+1}(t)$ se cruzan entre dos ventanas, el orden de los ejes cambia y la alineación
de signo no lo detecta. (c) La regla lineal del paso 3 no es la regla cuadrática del PCA global
(\cref{obs:signo-pca}): los dos signos pueden diferir en un mismo eje. (d) La partición de
\texttt{aede} no impone longitud mínima de segmento; la datación con la partición de Fisher sobre
otras trayectorias sí la impone (\cref{sec:fisher}).
\end{observacion}

La descomposición espectral de cada $\mR_t$ produce series temporales de valores propios
$\lambda_\ell(t)$ y de correlaciones entre variables y componentes. Que el primer valor propio
descienda significa que el conjunto pierde coherencia unidimensional; que cambien las
correlaciones indica qué series provocan la reconfiguración.
\textcite{camiz_maulucci_roig_2010} enmarcan el procedimiento en el análisis exploratorio: sus
resultados describen la muestra y no se generalizan a una población, en la línea del análisis
exploratorio de \textcite{tukey_1977}\verificar{metadatos verificados en dos catálogos; texto no
consultado. Se cita solo como encuadre general.}. En el mismo trabajo, el análisis STATIS dual, que compara las matrices de correlación de
todas las ventanas en un espacio común, identifica años de quiebre coherentes con el EPCA.

Un problema que estos trabajos no resuelven es la elección de $\nw$.
\textcite{camiz_maulucci_roig_2010} fijan la ventana sin un criterio estadístico explícito,
\textcite{camiz_roig_2011}\verificar{el nombre del editor de las actas (O.~Papú) se infirió por
analogía; el año figura como 2010/2011 en la ficha.} comparan dos anchos de manera informal y, como se señaló,
\textcite{camiz_cerda_atoche_2026} eligen la ventana por experimentación. El ancho regula un
compromiso: una ventana corta estima mal las matrices de correlación cuando $p$ se acerca a $\nw$,
y una ventana larga atenúa la sensibilidad al cambio estructural. Si $\nw\le p$ y los pesos son positivos, además,
$\rango\mR_t\le\nw-1<p$ y la matriz es singular, porque los $\nw$ vectores
$\tilde\vz_s$ de la \cref{prop:epca}(a) satisfacen la relación lineal
$\sum_s\omega_s\tilde\vz_s=\mathbf{0}$.

\subsection{Tamaño efectivo de una ventana y umbrales por ventana}
\label{sec:kish}

Con pesos no uniformes, una ventana de $\nw$ meses no aporta la información de $\nw$
observaciones de igual peso. La medida habitual de esa pérdida es el tamaño efectivo que se
atribuye a \textcite{kish_1965}\verificar{texto no consultado: no se verificaron el capítulo ni la
página de la fórmula; se cita solo como su atribución habitual.}.

\begin{definicion}[Tamaño efectivo de Kish]\label{def:kish}
Para pesos $\omega_1,\dots,\omega_{\nw}\ge0$ con $\sum_u\omega_u=1$, el tamaño efectivo es
\begin{equation}\label{eq:kish}
  \nomega=\Bigl(\sum_{u=1}^{\nw}\omega_u^2\Bigr)^{-1}.
\end{equation}
\end{definicion}

\begin{proposicion}[Tamaño efectivo de los pesos de la ventana]\label{prop:kish}
Sean $\omega_u\ge0$, $u=1,\dots,\nw$, con $\sum_u\omega_u=1$.
\begin{enumerate}[label=(\alph*)]
  \item $1\le\nomega\le\nw$. Además, $\nomega=\nw$ si y solo si $\omega_u=1/\nw$ para todo $u$, y
    $\nomega=1$ si y solo si un peso vale uno y los demás, cero.
  \item Con los pesos de la \cref{def:pesos}, $\nomega=\nw$ si $\nu=0$, y $\nomega<\nw$ si $\nu>0$
    y $\nw\ge3$.
  \item Con los pesos de la \cref{def:pesos}, $\nomega$ depende solo de $\nw$ y de $\nu$, no de los
    datos.
  \item Para $\nu>0$ fijo y $a_\nu=\nu/\sqrt2$, cuando $\nw\to\infty$,
    \[
      \frac{\nomega}{\nw}\;\longrightarrow\;c(\nu)=
      \frac{\Bigl(\int_{-a_\nu}^{a_\nu}e^{-x^2/2}\,dx\Bigr)^2}{2a_\nu\int_{-a_\nu}^{a_\nu}e^{-x^2}\,dx},
      \qquad 0<c(\nu)<1 .
    \]
\end{enumerate}
\end{proposicion}

\begin{proof}
(a) Por la desigualdad de Cauchy-Schwarz,
$1=\bigl(\sum_u\omega_u\cdot1\bigr)^2\le\nw\sum_u\omega_u^2$, es decir, $\nomega\le\nw$, con
igualdad si y solo si $(\omega_u)$ es proporcional a $(1,\dots,1)$, o sea, si $\omega_u=1/\nw$.
Como $0\le\omega_u\le1$, se tiene $\omega_u^2\le\omega_u$, con igualdad si y solo si
$\omega_u\in\{0, 1\}$; sumando, $\sum_u\omega_u^2\le1$, es decir, $\nomega\ge1$, con igualdad si y
solo si todos los pesos valen cero o uno, lo que, con suma uno, deja un único peso igual a uno.
(b) Con $\nu=0$ los pesos son uniformes y se aplica (a). Con $\nu>0$, $\tilde\omega_u$ es
estrictamente decreciente en $\abs{u-(\nw+1)/2}$. Si $\nw\ge3$, esa distancia vale $(\nw-1)/2$ en
$u=1$ y $(\nw-3)/2$ en $u=2$, de modo que $\tilde\omega_2>\tilde\omega_1$; los pesos no son todos
iguales y (a) da la desigualdad estricta. (c) La \cref{eq:pesos} solo contiene $u$, $\nw$ y $\nu$.
(d) Sea $\xi_u=a_\nu(2u-\nw-1)/(\nw-1)$, de modo que $\tilde\omega_u=e^{-\xi_u^2/2}$ y los $\xi_u$ son
$\nw$ puntos equiespaciados de $-a_\nu$ a $a_\nu$, con separación $2a_\nu/(\nw-1)$. Para $k\in\{1, 2\}$ sea
$S_k=\frac{2a_\nu}{\nw-1}\sum_u e^{-k\xi_u^2/2}$. Son sumas de Riemann de funciones continuas en
$[-a_\nu,a_\nu]$, salvo un término de orden $1/\nw$, de modo que $S_1\to I_1=\int_{-a_\nu}^{a_\nu}e^{-x^2/2}dx$ y
$S_2\to I_2=\int_{-a_\nu}^{a_\nu}e^{-x^2}dx$. Como $\nomega=(\sum_u\tilde\omega_u)^2/\sum_u\tilde\omega_u^2$,
\[
  \frac{\nomega}{\nw}=\frac{\nw-1}{\nw}\cdot\frac{S_1^2}{2a_\nu\,S_2}\;\longrightarrow\;
  \frac{I_1^2}{2a_\nu\,I_2}=c(\nu).
\]
Por la desigualdad de Cauchy-Schwarz en $[-a_\nu,a_\nu]$,
$I_1^2=\bigl(\int e^{-x^2/2}\cdot1\,dx\bigr)^2\le2a_\nu\int e^{-x^2}dx=2a_\nu\,I_2$, con igualdad solo si
$e^{-x^2/2}$ fuera constante en $[-a_\nu,a_\nu]$, lo que no ocurre con $a_\nu>0$. Luego $c(\nu)<1$, y
$c(\nu)>0$ porque $I_1>0$.
\end{proof}

Con $\nw=2$ los dos pesos gaussianos son iguales por simetría, y el caso queda fuera de (b). La
\cref{prop:kish}(c) y (d) muestra que la reducción de $\nomega$ frente a $\nw$ con pesos
gaussianos es una propiedad de los pesos, aproximadamente proporcional al ancho, y no un rasgo de
los datos. Por eso el tamaño efectivo, y no el ancho, fija la escala con que se juzgan los
valores propios de cada ventana.

\begin{definicion}[Umbrales por ventana]\label{def:umbral-ventana}
Sea $\lambda^{\mathrm H}_\ell(n',p)$ el umbral del análisis paralelo de la \cref{def:retencion}(b)
calculado con una tabla de $n'$ filas y $p$ columnas, con el mismo percentil que el análisis
global. Para la ventana de centro $t$ y pesos de tamaño efectivo $\nomega$, sean
$\nhorn=\lceil\nomega\rceil$ y
\begin{enumerate}[label=(\alph*)]
  \item el umbral de Kaiser, $\lambda_\ell(t)>1$;
  \item el umbral de Horn por ventana, $\lambda_\ell(t)>\lambda^{\mathrm H}_\ell(\nhorn,p)$.
\end{enumerate}
Para cada eje $\ell$, ancho y esquema de pesos, se informa el número de ventanas en que
$\lambda_\ell(t)$ \emph{no} supera cada umbral (decisión D14).
\end{definicion}

\begin{observacion}[Alcance de los umbrales por ventana]\label{obs:umbral-ventana}
(a) Con pesos uniformes, $\nhorn=\nw$ por la \cref{prop:kish}(a), y el umbral es el del análisis
paralelo con $\nw$ observaciones. Con pesos gaussianos, sustituir el ancho por el tamaño efectivo
es una aproximación: el nulo de Horn usa filas gaussianas independientes de igual peso, no la
correlación ponderada de la \cref{def:epca}, y la tesis no demuestra que ambos nulos coincidan.
(b) Como en la \cref{sec:retencion}, ese nulo supone meses independientes, que las series
autocorrelacionadas no cumplen. (c) Por la \cref{prop:epca}(c),
$\sum_{\ell\le d}\lambda_\ell(t)\le p$ en toda ventana que cumpla el \cref{sup:ventana}. Las
ventanas en que una serie es constante quedan fuera de la \cref{def:epca}: el código las señala,
y en ellas la matriz que diagonaliza \texttt{aede} tiene un cero en la diagonal y traza menor que
$p$. (d) En el cálculo, $\nhorn$ se obtiene como $\lceil\nomega-10^{-9}\rceil$, para que el error de
redondeo de \cref{eq:kish} no convierta $\nomega=\nw$ en $\nw+1$.
\end{observacion}

\subsection{La regla de transiciones del notebook 06}
\label{sec:regla06}

La regla del notebook \texttt{06\_epca.ipynb}, anterior a la iteración 3b, marca transiciones con
la estabilidad de un eje entre ventanas sucesivas. Se formaliza aquí porque las decisiones D13 y
D17 la usan como evidencia.

\begin{definicion}[Regla del 06]\label{def:regla06}
Sean $\nw$ el ancho, $\paso\ge1$ el paso entre ventanas y $\ell_0$ el eje que se sigue. Las ventanas
empiezan en los meses $i_k=1+(k-1)\paso$, $k=1,\dots,M$, con $i_M+\nw-1\le n$, y se fechan en
$t_{i_k}=i_k+\lfloor\nw/2\rfloor$ (\cref{def:ventana}). En cada ventana, las series se estandarizan
con la media y la desviación típica de la ventana, sin ponderar, y $\mathbf{a}_k\in\RR^p$ es el
vector propio unitario del eje $\ell_0$ de la matriz de correlaciones resultante. Con $\rho(\cdot,
\cdot)$ la correlación de Pearson entre las $p$ coordenadas de dos vectores, se alinea el signo de
forma sucesiva, $\tilde{\mathbf{a}}_1=\mathbf{a}_1$ y
$\tilde{\mathbf{a}}_k=\operatorname{sgn}^{+}\!\bigl(\rho(\mathbf{a}_k,\tilde{\mathbf{a}}_{k-1})\bigr)
\mathbf{a}_k$, donde $\operatorname{sgn}^{+}(x)=-1$ si $x<0$ y $1$ en otro caso, y se define la
\emph{estabilidad}
\[
  \rho^{(06)}_k=\rho(\tilde{\mathbf{a}}_k,\tilde{\mathbf{a}}_{k-1})
  =\abs{\rho(\mathbf{a}_k,\tilde{\mathbf{a}}_{k-1})},\qquad k=2,\dots,M.
\]
Con $\bar\rho$ y $s_\rho$ la media y la desviación típica (divisor $M-1$, el número de valores) de
las estabilidades, hay una \emph{transición candidata} en $t_{i_k}$ si
$\rho^{(06)}_k<\bar\rho-s_\rho$.
\end{definicion}

\begin{corolario}[La regla del 06 usa pesos uniformes]\label{cor:regla06}
Bajo el \cref{sup:ventana}, la matriz de correlaciones de cada ventana de la regla del 06 es la
$\mR_t$ de la \cref{def:epca} con pesos uniformes ($\nu=0$) y el mismo centro. En consecuencia, sus
valores propios coinciden con los de \texttt{aede} con pesos uniformes, y sus vectores propios
coinciden salvo el signo.
\end{corolario}

\begin{proof}
Con $\omega_s=1/\nw$, la media ponderada es la media simple de la ventana y $\mC_t$ es la
covarianza con divisor $\nw$, de modo que $\mR_t$ es la matriz de correlaciones de los datos
estandarizados sin ponderar, que no depende del divisor. La fecha coincide por la
\cref{def:ventana}. Los vectores propios de una matriz simétrica con valores propios simples están
determinados salvo el signo.
\end{proof}

El \cref{cor:regla06} es la base formal de la lectura de la decisión D17.3: la regla del 06
cuenta como evidencia de pesos uniformes, no de pesos gaussianos.

\begin{observacion}[Alcance de la regla del 06]\label{obs:regla06}
(a) El umbral $\bar\rho-s_\rho$ es relativo al propio perfil de estabilidades: la regla puede
marcar transiciones aunque la estructura sea constante, porque no compara con un nulo. Por eso la
decisión D10 la contrasta con el nulo de estructura (\cref{sec:nulos}). (b) La estabilidad usa la
correlación de Pearson, que centra los vectores de cargas, es decir, resta la carga media del eje;
es el mismo centrado que la decisión D17.2 descarta para el coeficiente RV
(\cref{prop:rv-cargas}(d)). (c) La regla sigue el eje $\ell_0$ por su posición. Si
$\lambda_{\ell_0}$ y $\lambda_{\ell_0+1}$ se cruzan entre dos ventanas, el eje seguido cambia, como
en la \cref{obs:aede}(b).
\end{observacion}

\subsection{Bordes de ventana}
\label{sec:bordes}

Una ventana móvil cambia de composición cada mes. Cuando un episodio con valores extremos entra en
la ventana o sale de ella, la correlación de la ventana puede saltar aunque la estructura no
cambie. Por eso las decisiones D7 y D13 descuentan las transiciones que coinciden con esos meses.

\begin{lema}[Ventanas que contienen un mes]\label{lem:ventana-mes}
Con la \cref{def:ventana}, el mes $s$ pertenece a $\Wt_t$ si y solo si
$s-\lceil\nw/2\rceil+1\le t\le s+\lfloor\nw/2\rfloor$. Hay, por tanto, $\nw$ centros cuya ventana
contiene a $s$, si ninguno queda fuera de la muestra.
\end{lema}

\begin{proof}
$\Wt_t=\{t-\lfloor\nw/2\rfloor,\dots,t+\lceil\nw/2\rceil-1\}$. La condición
$s\ge t-\lfloor\nw/2\rfloor$ equivale a $t\le s+\lfloor\nw/2\rfloor$, y
$s\le t+\lceil\nw/2\rceil-1$, a $t\ge s-\lceil\nw/2\rceil+1$. El intervalo tiene
$\lfloor\nw/2\rfloor+\lceil\nw/2\rceil=\nw$ enteros.
\end{proof}

\begin{definicion}[Borde de ventana]\label{def:borde}
Para el episodio $\epi$ y el ancho $\nw$, el \emph{centro de entrada} es
$c^{\mathrm{ent}}_{\epi}=\tini{\epi}-\lceil\nw/2\rceil+1$, el primer centro cuya ventana contiene el
inicio, y el \emph{centro de salida} es $c^{\mathrm{sal}}_{\epi}=\tfin{\epi}+\lfloor\nw/2\rfloor$, el
último cuya ventana contiene el fin (\cref{lem:ventana-mes}). Una transición fechada en el centro
$\tau$ es un \emph{borde de ventana} con tolerancia $h$ si
$\abs{\tau-c^{\mathrm{ent}}_{\epi}}\le h$ o $\abs{\tau-c^{\mathrm{sal}}_{\epi}}\le h$ para algún
$\epi\in\Epis$. La tolerancia principal es $h=3$ y la de sensibilidad, $h=6$ (decisión D7).
\end{definicion}

\subsection{Lectura multiescala del ancho de ventana}
\label{sec:multiescala}

La decisión D10 fija la grilla de anchos $\nw\in\{24, 36, 48, 60\}$ y su lectura. El ancho de 48 meses
es el de \textcite{camiz_cerda_atoche_2026} y da comparabilidad con el antecedente. Los otros tres
no compiten con él: cada ancho es una escala de observación distinta, y la grilla no busca la
ventana «correcta».

\begin{observacion}[Qué cambia con el ancho]\label{obs:multiescala}
(a) Cada ancho se analiza por separado (\texttt{aede}, regla del 06 y coeficiente RV) y se compara
con el nulo de estructura constante del \emph{mismo} ancho (\cref{sec:nulos}); ningún criterio
compara anchos entre sí para elegir uno. (b) Con pesos gaussianos, la escala efectiva es
$\nomega<\nw$ (\cref{prop:kish}). (c) Un par de ventanas adyacentes no solapadas
(\cref{def:pares}) compara dos bloques de $\nw$ meses. Si dos transiciones distan menos de $\nw$
meses, ningún par compara dos tramos de estructura constante a ambos lados de la primera, porque
la segunda ventana contiene ya la segunda transición. El ancho fija, así, la resolución temporal
con que pueden separarse dos cambios. (d) Las ventanas largas estiman mejor cada $\mR_t$, porque
la razón entre meses y series es mayor, pero promedian más meses y atenúan los cambios breves.
Es el compromiso descrito antes de esta subsección, que la grilla expone en lugar de resolver.
\end{observacion}

% -----------------------------------------------------------------------------
\section{Partición óptima de Fisher y número de clases}
\label{sec:fisher}

\textcite{camiz_roig_2011} completan la cadena con un procedimiento de datación. Sobre la serie del
primer valor propio del EPCA aplican el algoritmo de \textcite{fisher_1958}, que divide una
sucesión ordenada en intervalos contiguos de máxima homogeneidad interna, y luego realizan un PCA
independiente en cada período.

\begin{definicion}[Partición óptima de una trayectoria]\label{def:fisher}
Sea $\vy_1,\dots,\vy_N\in\RR^{q}$ una trayectoria y $n_{\min}\ge1$ una longitud mínima. Para
$1\le a\le b\le N$, sean $\bar\vy_{a:b}$ la media de $\vy_a,\dots,\vy_b$ y
\[
  Q(a,b)=\sum_{t=a}^{b}\norm{\vy_t-\bar\vy_{a:b}}^2 .
\]
Una partición admisible de $\{1,\dots,b\}$ en $K$ segmentos es una sucesión
$0=\tau_0<\tau_1<\dots<\tau_K=b$ con $\tau_k-\tau_{k-1}\ge n_{\min}$. El problema de Fisher es
\begin{equation}\label{eq:fisher}
  Q^\ast_K(b)=\min_{\text{admisibles}}\;\sum_{k=1}^{K}Q(\tau_{k-1}+1,\tau_k),
\end{equation}
con $Q^\ast_K(b)=+\infty$ si no hay particiones admisibles ($b<Kn_{\min}$). La partición óptima
de la trayectoria en $K$ segmentos es la que alcanza $Q^\ast_K(N)$.
\end{definicion}

El aporte general de \textcite{fisher_1958} consiste en mostrar que, cuando los datos están
ordenados, la partición óptima puede hallarse exactamente. A diferencia de los algoritmos de
clasificación sin restricción, el resultado no depende de una inicialización. El método no
determina $K$, que requiere un criterio adicional. La formulación con $q>1$ y con longitud mínima
se declara como adaptación mientras no se verifique el alcance del artículo original; la
\cref{prop:fisher} la justifica sin depender de él. En \texttt{aede}, $q=1$ y $n_{\min}=1$
(\texttt{fisher\_r}); en la datación del documento en R, $q\ge1$ y $n_{\min}$ es un parámetro
(\texttt{particion\_fisher}).

\begin{lema}[Costo de un segmento]\label{lem:costo}
$Q(a,b)=\sum_{t=a}^{b}\norm{\vy_t}^2-\frac{1}{b-a+1}\Bigl\lVert\sum_{t=a}^{b}\vy_t\Bigr\rVert^2$.
Con sumas acumuladas de $\vy_t$ y de $\norm{\vy_t}^2$, cada $Q(a,b)$ se calcula en $O(q)$
operaciones.
\end{lema}

\begin{proof}
Como $\sum_{t=a}^{b}\vy_t=(b-a+1)\bar\vy_{a:b}$, al desarrollar el cuadrado se obtiene
$\sum_t\norm{\vy_t-\bar\vy_{a:b}}^2=\sum_t\norm{\vy_t}^2-(b-a+1)\norm{\bar\vy_{a:b}}^2$, que es la
expresión del enunciado. Las dos sumas son diferencias de sumas acumuladas.
\end{proof}

\begin{proposicion}[Subestructura óptima]\label{prop:fisher}
Para $b\ge n_{\min}$, $Q^\ast_1(b)=Q(1,b)$, y para $K\ge2$ y $b\ge Kn_{\min}$,
\begin{equation}\label{eq:bellman}
  Q^\ast_K(b)=\min_{(K-1)n_{\min}\le a\le b-n_{\min}}\Bigl[\,Q^\ast_{K-1}(a)+Q(a+1,b)\,\Bigr].
\end{equation}
Además, si $a^\ast$ alcanza el mínimo, una partición óptima de $\{1,\dots,b\}$ en $K$ segmentos
se obtiene agregando el segmento $\{a^\ast+1,\dots,b\}$ a una partición óptima de
$\{1,\dots,a^\ast\}$ en $K-1$ segmentos. En consecuencia, \cref{eq:bellman} calcula la solución
exacta de \cref{eq:fisher} en $O(K_{\max}N^2q)$ operaciones, sin inicialización.
\end{proposicion}

\begin{proof}
El caso $K=1$ es la definición. Sea $K\ge2$. ($\ge$) Tómese una partición admisible óptima de
$\{1,\dots,b\}$ y sea $a=\tau_{K-1}$. Sus primeros $K-1$ segmentos forman una partición admisible
de $\{1,\dots,a\}$, de modo que $a\ge(K-1)n_{\min}$, y el último tiene longitud $b-a\ge n_{\min}$.
Como el objetivo es aditivo por segmentos,
$Q^\ast_K(b)=\sum_{k<K}Q(\tau_{k-1}+1,\tau_k)+Q(a+1,b)\ge Q^\ast_{K-1}(a)+Q(a+1,b)$, que es al
menos el mínimo de \cref{eq:bellman}. ($\le$) Para cada $a$ del rango, una partición óptima de
$\{1,\dots,a\}$ en $K-1$ segmentos seguida del segmento $\{a+1,\dots,b\}$ es admisible en $K$
segmentos, de modo que $Q^\ast_K(b)\le Q^\ast_{K-1}(a)+Q(a+1,b)$ para todo $a$. Las dos
desigualdades dan \cref{eq:bellman}, y la construcción de ($\le$) con $a=a^\ast$ alcanza el
óptimo. Para la complejidad, se calculan $Q^\ast_K(b)$ para $K\le K_{\max}$ y $b\le N$, cada uno
como un mínimo sobre a lo sumo $N$ valores de $a$, y cada $Q(a+1,b)$ cuesta $O(q)$ por el
\cref{lem:costo}; los cortes se recuperan guardando el $a^\ast$ de cada par $(K,b)$.
\end{proof}

La recurrencia \cref{eq:bellman} es la que implementan las funciones
\texttt{fisher\_r} y \texttt{particion\_fisher}. Cuando el mínimo se alcanza en varios $a$, ambas eligen el primero.

\subsection{La longitud mínima de segmento}
\label{sec:nmin}

La longitud mínima $n_{\min}$ restringe las particiones admisibles. Esa restricción tiene dos
consecuencias que importan para leer la grilla de $K_{\max}$.

\begin{corolario}[Factibilidad y segmentos en el mínimo]\label{cor:fisher-factible}
\begin{enumerate}[label=(\alph*)]
  \item Existe una partición admisible de $\{1,\dots,N\}$ en $K$ segmentos si y solo si
    $Kn_{\min}\le N$. El mayor número de segmentos que puede evaluarse con un máximo $K_{\max}$
    es, por tanto, $\Kef=\min\bigl(K_{\max},\lfloor N/n_{\min}\rfloor\bigr)$.
  \item Existe una partición admisible en $K$ segmentos de longitud exactamente $n_{\min}$ si y
    solo si $N=Kn_{\min}$.
\end{enumerate}
\end{corolario}

\begin{proof}
(a) Si la partición es admisible, $N=\sum_k(\tau_k-\tau_{k-1})\ge Kn_{\min}$. Recíprocamente, si
$Kn_{\min}\le N$, los cortes $\tau_k=kn_{\min}$ para $k<K$ dan $K-1$ segmentos de longitud
$n_{\min}$ y un último de longitud $N-(K-1)n_{\min}\ge n_{\min}$. (b) La suma de las longitudes
es $N$.
\end{proof}

\begin{proposicion}[Monotonía del costo óptimo]\label{prop:fisher-monotonia}
Sea $Q_K=Q^\ast_K(N)$.
\begin{enumerate}[label=(\alph*)]
  \item Si $(K+1)n_{\min}\le N$ y alguna partición óptima en $K$ segmentos tiene un segmento de
    longitud al menos $2n_{\min}$, entonces $Q_{K+1}\le Q_K$. La condición se cumple siempre que
    $N>K(2n_{\min}-1)$.
  \item En particular, con $n_{\min}=1$, $Q_{K+1}\le Q_K$ para todo $1\le K<N$.
  \item Sin la condición de (a), la desigualdad puede fallar.
\end{enumerate}
\end{proposicion}

\begin{proof}
(a) Sea $\{a,\dots,b\}$ un segmento de longitud $b-a+1\ge2n_{\min}$ de una partición óptima en $K$
segmentos. Al partirlo en $\{a,\dots,a+n_{\min}-1\}$ y $\{a+n_{\min},\dots,b\}$, ambos de longitud
al menos $n_{\min}$, se obtiene una partición admisible en $K+1$ segmentos. El costo no aumenta,
porque la media de cada parte minimiza la suma de cuadrados de esa parte:
$\sum_{t\in A}\norm{\vy_t-\bar\vy_{a:b}}^2\ge\sum_{t\in A}\norm{\vy_t-\bar\vy_A}^2$ para cada parte
$A$. Luego $Q_{K+1}\le Q_K$. Si todos los segmentos de una partición en $K$ segmentos midieran a
lo sumo $2n_{\min}-1$, se tendría $N\le K(2n_{\min}-1)$; por tanto, si $N>K(2n_{\min}-1)$, toda
partición, y en particular la óptima, tiene un segmento de longitud $2n_{\min}$ o más.
(b) Con $n_{\min}=1$, la condición $N>K$ equivale a $K<N$. (c) Sean $q=1$, $n_{\min}=10$, $N=40$ e
$y_t$ igual a $0$ para $t\le13$, a $5$ para $14\le t\le26$ y a $10$ para $t\ge27$. La partición
con cortes tras $13$ y $26$ es admisible en tres segmentos y tiene costo cero, de modo que
$Q_3=0$. En cuatro segmentos, la única partición admisible tiene longitudes $10$, $10$, $10$ y
$10$ por el \cref{cor:fisher-factible}(b), y su segundo segmento, $\{11,\dots,20\}$, contiene
valores $0$ y $5$; luego $Q_4>0=Q_3$.
\end{proof}

Cuando falla la monotonía, la reducción relativa del costo de la \cref{def:criterios-k} puede ser
negativa. Las reglas que la usan deben leerse con esa salvedad cerca de $\Kef$.

\subsection{Datación de regímenes con los \emph{scores} del PCA}
\label{sec:datacion}

\begin{definicion}[Trayectoria PCA-MM12 y datación]\label{def:datacion}
Sean $c_{t\ell}$, $\ell\le d$, las $d$ primeras componentes principales del PCA global
(\cref{def:pca}) y $w_{\mathrm m}$ la longitud de una media móvil causal. Para
$t=w_{\mathrm m},\dots,n$ se definen
\[
  \bar c_{t\ell}=\frac{1}{w_{\mathrm m}}\sum_{s=t-w_{\mathrm m}+1}^{t}c_{s\ell},
\]
y $\vy_t\in\RR^d$, el vector $(\bar c_{t1},\dots,\bar c_{td})$ con cada coordenada estandarizada con
divisor $N=n-w_{\mathrm m}+1$. La \emph{datación} es la partición óptima de la trayectoria
$\vy_{w_{\mathrm m}},\dots,\vy_n$ (\cref{def:fisher}) con longitud mínima $n_{\min}$, para
$K=1,\dots,\Kef$. Cada corte se fecha en el primer mes del segmento siguiente.
\end{definicion}

La media móvil es causal: $\bar c_{t\ell}$ solo usa meses hasta $t$. El cambio de signo de una
componente no altera la partición, porque $\norm{\vy_t-\bar\vy_{a:b}}$ no cambia si se invierte el
signo de una coordenada en todos los meses; la regla de signo del PCA
(\cref{obs:signo-pca}) es, por tanto, irrelevante aquí. La etiqueta de fase de cada régimen
(\cref{def:fase}) se asigna después de la partición.

\begin{definicion}[Criterios para el número de segmentos]\label{def:criterios-k}
Sean $Q_K=Q^\ast_K(N)$, $K=1,\dots,\Kef$, y, si $Q_K>0$,
$\red_K=(Q_K-Q_{K+1})/Q_K$, la reducción relativa del costo al pasar de $K$ a $K+1$ segmentos,
$K=1,\dots,\Kef-1$.
\begin{enumerate}[label=(\alph*)]
  \item \emph{Codo por caída máxima}: $\Kest_{\mathrm{codo}}$ es el menor
    $K\in\{2,\dots,\Kef-1\}$ que maximiza $\red_{K-1}-\red_K$.
  \item \emph{Regla K} (decisión D12), con un umbral $\upsilon$, en dos pasos aplicados en este
    orden: (1) $\Kest_{\mathrm K}=\min\{K\in\{1,\dots,\Kef-1\}:\red_K<\upsilon\}$; (2) si ese
    conjunto es vacío, $\Kest_{\mathrm K}=\Kest_{\mathrm{codo}}$. D12 fija $\upsilon=10$\,\%.
  \item Un número elegido $\Kest$ está \emph{en el tope} si $\Kest=\Kef$.
  \item Una solución es un \emph{artefacto de la longitud mínima} (decisión D17.1) si todos sus
    segmentos miden exactamente $n_{\min}$.
\end{enumerate}
\end{definicion}

\begin{observacion}[Lectura de las reglas]\label{obs:criterios-k}
(a) D12 describe el primer paso como «el menor $k$ a partir del cual la reducción relativa del
costo es menor que 10\,\%». La \cref{def:criterios-k}(b) toma el primer $K$ con $\red_K<\upsilon$,
sin exigir que la reducción siga bajo el umbral para todo $K$ posterior. Las dos lecturas
coinciden si $\red_K$ es decreciente; es la lectura del código (\texttt{datar}), que aquí se
declara. (b) Por la \cref{prop:fisher-monotonia}, $\red_K\ge0$ está garantizado solo mientras
$N>K(2n_{\min}-1)$; más allá, una reducción negativa satisface el primer paso por razones
ajenas a la estructura. (c) Por el \cref{cor:fisher-factible}(b), la marca de (d) solo puede
activarse si $N$ es múltiplo de $n_{\min}$. (d) El codo de \texttt{aede} (\cref{alg:aede}) es el
criterio (a) aplicado a la trayectoria $y_i=\lambda_1(t_i)$, con $n_{\min}=1$.
\end{observacion}

\subsection{Índice de Caliński y Harabasz}

\begin{definicion}[Índice de Caliński y Harabasz en dos variantes]\label{def:ch}
Sea $T=Q(1,N)$ la suma de cuadrados total de la trayectoria y $Q_K=Q^\ast_K(N)$ la intrasegmento
de la partición óptima en $K$ segmentos, $2\le K<N$. Se definen
\[
  \CH^{\mathrm h}_K=\frac{(T-Q_K)/(K-1)}{Q_K/(N-K)}
  \qquad\text{y}\qquad
  \CH^{\mathrm a}_K=\frac{(T-Q_K)/K}{Q_K/(N-K)} .
\]
$\CH^{\mathrm h}_K$ es la forma habitual del índice de
\textcite{calinski_harabasz_1974}\verificar{texto no consultado: contrastar en la fuente los grados
de libertad ($K-1$ y $N-K$) de la forma original antes de atribuírsela.}; $\CH^{\mathrm a}_K$ es la que calcula
\texttt{aede}, con $K$ y $N-K$ grados de libertad. El número de segmentos elegido es el menor
$K\in\{2,\dots,K_{\max}\}$ que maximiza el índice.
\end{definicion}

\begin{proposicion}[La variante de \texttt{aede} no elige menos segmentos]\label{prop:ch}
Para $2\le K<N$, $\CH^{\mathrm a}_K=\frac{K-1}{K}\,\CH^{\mathrm h}_K$. Sean $K^{\mathrm h}$ y
$K^{\mathrm a}$ los números de segmentos que eligen $\CH^{\mathrm h}$ y $\CH^{\mathrm a}$ en
$\{2,\dots,K_{\max}\}$. Si $Q_K<T$ y $Q_K>0$ para todo $K$ de la grilla, entonces
$K^{\mathrm a}\ge K^{\mathrm h}$.
\end{proposicion}

\begin{proof}
La identidad es inmediata. Las hipótesis garantizan que $\CH^{\mathrm h}_K$ está definido y es
positivo en la grilla. Supóngase $K^{\mathrm a}<K^{\mathrm h}$. Como $K^{\mathrm a}$ maximiza
$\CH^{\mathrm a}$,
\[
  \CH^{\mathrm a}_{K^{\mathrm a}}\ge\CH^{\mathrm a}_{K^{\mathrm h}}
  =\frac{K^{\mathrm h}-1}{K^{\mathrm h}}\CH^{\mathrm h}_{K^{\mathrm h}}
  \ge\frac{K^{\mathrm h}-1}{K^{\mathrm h}}\CH^{\mathrm h}_{K^{\mathrm a}}
  >\frac{K^{\mathrm a}-1}{K^{\mathrm a}}\CH^{\mathrm h}_{K^{\mathrm a}}
  =\CH^{\mathrm a}_{K^{\mathrm a}},
\]
donde la segunda desigualdad usa que $K^{\mathrm h}$ maximiza $\CH^{\mathrm h}$, y la tercera, que
$K\mapsto(K-1)/K$ es estrictamente creciente y que $\CH^{\mathrm h}_{K^{\mathrm a}}>0$. La
contradicción prueba la afirmación.
\end{proof}

\begin{observacion}[Supuestos del índice]\label{obs:ch}
El índice de Caliński y Harabasz compara la dispersión entre e intra segmentos con los grados de
libertad de observaciones independientes. Los puntos de la trayectoria de \texttt{aede}
provienen de ventanas consecutivas que comparten $\nw-1$ meses, de modo que valores consecutivos
de $\lambda_1$ se calculan con casi los mismos datos y no son independientes. En esas
condiciones, ninguna de las dos formas tiene la distribución de referencia de observaciones
independientes, y la \cref{prop:ch} muestra que la variante desplaza la elección hacia más
segmentos. Por la decisión D11, las dos variantes se informan, pero no se usan como selector del
número de segmentos, salvo que el índice muestre un óptimo interior con $K_{\max}=20$ y fuera del
rango del nulo de estructura constante (\cref{sec:nulos}).
\end{observacion}

% -----------------------------------------------------------------------------
\section{Comparación de configuraciones}
\label{sec:comparacion}

\begin{definicion}[Coeficiente RV]\label{def:rv}
Para dos configuraciones no nulas $\mF_1\in\RR^{m\times q}$ y $\mF_2\in\RR^{m\times q'}$ de las
mismas $m$ filas,
\textcite{robert_escoufier_1976}\verificar{volumen, número, páginas y DOI en JSTOR. La ficha es
preliminar: el texto no se ha leído y solo se cita su aporte general.} definen
\begin{equation}\label{eq:rv}
  \RV(\mF_1,\mF_2)=\frac{\tr\bigl(\mF_1\mF_1^{\T}\mF_2\mF_2^{\T}\bigr)}
  {\sqrt{\tr\bigl[(\mF_1\mF_1^{\T})^2\bigr]\,\tr\bigl[(\mF_2\mF_2^{\T})^2\bigr]}} .
\end{equation}
\end{definicion}

\begin{proposicion}[Propiedades del coeficiente RV]\label{prop:rv}
\begin{enumerate}[label=(\alph*)]
  \item $0\le\RV(\mF_1,\mF_2)\le1$, y $\RV=1$ si y solo si $\mF_1\mF_1^{\T}=c\,\mF_2\mF_2^{\T}$
    para algún $c>0$.
  \item Si $a_1,a_2\ne0$ y $\mO_1$, $\mO_2$ son ortogonales, entonces
    \[\RV(a_1\mF_1\mO_1,a_2\mF_2\mO_2)=\RV(\mF_1,\mF_2).\]
\end{enumerate}
\end{proposicion}

\begin{proof}
Sean $\mathbf{G}=\mF_1\mF_1^{\T}$ y $\mathbf{H}=\mF_2\mF_2^{\T}$, simétricas, semidefinidas
positivas y no nulas. El numerador es
$\prodF{\mathbf{G}}{\mathbf{H}}=\tr(\mF_1\mF_1^{\T}\mF_2\mF_2^{\T})=\normF{\mF_1^{\T}\mF_2}^2\ge0$, y
el denominador es $\normF{\mathbf{G}}\normF{\mathbf{H}}>0$. Por la desigualdad de Cauchy-Schwarz,
$\prodF{\mathbf{G}}{\mathbf{H}}\le\normF{\mathbf{G}}\normF{\mathbf{H}}$, con igualdad si y solo si
$\mathbf{G}$ y $\mathbf{H}$ son proporcionales; la constante es positiva porque ambas tienen traza
positiva. (b) $(a_1\mF_1\mO_1)(a_1\mF_1\mO_1)^{\T}=a_1^2\mathbf{G}$, y el factor $a_1^2$ se cancela;
igual con $\mF_2$.
\end{proof}

La invariancia de la \cref{prop:rv}(b) es la adecuada cuando las coordenadas solo están
identificadas salvo una transformación ortogonal, como ocurre con los signos de los ejes en el EPCA
o con los códigos de un autoencoder \textbf{lineal} (\cref{prop:ae-lineal}). En un autoencoder no
lineal, los códigos están identificados, a lo sumo, salvo transformaciones no lineales invertibles,
frente a las que el RV no es invariante. En ese caso, el RV mide solo la concordancia lineal entre
configuraciones y debe interpretarse como una cota de su similitud, no como una medida invariante.
Es también la medida que subyace a STATIS \parencite{lavit_1988}\verificar{la autoría consta en la
\emph{Three-Mode Bibliography} de Leiden y no en un catálogo; falta el ISBN. Texto no
consultado: solo se cita como fuente del método.}. En esta tesis, el RV permite usar una misma métrica para comparar ventanas sucesivas del
EPCA, réplicas de un autoencoder y configuraciones de métodos distintos, con la salvedad anterior.
Queda por verificar en la fuente cómo se comporta cuando $m$ es pequeño respecto del número de
columnas.

La \cref{def:rv} no exige centrar las columnas, y la demostración de la \cref{prop:rv} no lo usa.
Centrar es un paso previo que depende de qué representan las filas. Si las filas son meses y las
columnas, coordenadas de esos meses (\emph{scores} o códigos), las columnas se centran, como en el
uso habitual. Si las filas son series y las columnas, cargas de los ejes de una ventana, el
centrado resta la carga media de cada eje. La decisión D17.2 adopta para ese caso el RV sin
centrar.

\begin{definicion}[RV de cargas entre ventanas]\label{def:rv-cargas}
Para la ventana de centro $t$ y $d$ ejes, sea
$\mF_t=\bigl[\sqrt{\lambda_1(t)}\,\vv_1(t)\ \cdots\ \sqrt{\lambda_d(t)}\,\vv_d(t)\bigr]\in
\RR^{p\times d}$ la matriz de cargas (salida \texttt{correlationlsq} de \texttt{aede}). El
\emph{RV de cargas} entre las ventanas de centros $t$ y $t'$ es $\RV^{(d)}(t,t')=\RV(\mF_t,
\mF_{t'})$, sin centrar las columnas. Como sensibilidad (decisión D17.2) se usa el coeficiente
entre las matrices de correlación completas,
\[
  \rho_{\mR}(t,t')=\frac{\prodF{\mR_t}{\mR_{t'}}}{\normF{\mR_t}\,\normF{\mR_{t'}}}.
\]
\end{definicion}

\begin{proposicion}[El RV de cargas compara aproximaciones de rango $d$]\label{prop:rv-cargas}
Sea $\mR^{[d]}_t=\mF_t\mF_t^{\T}=\sum_{\ell\le d}\lambda_\ell(t)\vv_\ell(t)\vv_\ell(t)^{\T}$.
\begin{enumerate}[label=(\alph*)]
  \item $\RV^{(d)}(t,t')=\prodF{\mR^{[d]}_t}{\mR^{[d]}_{t'}}\big/\bigl(\normF{\mR^{[d]}_t}
    \normF{\mR^{[d]}_{t'}}\bigr)$.
  \item $\RV^{(d)}(t,t')$ no cambia si las columnas de $\mF_t$ o de $\mF_{t'}$ se permutan, cambian de
    signo o se rotan con cualquier matriz ortogonal $d\times d$. En particular, no depende de las
    reglas de signo de la \cref{alg:aede} ni del orden de los $d$ primeros ejes.
  \item Con $d=p$, $\mR^{[p]}_t=\mR_t$ y $\RV^{(p)}(t,t')=\rho_{\mR}(t,t')$.
  \item Si un eje tiene cargas iguales en todas las series, $\vv_\ell(t)=\vuno/\sqrt p$, centrar la
    columna $\ell$ de $\mF_t$ la anula.
\end{enumerate}
\end{proposicion}

\begin{proof}
(a) Es la \cref{eq:rv} con $\mF_t\mF_t^{\T}=\mR^{[d]}_t$, simétrica; la segunda igualdad del
enunciado es la descomposición espectral truncada. (b) Las matrices de permutación, las diagonales
de signos y las rotaciones son ortogonales: se aplica la \cref{prop:rv}(b) con $a_1=a_2=1$.
(c) La descomposición espectral completa da $\mR^{[p]}_t=\mR_t$, y (a) da $\rho_{\mR}$. (d) La
columna centrada es $\sqrt{\lambda_\ell(t)}\,(\mI_p-\vuno\vuno^{\T}/p)\vuno/\sqrt p=\mathbf{0}$.
\end{proof}

Por la \cref{prop:rv-cargas}(b), un cruce entre $\lambda_2$ y $\lambda_3$ (\cref{obs:aede}(b)) no
altera el RV de cargas con $d=3$, porque el subespacio de los tres primeros ejes es el mismo. Sí lo
altera un cruce entre $\lambda_d$ y $\lambda_{d+1}$, que cambia ese subespacio. La
\cref{prop:rv-cargas}(d) es el caso extremo de lo que motiva D17.2: si un eje tiene cargas casi
iguales y del mismo signo, como un eje de tendencia común, el centrado borra justo lo que las
series comparten.

\begin{definicion}[Pares de ventanas adyacentes no solapadas]\label{def:pares}
Para el ancho $\nw$ y las ventanas $i=1,\dots,N$ de la \cref{def:ventana}, el par $i$ está formado
por las ventanas $i$ e $i+\nw$, $i=1,\dots,N-\nw$. La primera contiene los meses
$i,\dots,i+\nw-1$ y la segunda, los meses $i+\nw,\dots,i+2\nw-1$: son disjuntas y contiguas. La
\emph{frontera} del par es $\fr(i)=i+\nw$, el primer mes de la segunda ventana, que con $\nw$ par
es el punto medio entre los centros $t_i$ y $t_{i+\nw}$. El \emph{RV adyacente} del par es
$\RVad{i}=\RV^{(d)}(t_i,t_{i+\nw})$. Hay $N-\nw=n-2\nw+1$ pares.
\end{definicion}

Dos ventanas consecutivas con paso uno comparten $\nw-1$ meses, y su RV es cercano a uno por
construcción. Los pares adyacentes no comparten ningún mes, que es la razón de la decisión D10.

\begin{definicion}[Rotación y distancia de Procrustes]\label{def:procrustes}
Para $\mF_1,\mF_2\in\RR^{m\times q}$ centradas, la rotación de Procrustes de $\mF_1$ hacia
$\mF_2$ es una matriz ortogonal $\mO^\ast\in\argmin_{\mO^{\T}\mO=\mI_q}\normF{\mF_1\mO-\mF_2}$, y
$\normF{\mF_1\mO^\ast-\mF_2}$ es la distancia de Procrustes.
\end{definicion}

\begin{proposicion}[Solución de Procrustes]\label{prop:procrustes}
Si $\mF_1^{\T}\mF_2=\mO_1\boldsymbol{\Sigma}\mO_2^{\T}$ es una descomposición en valores singulares,
con $\mO_1,\mO_2$ ortogonales y $\boldsymbol{\Sigma}$ diagonal no negativa, entonces
$\mO^\ast=\mO_1\mO_2^{\T}$ y
$\normF{\mF_1\mO^\ast-\mF_2}^2=\normF{\mF_1}^2+\normF{\mF_2}^2-2\tr\boldsymbol{\Sigma}$.
\end{proposicion}

\begin{proof}
$\normF{\mF_1\mO-\mF_2}^2=\normF{\mF_1}^2+\normF{\mF_2}^2-2\tr(\mO^{\T}\mF_1^{\T}\mF_2)$, porque
$\normF{\mF_1\mO}=\normF{\mF_1}$. Con $\boldsymbol{\Theta}=\mO_2^{\T}\mO^{\T}\mO_1$, que es ortogonal,
$\tr(\mO^{\T}\mO_1\boldsymbol{\Sigma}\mO_2^{\T})=\tr(\boldsymbol{\Sigma}\boldsymbol{\Theta})=
\sum_i\Sigma_{ii}\Theta_{ii}\le\sum_i\Sigma_{ii}$, porque $\abs{\Theta_{ii}}\le1$. La igualdad se alcanza con
$\boldsymbol{\Theta}=\mI_q$, es decir, con $\mO=\mO_1\mO_2^{\T}$.
\end{proof}

Es la rotación de la función \texttt{procrustes} del documento en R. La normalización de la
distancia para el criterio de ventana (\cref{oe:ventana}) queda por definir.

\begin{definicion}[Coeficiente de congruencia de Tucker]\label{def:tucker}
Para dos vectores de cargas no nulos $\mathbf{a},\mathbf{b}\in\RR^{p}$,
$\congr(\mathbf{a},\mathbf{b})=\mathbf{a}^{\T}\mathbf{b}/(\norm{\mathbf{a}}\norm{\mathbf{b}})$
\parencite{tucker_1951}\verificar{ficha preliminar: entidad emisora y lugar del informe sin
confirmar en fuente oficial; texto no consultado, de modo que no se ha comprobado que el informe
defina el coeficiente. Se cita solo como su fuente habitual.}. Su uso como índice de
similitud entre factores y los umbrales de interpretación se discuten en
\textcite{lorenzo_seva_ten_berge_2006}.
\end{definicion}

A diferencia de la correlación, la congruencia no centra los vectores, y su valor absoluto no
depende del signo con que se oriente cada eje.

% -----------------------------------------------------------------------------
\section{Nulos de referencia y criterio del ancho de ventana}
\label{sec:nulos}

Los métodos de las secciones anteriores describen la muestra y no tienen una distribución de
referencia. Las hipótesis de trabajo, sin embargo, se formulan contra lo que cabría esperar sin
estructura común (el PCA estático) o con una estructura constante (la parte dinámica; HT1). Esta
sección define los dos nulos que dan esa referencia y el criterio de OE4 que se apoya en el
segundo. Los dos se construyen con remuestreo por bloques, cuyo fundamento general para datos
dependientes se remite a \textcite{kunsch_1989}\verificar{ficha preliminar: texto no consultado;
no se verificó que el artículo defina el esquema de bloques móviles solapados del código.} y a
\textcite{lahiri_2003}\verificar{metadatos verificados; texto no consultado. Solo se cita como
manual del bootstrap por bloques.}.

\begin{definicion}[Índices por bloques]\label{def:bloques}
Para una longitud $n'$ y una longitud de bloque $\Lb\le n'$, sean $\mb=\lceil n'/\Lb\rceil$ y
$\iota_1,\dots,\iota_{\mb}$ inicios independientes y uniformes en $\{1,\dots,n'-\Lb+1\}$. La sucesión de
índices $\pib(1),\dots,\pib(n')$ son los $n'$ primeros términos de
$(\iota_1,\iota_1+1,\dots,\iota_1+\Lb-1,\;\iota_2,\dots,\iota_2+\Lb-1,\;\dots)$.
\end{definicion}

\subsection{Nulo de tendencias mixto}
\label{sec:nulo-mixto}

El nulo de tendencias (decisión D2) compara el PCA de las series con el de series independientes
que conservan cada una su propia dinámica. La decisión D5 distingue si la tendencia de una serie es
determinística (TS) o estocástica (DS), pero no si una serie DS tiene deriva; la decisión D8 lo
resuelve con el estadístico siguiente, cuyo fundamento se remite a
\textcite{newey_west_1987}\verificar{orden de los autores (la Econometric Society lista primero a
West) y DOI sin comprobar; texto no consultado. El núcleo, los 12 rezagos y el umbral de $\abs t>2$
son elecciones de la tesis (D8) y no se atribuyen a la fuente.}.

\begin{definicion}[Estadístico HAC de la deriva]\label{def:hac}
Para la serie $j$, sean $\dot x_t=\Delta x_{t+1,j}-\overline{\Delta x}_j$, $t=1,\dots,n'$, sus
incrementos centrados, con $n'=n-1$ y $\overline{\Delta x}_j$ la media de los incrementos, y sea
$m_{\mathrm{HAC}}\ge0$ un número de rezagos. Con $\hat\Gamma_k=\frac1{n'}\sum_{t=k+1}^{n'}\dot x_t
\dot x_{t-k}$,
\[
  \hat\sigma^2_{\mathrm{HAC},j}=\hat\Gamma_0+2\sum_{k=1}^{m_{\mathrm{HAC}}}
  \Bigl(1-\frac{k}{m_{\mathrm{HAC}}+1}\Bigr)\hat\Gamma_k ,
  \qquad
  t^{\mathrm{HAC}}_j=\frac{\overline{\Delta x}_j}{\sqrt{\hat\sigma^2_{\mathrm{HAC},j}/n'}} .
\]
La decisión D8 fija $m_{\mathrm{HAC}}=12$, sin corrección por grados de libertad.
\end{definicion}

\begin{lema}[El estimador es no negativo]\label{lem:hac}
Con $\dot x_t=0$ para $t\notin\{1,\dots,n'\}$ y $\bar m=m_{\mathrm{HAC}}+1$,
\[
  \hat\sigma^2_{\mathrm{HAC},j}=\frac{1}{n'\bar m}\sum_{t=1}^{n'+\bar m-1}
  \Bigl(\sum_{k=0}^{\bar m-1}\dot x_{t-k}\Bigr)^2\ge0,
\]
con igualdad si y solo si todos los $\dot x_t$ son nulos, es decir, si los incrementos de la serie
son constantes. En otro caso, $t^{\mathrm{HAC}}_j$ está bien definido.
\end{lema}

\begin{proof}
Al desarrollar el cuadrado, el producto $\dot x_s\dot x_{s'}$ con $\abs{s-s'}=k\le\bar m-1$
aparece una vez por cada $t$ con $s,s'\in\{t-\bar m+1,\dots,t\}$, es decir, $\bar m-k$ veces. Por
tanto, la suma vale
$\bar m\sum_s\dot x_s^2+2\sum_{k=1}^{\bar m-1}(\bar m-k)\sum_s\dot x_s\dot x_{s-k}
=n'\bar m\bigl[\hat\Gamma_0+2\sum_{k=1}^{\bar m-1}(1-k/\bar m)\hat\Gamma_k\bigr]$, que es
$n'\bar m\,\hat\sigma^2_{\mathrm{HAC},j}$. Si la suma es nula, todas las sumas móviles lo son; la
primera, $t=1$, da $\dot x_1=0$, y por inducción $\dot x_t=0$ para todo $t$.
\end{proof}

La propiedad del \cref{lem:hac} es la que da nombre al estimador de \textcite{newey_west_1987}; la
demostración anterior es propia y solo usa la forma del núcleo.

\begin{definicion}[Nulo de tendencias mixto]\label{def:nulo-mixto}
Cada serie $j$ recibe un tratamiento según su clase de la decisión D5 y el estadístico de la
\cref{def:hac}:
\begin{enumerate}[label=(\alph*)]
  \item serie TS: $x^\ast_{tj}=\hat x_{tj}+\hat\epsilon_{\pib(t)j}$, donde $\hat x_{tj}$ es la recta
    de mínimos cuadrados de la serie sobre el tiempo y $\hat\epsilon_{tj}=x_{tj}-\hat x_{tj}$ son sus
    residuos, remuestreados en bloques;
  \item serie DS con $\abs{t^{\mathrm{HAC}}_j}>2$ (con deriva): $x^\ast_{1j}=0$ y
    $x^\ast_{t+1,j}=x^\ast_{tj}+\Delta x_{\pib(t)+1,j}$;
  \item serie DS con $\abs{t^{\mathrm{HAC}}_j}\le2$ (sin deriva): como (b), con los incrementos menos
    su media;
  \item serie indeterminada: el nulo se calcula dos veces, una tratándola como TS y otra como DS.
\end{enumerate}
Cada serie usa sus propios inicios de bloque (\cref{def:bloques}), independientes de los de las
demás, de modo que toda relación entre series se rompe por construcción. En cada réplica se
calculan las inercias de las primeras componentes, y se informa la proporción de réplicas que
igualan o superan la inercia observada.
\end{definicion}

\begin{observacion}[Supuestos del nulo mixto]\label{obs:nulo-mixto}
(a) El nulo supone que el remuestreo por bloques reproduce la dinámica de corto plazo de cada
serie, es decir, que sus incrementos (o sus residuos, si es TS) son débilmente dependientes. El
umbral $\abs{t^{\mathrm{HAC}}_j}>2$ hereda ese supuesto. (b) D8 pide declarar las series con
$\abs{t^{\mathrm{HAC}}_j}$ cercano a 2, sin fijar qué distancia cuenta como cercana; el texto no la
fija. (c) La rama TS usa, para $n$ residuos, los inicios sorteados para $n-1$ incrementos: la
sucesión de índices tiene $\mb\Lb$ términos y cubre los $n$ meses solo si $\mb\Lb\ge n$, lo que
depende de que $\Lb$ no divida a $n-1$.
\end{observacion}

\subsection{Nulo de estructura constante}
\label{sec:nulo-estructura}

El nulo de la decisión D7 responde a HT1: qué perfil de similitud entre ventanas cabría esperar si
la estructura de correlación fuera la misma en toda la muestra. En el análisis oficial, con las
series sin transformar, se remuestrean los incrementos y se acumulan.

\begin{definicion}[Nulo de estructura constante]\label{def:nulo-estructura}
Sean $\Delta\vx_{t}=\vx_{t}-\vx_{t-1}$, $t=2,\dots,n$, los incrementos mensuales de las $p$ series y
$\pib$ una sucesión de índices por bloques (\cref{def:bloques}) con $n'=n-1$, \emph{la misma para
todas las series}. La réplica es la matriz $\mX^{(\zeta)}$ con filas
\[
  \vx^{(\zeta)}_1=\mathbf{0},\qquad \vx^{(\zeta)}_{t+1}=\vx^{(\zeta)}_t+\Delta\vx_{\pib(t)+1},\qquad
  t=1,\dots,n-1 .
\]
Con $B$ réplicas independientes, $\zeta=1,\dots,B$, se calculan para cada ancho y esquema de pesos
el número $\Kest$ que elige cada criterio de \texttt{aede} con $K_{\max}=20$ y el perfil
$(\RVad{i})_i$ de la \cref{def:pares}.
\end{definicion}

\begin{proposicion}[Propiedades del nulo de estructura]\label{prop:nulo-estructura}
\begin{enumerate}[label=(\alph*)]
  \item Cada incremento de la réplica es un vector de incrementos observado completo,
    $\vx^{(\zeta)}_{t+1}-\vx^{(\zeta)}_t=\Delta\vx_{\pib(t)+1}$, y dentro de cada bloque se conservan
    $\Lb$ incrementos consecutivos en su orden.
  \item Si cada serie se transforma como $x_{tj}\mapsto a_j+b_jx_{tj}$, con $b_j>0$, y se usan los
    mismos índices, cada réplica pasa a ser $\mX^{(\zeta)}\mathbf{G}_b$. Por tanto, las matrices
    $\mR_t$ de todas sus ventanas, los valores propios de \texttt{aede}, los $\Kest$ y los perfiles
    $\RVad{i}$ no cambian: el nulo no depende de las unidades ni del nivel de las series.
  \item Acumular desde cero en lugar de desde $\vx_1$ no cambia ninguna de esas medidas.
\end{enumerate}
\end{proposicion}

\begin{proof}
(a) Es la construcción. (b) Los incrementos pasan a $\Delta\vx_t\mathbf{G}_b$, porque las $a_j$ se
cancelan, y la acumulación desde cero es lineal; luego la réplica es $\mX^{(\zeta)}\mathbf{G}_b$.
Por la \cref{prop:epca}(d), cada $\mR_t$ es la misma, y todas las medidas citadas son funciones de
las $\mR_t$ y de sus elementos propios; el RV no depende del signo de los vectores
(\cref{prop:rv-cargas}(b)). (c) Sumar $\vx_1$ a todas las filas es la transformación de (b) con
$b_j=1$.
\end{proof}

\begin{observacion}[Supuestos y límites del nulo de estructura]\label{obs:nulo-estructura}
(a) Bajo la hipótesis de estructura constante, los incrementos se tratan como una sucesión
débilmente dependiente cuya distribución conjunta no cambia en el tiempo. El remuestreo conserva
la correlación contemporánea y la dinámica corta, pero dispersa en el tiempo cualquier cambio de
estructura más largo que un bloque. (b) Si la muestra tiene tramos de varianza muy distinta, las
réplicas los mezclan, y el nulo puede ser más fácil de superar que una estructura constante
genuina. (c) La acumulación conserva la deriva de cada serie, como la versión con deriva del nulo
de tendencias; D7 no pide el tratamiento mixto de D8 para este nulo. (d) El esquema con los
mismos inicios para todas las series es una elección de la tesis
\pendiente{fuente del bootstrap por bloques conjunto, con los mismos inicios de bloque para todas
las series; sin ficha que lo trate}.
\end{observacion}

\subsection{Criterio de OE4}
\label{sec:criterio-oe4}

La decisión D10 formula el criterio de \cref{oe:ventana} con el RV entre ventanas adyacentes no
solapadas, comparado con el nulo de estructura constante del mismo ancho, y la decisión D17.2 fija
el objeto del RV y la forma de la comparación.

\begin{definicion}[Criterio de OE4]\label{def:oe4}
Para un ancho $\nw$ y un esquema de pesos, sean $\RVad{i}$, $i=1,\dots,m$, el perfil observado
(\cref{def:pares}, con $d=3$ por la decisión D9) y $\RVadr{i}{\zeta}$ el de la réplica $\zeta$ del
nulo de estructura (\cref{def:nulo-estructura}).
\begin{enumerate}[label=(\alph*)]
  \item El umbral \emph{por posición} es $\umbralRV{i}$, el cuantil empírico de nivel 5\,\% de
    $\{\RVadr{i}{\zeta}\}_{\zeta=1}^{B}$, con interpolación lineal.
  \item La proporción observada es $\varpi_0=\frac1m\sum_i\ind\{\RVad{i}<\umbralRV{i}\}$.
  \item La proporción de la réplica $\zeta$ es
    $\varpi_\zeta=\frac1m\sum_i\ind\{\RVadr{i}{\zeta}<\umbralRVr{i}{-\zeta}\}$, donde
    $\umbralRVr{i}{-\zeta}$ es el mismo cuantil calculado con las otras $B-1$ réplicas.
  \item El p-valor es $\pval=\frac1B\#\{\zeta:\varpi_\zeta\ge\varpi_0\}$.
\end{enumerate}
El criterio principal usa $\RV^{(3)}$; la sensibilidad, $\rho_{\mR}$ (\cref{def:rv-cargas}).
\end{definicion}

El umbral se toma por posición porque, sin transformar, el nulo acumula incrementos desde cero y
no es estacionario en el tiempo: un umbral único mezclaría posiciones distintas. La proporción se
compara con su propia distribución nula, y no par por par, porque los pares contiguos son
dependientes y porque comparar $m$ pares con un mismo nivel multiplica las falsas alarmas. Así, el
criterio responde a OE4 por ancho de ventana, sin elegir uno.

\begin{supuesto}[Intercambiabilidad bajo el nulo]\label{sup:intercambiable}
Bajo la hipótesis de estructura constante, el perfil observado $\mathcal{P}_0$ y los perfiles de las réplicas
$\mathcal{P}_1,\dots,\mathcal{P}_B$ son intercambiables: su distribución conjunta no cambia al permutarlos.
\end{supuesto}

\begin{proposicion}[Validez de la versión simétrica]\label{prop:oe4}
Para $j=0,\dots,B$, sea $\tilde\varpi_j$ la proporción de posiciones $i$ en que $\mathcal{P}_{j,i}$ queda bajo el
cuantil de nivel 5\,\% de $\{\mathcal{P}_{c,i}\}_{c\ne j}$, calculado con los otros $B$ perfiles, y sea
$\widetilde{\pval}=\bigl(1+\#\{\zeta\ge1:\tilde\varpi_\zeta\ge\tilde\varpi_0\}\bigr)/(B+1)$. Bajo el \cref{sup:intercambiable},
$\Pr(\widetilde{\pval}\le\alpha_0)\le\alpha_0$ para todo $\alpha_0\in[0,1]$.
\end{proposicion}

\begin{proof}
Cada $\tilde\varpi_j$ es la misma función de $\mathcal{P}_j$ y del conjunto de los demás perfiles, de modo que permutar
los perfiles permuta los $\tilde\varpi_j$; por el supuesto, $(\tilde\varpi_0,\dots,\tilde\varpi_B)$ es intercambiable. Para
$k\in\{1,\dots,B+1\}$ sea $\mathcal{J}=\{j:\#\{c\ne j:\tilde\varpi_c\ge\tilde\varpi_j\}\le k-1\}$. Si $\mathcal{J}$ tuviera $k+1$ elementos
o más, el de menor $\tilde\varpi$ en $\mathcal{J}$ tendría al menos $k$ índices distintos con $\tilde\varpi$ mayor o igual, contra
la definición de $\mathcal{J}$; luego $\#\mathcal{J}\le k$. Por la intercambiabilidad, cada índice pertenece a $\mathcal{J}$ con
la misma probabilidad, de modo que $\Pr(0\in\mathcal{J})=\mathbb{E}[\#\mathcal{J}]/(B+1)\le k/(B+1)$. Como $0\in\mathcal{J}$
equivale a $\widetilde{\pval}\le k/(B+1)$, y $\widetilde{\pval}$ solo toma valores $k/(B+1)$, se sigue
$\Pr(\widetilde{\pval}\le\alpha_0)\le\lfloor\alpha_0(B+1)\rfloor/(B+1)\le\alpha_0$.
\end{proof}

\begin{observacion}[La versión implementada]\label{obs:oe4}
La \cref{def:oe4} difiere de la versión simétrica en dos puntos: el umbral de cada réplica se
calcula sin el perfil observado ($B-1$ perfiles en lugar de $B$), y el p-valor no suma uno ni
divide entre $B+1$. Las dos diferencias son de orden $1/B$, pero la \cref{prop:oe4} no cubre
exactamente la \cref{def:oe4}. Además, el \cref{sup:intercambiable} es una idealización: las
réplicas se remuestrean a partir de los propios incrementos observados.
\end{observacion}

\begin{definicion}[Transición fuera del nulo]\label{def:fuera-nulo}
Una transición candidata fechada en $\tau$, con un ancho y un esquema de pesos, está \emph{fuera del
nulo} si algún par $i$ del mismo ancho y pesos, con $\abs{\fr(i)-\tau}\le h$, cumple
$\RVad{i}<\umbralRV{i}$. Para la regla del 06, que no tiene pesos, se usan los pares con pesos
uniformes (\cref{cor:regla06}; decisión D17.3). La tolerancia es la principal de D7, $h=3$.
\end{definicion}

La \cref{def:fuera-nulo} es la definición operativa del código; no figura en el registro de
decisiones, y queda como pregunta para el asesor.

\begin{definicion}[Transición confirmada]\label{def:confirmada}
Sea $\tau^\circ$ un mes. Una transición candidata a lo sumo a $h=3$ meses de $\tau^\circ$ \emph{cuenta} si no es
un borde de ventana (\cref{def:borde}, $h=3$) y está fuera del nulo (\cref{def:fuera-nulo}). La
transición de $\tau^\circ$ está \emph{confirmada} si se cumplen las dos condiciones siguientes; en otro
caso es una \emph{transición no confirmada} (decisiones D13 y D17.3):
\begin{enumerate}[label=(\alph*)]
  \item hay candidatas que cuentan con al menos dos anchos distintos, entre la regla del 06 y
    \texttt{aede} con pesos uniformes;
  \item hay al menos una candidata que cuenta en \texttt{aede} con pesos gaussianos, con cualquier
    ancho.
\end{enumerate}
\end{definicion}

Las candidatas de \texttt{aede} son los cortes de su codo por caída máxima
(\cref{def:criterios-k}(a)) con $K_{\max}=10$. D17.1 fija ese máximo para la regla K de la
datación, no el criterio de \texttt{aede}; tomar el codo es la lectura del código, que se declara
como ambigüedad.

% -----------------------------------------------------------------------------
\section[Aprendizaje profundo]{Aprendizaje profundo: autoencoders, LSTM y\texorpdfstring{\\}{} LSTM-autoencoder}
\label{sec:dl}

\subsection{Autoencoders}

\begin{definicion}[Autoencoder]\label{def:ae}
Un autoencoder es un par formado por un codificador $f_\theta:\RR^p\to\RR^d$ y un decodificador
$g_\phi:\RR^d\to\RR^p$, con $d<p$, que se entrenan para reconstruir la entrada:
\begin{equation}\label{eq:ae}
  (\hat\theta,\hat\phi)\in\argmin_{\theta,\phi}\;\frac1n\sum_{t=1}^{n}
  \bigl\lVert\vz_t-g_\phi\bigl(f_\theta(\vz_t)\bigr)\bigr\rVert^2+\Omega(\theta,\phi),
\end{equation}
donde $\Omega$ es un término de regularización. El código del mes $t$ es
$\veta_t=f_{\hat\theta}(\vz_t)$. El autoencoder es \emph{lineal} si $f(\vz)=\mA\vz$ y
$g(\veta)=\mB\veta$, con $\mA\in\RR^{d\times p}$ y $\mB\in\RR^{p\times d}$, y tiene \emph{pesos
atados} si $\mB=\mA^{\T}$.
\end{definicion}

El vínculo con el PCA es formal. \textcite{bourlard_kamp_1988}\verificar{Bourlard y Kamp (1988),
número del fascículo y alcance exacto del resultado sobre las no linealidades; ficha preliminar,
texto no consultado.} muestran que la solución
óptima de un perceptrón autoasociativo en el sentido de mínimos cuadrados se obtiene por
descomposición en valores singulares. \textcite{baldi_hornik_1989}\verificar{solo el DOI; el texto completo ya se leyó.} muestran que, en el caso lineal,
el mínimo global proyecta sobre el subespacio de los primeros componentes principales y que la
superficie de error no tiene mínimos locales espurios. La proposición siguiente demuestra la
primera parte en la forma que usa la tesis; la ausencia de mínimos locales espurios se cita sin
demostrar.

\begin{proposicion}[Autoencoder lineal y PCA]\label{prop:ae-lineal}
Considérese \cref{eq:ae} sin regularización ($\Omega=0$) y con codificador y decodificador
lineales. Entonces:
\begin{enumerate}[label=(\alph*)]
  \item para todo $\mA,\mB$, el error es al menos $\sum_{\ell>d}\lambda_\ell$, y la cota se
    alcanza con $\mA=\mV_d^{\T}$, $\mB=\mV_d$;
  \item si se cumplen el \cref{sup:regular} y $\lambda_d>\lambda_{d+1}$, todo minimizador
    satisface $\mB\mA=\mV_d\mV_d^{\T}$;
  \item bajo las mismas condiciones, con pesos atados, todo minimizador es de la forma
    $\mA=\mO\mV_d^{\T}$ con $\mO\in\RR^{d\times d}$ ortogonal; en particular,
    $\mA^{\T}\mA=\mV_d\mV_d^{\T}$ y el código $\veta_t=\mO\mV_d^{\T}\vz_t$ es el vector de las $d$
    primeras componentes principales del mes $t$, rotado por $\mO$.
\end{enumerate}
\end{proposicion}

\begin{proof}
(a) Sea $\mM=\mB\mA$, de rango a lo sumo $d$. Elíjase $\mE\in\RR^{p\times d}$ con columnas
ortonormales cuyo rango contenga la imagen de $\mM$. Para cada $t$, $\mM\vz_t$ pertenece a ese
subespacio, y la proyección ortogonal $\mE\mE^{\T}\vz_t$ es el punto del subespacio más cercano a
$\vz_t$; luego $\norm{\vz_t-\mM\vz_t}\ge\norm{\vz_t-\mE\mE^{\T}\vz_t}$. Promediando y aplicando la
\cref{prop:pca}, el error es al menos $J(\mE)\ge\sum_{\ell>d}\lambda_\ell$. Con $\mA=\mV_d^{\T}$ y
$\mB=\mV_d$, $\mM=\mV_d\mV_d^{\T}$ y el error es $J(\mV_d)=\sum_{\ell>d}\lambda_\ell$.
(b) En un minimizador, las dos desigualdades de (a) son igualdades. La segunda, por la
\cref{prop:pca} y $\lambda_d>\lambda_{d+1}$, da $\mE\mE^{\T}=\mV_d\mV_d^{\T}$; la primera exige
$\mM\vz_t=\mV_d\mV_d^{\T}\vz_t$ para todo $t$, por la unicidad del punto más cercano, es decir, $\mM\mZ^{\T}=\mV_d\mV_d^{\T}\mZ^{\T}$.
Multiplicando por $\mZ/n$ a la derecha, $\mM\mR=\mV_d\mV_d^{\T}\mR$, y como $\mR$ es invertible,
$\mM=\mV_d\mV_d^{\T}$. (c) Con $\mB=\mA^{\T}$, la cota de (a) se alcanza con $\mA=\mV_d^{\T}$, de
modo que los minimizadores con pesos atados son minimizadores del problema general, y por (b)
$\mA^{\T}\mA=\mV_d\mV_d^{\T}=:\mP$. Entonces $\rango\mA=d$, de modo que $\mA\mA^{\T}$ es invertible.
De $\mP^2=\mP$ se sigue $\mA^{\T}(\mA\mA^{\T}-\mI_d)\mA=\mathbf{0}$, y multiplicando por $\mA$ a
la izquierda y por $\mA^{\T}$ a la derecha, $(\mA\mA^{\T})(\mA\mA^{\T}-\mI_d)(\mA\mA^{\T})=
\mathbf{0}$, de donde $\mA\mA^{\T}=\mI_d$. Sea $\mO=\mA\mV_d$. Las filas de $\mA$ son ortonormales y
generan la imagen de $\mA^{\T}$, que es la de $\mP$, es decir, la de $\mV_d$; luego
$\mA=\mA\mV_d\mV_d^{\T}=\mO\mV_d^{\T}$ y $\mO\mO^{\T}=\mA\mV_d\mV_d^{\T}\mA^{\T}=\mA\mA^{\T}=\mI_d$.
\end{proof}

De este resultado se sigue que el PCA es el caso particular del autoencoder con activación
identidad. Las filas de $\mA$ forman una base del subespacio principal, pero solo salvo una
rotación $\mO$ arbitraria. Como apoyo técnico,
\textcite{plaut_2018}\verificar{\emph{preprint} sin revisión por pares (fuente de exploración);
se cita solo como apoyo técnico secundario de Baldi y Hornik (1989).}, en un \emph{preprint},
muestra que las direcciones principales pueden recuperarse con una descomposición en valores
singulares de los pesos aprendidos. La \cref{prop:ae-lineal} caracteriza el mínimo global; que el
entrenamiento por gradiente lo alcance es un resultado distinto, el de ausencia de mínimos locales
espurios de \textcite{baldi_hornik_1989}, que no se demuestra aquí y que no se extiende a los
autoencoders no lineales.

\textcite{hinton_salakhutdinov_2006}\verificar{la ficha se elaboró con el material suplementario y
no con el artículo principal de \emph{Science}.} mostraron que los autoencoders profundos, con
preentrenamiento capa por capa y ajuste fino posterior, pueden producir códigos de baja dimensión
mejores que los del PCA cuando la estructura de los datos no es lineal. Sus aplicaciones (imágenes
y documentos) usan, no obstante, conjuntos de datos grandes, y el sobreajuste solo se observó en el
conjunto de rostros, pequeño en relación con la capacidad del modelo.
\textcite{vincent_et_al_2008}\verificar{DOI en ACM Digital Library.} proponen el autoencoder con
eliminación de ruido, que reconstruye la entrada limpia a partir de una versión corrompida. Los
autores ofrecen una lectura geométrica: la red aprende a proyectar los puntos perturbados de vuelta
hacia la región que ocupan los datos. A partir de esa lectura, esta tesis infiere que una
observación alejada de la estructura típica tendría un error de reconstrucción alto. La inferencia
es propia, no de los autores. \textcite{gal_ghahramani_2016}\verificar{editores del volumen; se
leyó el resumen oficial, pero no el PDF.} interpretan el entrenamiento con \emph{dropout} como
inferencia variacional aproximada, de modo que mantener el \emph{dropout} activo durante la
inferencia y promediar varias pasadas estocásticas produce una medida de la incertidumbre del
modelo.

\subsection{Redes LSTM}

Una red recurrente procesa la secuencia $\vz_1,\dots,\vz_L$ y mantiene un estado oculto
$\vh_t\in\RR^{n_h}$. \textcite{hochreiter_schmidhuber_1997}\verificar{se procesó el informe
técnico FKI-207-95 y no la versión publicada en \emph{Neural Computation}.} introdujeron la celda
de memoria LSTM para resolver el problema del gradiente que se desvanece o explota en las
dependencias temporales largas. Su núcleo es un «carrusel de error constante», una autoconexión
lineal que deja fluir el error sin atenuarse, protegido por compuertas multiplicativas de entrada y
de salida. \textcite{gers_schmidhuber_cummins_2000}\verificar{DOI contra el registro de MIT Press.
La ficha es preliminar: el texto no se ha leído y solo se cita su aporte general (la compuerta de
olvido).} añadieron la compuerta de olvido, que permite a la red reiniciar su memoria en secuencias
continuas.

\begin{definicion}[Celda LSTM con compuerta de olvido]\label{def:lstm}
Dados $\vh_{t-1}$ y $\vcel_{t-1}$, la celda calcula
\begin{equation}\label{eq:lstm}
\begin{aligned}
  \vig_t&=\sigma(\mW_{\mathrm e}\vz_t+\mU_{\mathrm e}\vh_{t-1}+\vb_{\mathrm e}), &
  \vfg_t&=\sigma(\mW_{\mathrm o}\vz_t+\mU_{\mathrm o}\vh_{t-1}+\vb_{\mathrm o}),\\
  \vog_t&=\sigma(\mW_{\mathrm s}\vz_t+\mU_{\mathrm s}\vh_{t-1}+\vb_{\mathrm s}), &
  \tilde\vcel_t&=\tanh(\mW_{\mathrm c}\vz_t+\mU_{\mathrm c}\vh_{t-1}+\vb_{\mathrm c}),\\
  \vcel_t&=\vfg_t\odot\vcel_{t-1}+\vig_t\odot\tilde\vcel_t, &
  \vh_t&=\vog_t\odot\tanh(\vcel_t),
\end{aligned}
\end{equation}
donde $\vig_t$, $\vfg_t$ y $\vog_t$ son las compuertas de entrada, de olvido y de salida, $\sigma$
es la función logística aplicada elemento a elemento y $\odot$ es el producto elemento a elemento.
Es la forma que implementan las bibliotecas actuales.
\end{definicion}

La forma aditiva de $\vcel_t$ hace que $\partial\vcel_t/\partial\vcel_{t-1}=\diag(\vfg_t)$, si se
omiten las dependencias indirectas a través de $\vh_{t-1}$: cuando $\vfg_t\approx\vuno$ la memoria
persiste, y cuando $\vfg_t\approx\mathbf{0}$ se reinicia. Este segundo comportamiento es relevante
cuando la serie cruza un cambio de régimen. \textcite{hochreiter_schmidhuber_1997} no aplican la
arquitectura a series económicas: sus experimentos son tareas sintéticas que requieren muchas
repeticiones de entrenamiento.

\subsection{LSTM-autoencoder}

\textcite{srivastava_mansimov_salakhutdinov_2015}\verificar{páginas 843--852 y editores del
volumen en PMLR.} combinan ambas ideas. Un LSTM codificador lee una secuencia, y su estado final
constituye la representación. Un LSTM decodificador reconstruye la secuencia en orden inverso. Los
autores comparan decodificadores condicionados, que reciben su salida previa, y no condicionados,
y proponen además un modelo compuesto que reconstruye el pasado y predice el futuro.

\begin{definicion}[Código de un LSTM-autoencoder]\label{def:lstm-ae}
En la adaptación de esta tesis, el codificador lee la ventana $(\vz_{t-L+1},\dots,\vz_t)$ y su
estado final $\vh_t$ se proyecta a un código $\veta_t\in\RR^{d}$, que se fecha en el mes $t$, el
último de la ventana.
\end{definicion}

A diferencia del PCA o de un autoencoder estático, el código $\veta_t$ resume una
\textbf{trayectoria}: dos meses con el mismo vector $\vz_t$ pero historias distintas pueden recibir
códigos distintos. Con el fechado de la \cref{def:lstm-ae}, el código solo usa información
disponible hasta $t$, a diferencia de las ventanas del EPCA, que se fechan en su centro
(\cref{def:ventana}).

\textcite{tuhin_et_al_2025}\verificar{lista completa de autores contra la fuente original.}
constituyen el precedente reciente más cercano en el dominio financiero. Entrenan un
LSTM-autoencoder sobre los retornos diarios de acciones del S\&P 500 entre 2000 y 2022, interpretan
los pesos de la capa central como una matriz de influencia direccional entre series y construyen con
ella redes de umbral por año. Las métricas topológicas de esas redes cambian de manera distintiva en
las crisis estudiadas, con mayor claridad que las redes construidas por métodos clásicos. Los
propios autores reconocen dos limitaciones que importan aquí: el enfoque exige muchos datos y su
interpretación económica es difícil.

% -----------------------------------------------------------------------------
\section{Criterios de comparación entre métodos}
\label{sec:criterios}

La comparación entre métodos no busca determinar cuál «funciona mejor», sino qué tipo de estructura
revela cada uno y bajo qué criterio resulta útil. Se adoptan los tres criterios de la ficha de
inscripción (interpretabilidad, estabilidad temporal y capacidad descriptiva), ampliados con la
robustez:

\begin{itemize}[leftmargin=*]
  \item \textbf{Interpretabilidad.} Indica si el resultado puede asociarse a dimensiones económicas
    comprensibles. Se evalúa con la concentración de cargas o correlaciones con las variables
    originales, la coherencia económica de los grupos y la posibilidad de explicar episodios sin
    forzar la interpretación. En los modelos no lineales, que carecen de cargas constantes, se usan
    las correlaciones entre cada variable y cada eje latente tras una rotación de Procrustes
    (\cref{def:procrustes}) hacia los \emph{scores} del PCA.
  \item \textbf{Estabilidad temporal.} Indica si la estructura se mantiene dentro de un régimen y
    cambia en las transiciones. Se evalúa con el coeficiente RV y la distancia entre subespacios de
    ventanas sucesivas, y con la persistencia de los grupos.
  \item \textbf{Capacidad descriptiva.} Indica cuánto de la información original conserva la
    representación. Se evalúa con la inercia retenida en el PCA y el error de reconstrucción,
    medido fuera de muestra con validación por bloques temporales contiguos.
  \item \textbf{Robustez.} Indica si los resultados resisten variaciones razonables del diseño:
    ancho de ventana (leído como multiescala, \cref{sec:multiescala}), conjunto de variables,
    transformación (el anexo de robustez de la decisión D15), semilla aleatoria y regla de
    selección del número de regímenes.
\end{itemize}

Además de estos criterios, la ficha de inscripción compromete una comparación \textbf{sobre un
contexto simulado}, con regímenes conocidos por construcción, que permita medir la capacidad de
recuperación de cada método y calibrar sus parámetros: número de factores, número de clases, ancho
de ventana y pesos de los análisis evolutivos. Es el único escenario en el que la «verdad» del
régimen está disponible sin circularidad. Por ello, complementa y no sustituye la validación con
episodios externos.

% -----------------------------------------------------------------------------
\section{Vacíos identificados y aporte de la tesis}
\label{sec:vacios}

La revisión permite identificar cinco vacíos:

\begin{enumerate}[leftmargin=*]
  \item \textbf{Ausencia de aprendizaje profundo en el antecedente peruano.}
    \textcite{camiz_cerda_atoche_2026} se mantiene dentro del marco exploratorio clásico.
  \item \textbf{Ausencia de una HFC evolutiva unificada.} No existe en la literatura revisada un
    algoritmo EHFC publicado. \textcite{camiz_roig_2011} y \textcite{camiz_cerda_atoche_2026}
    aplican la HFC y el EPCA de forma secuencial, y la integración se produce en la interpretación
    conjunta, no en un algoritmo.
  \item \textbf{Ausencia de una clasificación explícita de fases.}
    \textcite{camiz_cerda_atoche_2026} interpretan indicadores por grupo, pero no clasifican
    intervalos en expansión, contracción o transición ni los contrastan con episodios fechados de
    manera independiente.
  \item \textbf{Muestra anterior a la pospandemia.} La muestra de
    \textcite{camiz_cerda_atoche_2026} termina en 2020.
  \item \textbf{Ausencia de un criterio objetivo para el ancho de ventana del EPCA.} Las tres
    aplicaciones de la escuela fijan el ancho de manera informal.
\end{enumerate}

El aporte de la tesis responde a esos vacíos con cuatro contribuciones diferenciales: (1) la
incorporación de autoencoders y redes LSTM, anclados formalmente en el PCA mediante la
\cref{prop:ae-lineal}; (2) la clasificación explícita de fases, validada contra episodios fechados
con fuentes externas; (3) la actualización de la muestra hasta \fechaFinDatos{}, que cubre la
dinámica pospandemia; y (4) un criterio para el ancho de ventana del EPCA, basado en la
similitud entre ventanas adyacentes no solapadas, medida con el coeficiente RV de cargas
(\cref{def:rv-cargas,def:pares}) y contrastada con un nulo de estructura constante del mismo ancho
(\cref{def:oe4}); por la decisión D10, la grilla de anchos se lee como multiescala y no como la
búsqueda de un ancho único. La formalización de una EHFC queda como objetivo
exploratorio condicional, porque no existe una especificación previa contra la cual validarla.

% Capítulo 3 · Vacío en la iteración 3a. Se redacta cuando se discuta con el asesor la
% propuesta de la Parte C (decisión D2) y se fijen las decisiones de la iteración 3b.
\chapter{Metodología y desarrollo del modelo}
\label{cap:metodologia}

\pendiente{capítulo por redactar: diseño, datos, preprocesamiento y análisis en dos vías
(decisiones D1 a D6)}

% Capítulo 4 · Vacío por la regla 7 de CLAUDE.md: no se redacta ni se publica sin la
% aprobación explícita de Roberto y sin interpretación consolidada.
\chapter{Resultados y discusión}
\label{cap:resultados}

\pendiente{capítulo por redactar: requiere la discusión de la Parte C con el asesor y la
aprobación explícita de Roberto (regla 7)}

% Capítulo 5 · Vacío: se redacta a partir del capítulo 4.
\chapter{Conclusiones y recomendaciones}
\label{cap:conclusiones}

\pendiente{capítulo por redactar a partir del \cref{cap:resultados}}


% ---- Referencias y apéndices -------------------------------------------------
\backmatter
\printbibliography[heading=bibintoc, title={Referencias}]

\appendix
% Apéndices · Vacío en la iteración 3a. Cada apéndice va en su propio archivo y se
% incluye aquí con \input{apendices/NOMBRE}.


\end{document}
